\section{環境会計と評価・報告・監査を区別}

\subsection{平成26年度 問5-2}
\label{q:h26-reg-basic-5-2}

\noindent\textbf{原典：} sources/H26/basic.pdf、PDF物理16頁。

\textbf{公式正答：} 5（single）。

\paragraph{問題の要旨} 環境会計と評価・報告・監査を区別。

\Needspace{4\baselineskip}\paragraph{確認した条件・記述（独自要約）}\mbox{}\par

\noindent\textbf{時点：} 平成26年度に指定された資料・制度。\par

\noindent\textbf{照合：} 原典21画像で否定・五肢・組合せ表を確認。後年値で置換しない。\par

\paragraph{問われる核心} 環境会計と評価・報告・監査を区別。

\paragraph{必要知識} 対象、評価目的、会計単位、環境管理手段。

\paragraph{考え方・正答を選ぶ理由（査読待ち初稿）} 環境会計はコストと効果を認識して、貨幣単位又は物量単位で定量化・伝達する仕組みで5。LCAは設備の更新周期だけでなく製品等の全過程。環境報告は社会への情報開示、グリーン購入は必要性と環境負荷、環境監査は環境管理の点検であり、税務監査だけではない。

\Needspace{4\baselineskip}\paragraph{選択肢の検討（査読待ち初稿）}\mbox{}\par

\noindent\textbf{1：不適切} LCAは設備更新時のみ。式・定義・対応の独立判定と一致しない。\par

\noindent\textbf{2：不適切} 環境報告書は設備の自治体届出。式・定義・対応の独立判定と一致しない。\par

\noindent\textbf{3：不適切} グリーン購入は生物資源優先のみ。式・定義・対応の独立判定と一致しない。\par

\noindent\textbf{4：不適切} 環境監査は税務処理のみ。式・定義・対応の独立判定と一致しない。\par

\noindent\textbf{5：適切・選択対象} 環境会計は費用効果を金額物量で。判定と一致する。\par

\paragraph{再利用する考え方} 対象、評価目的、会計単位、環境管理手段。 指定年度・対象範囲・単位・否定を先に固定し、他年度の番号や現行値を転用しない。

\paragraph{共通理論} \hyperref[theory:env:accounting]{第\ref*{theory:env:accounting}節}

\par 参考：\cite{env58_accounting}。

\par 難易度 2（校準初稿）、優先度未評価。

\subsection{平成30年度 問5-2}
\label{q:h30-reg-basic-5-2}

\noindent\textbf{原典：} sources/H30/basic.pdf、PDF物理18頁。

\textbf{公式正答：} 3（single）。

\paragraph{問題の要旨} 環境会計の定義と他手法の誤対応。

\Needspace{4\baselineskip}\paragraph{確認した条件・記述（独自要約）}\mbox{}\par

\noindent\textbf{時点：} 平成30年度に指定された資料・制度。\par

\noindent\textbf{照合：} 原典18画像で否定・五肢・組合せ表を確認。後年値で置換しない。\par

\paragraph{問われる核心} 環境会計の定義と他手法の誤対応。

\paragraph{必要知識} 目的・対象・単位による管理手法の区別。

\paragraph{考え方・正答を選ぶ理由（査読待ち初稿）} 3の費用と効果を貨幣又は物量単位で表す仕組みが環境会計。グリーン購入は必要性と環境負荷全般、環境報告は社会への情報開示、環境監査は環境管理の評価、LCAは全過程の影響評価で、残る肢の限定が違う。

\Needspace{4\baselineskip}\paragraph{選択肢の検討（査読待ち初稿）}\mbox{}\par

\noindent\textbf{1：不適切} グリーン購入はバイオマスのみ。式・定義・対応の独立判定と一致しない。\par

\noindent\textbf{2：不適切} 環境報告は設備の自治体届出。式・定義・対応の独立判定と一致しない。\par

\noindent\textbf{3：適切・選択対象} 会計は費用効果を金額物量。判定と一致する。\par

\noindent\textbf{4：不適切} 環境監査は税務監査。式・定義・対応の独立判定と一致しない。\par

\noindent\textbf{5：不適切} LCAは設備の更新周期。式・定義・対応の独立判定と一致しない。\par

\paragraph{再利用する考え方} 目的・対象・単位による管理手法の区別。 集計対象・時点・法の主体・選択方向を固定し、同型でも元の五肢へ戻す。

\paragraph{共通理論} \hyperref[theory:env:accounting]{第\ref*{theory:env:accounting}節}

\par 参考：\cite{env58_accounting}。

\par 難易度 2（校準初稿）、優先度未評価。

\section{粒径範囲と大気汚染の生成経路}

\subsection{平成26年度 問5-1}
\label{q:h26-reg-basic-5-1}

\noindent\textbf{原典：} sources/H26/basic.pdf、PDF物理16頁。

\textbf{公式正答：} 1（single）。

\par\textbf{検討未確定：} COによるメタン寿命の化学過程と各健康影響は追加一次資料照合。肺沈着と粒径の関係を2.5μm一値へ固定しない。

\paragraph{問題の要旨} 粒径範囲と大気汚染の生成経路。

\Needspace{4\baselineskip}\paragraph{確認した条件・記述（独自要約）}\mbox{}\par

\noindent\textbf{時点：} 平成26年度に指定された資料・制度。\par

\noindent\textbf{照合：} 原典21画像で否定・五肢・組合せ表を確認。後年値で置換しない。\par

\paragraph{問われる核心} 粒径範囲と大気汚染の生成経路。

\paragraph{必要知識} 大きさの範囲と測定、一次・二次汚染物質、生成と除去。

\paragraph{考え方・正答を選ぶ理由（査読待ち初稿）} PM2.5は概ね粒径2.5μm以下の微小粒子で、2.5付近だけに限定する1が誤り。規格上は2.5μmを50\%分離する装置で採取するという測定定義もある。SO2は硫黄、NOxとVOCは二次生成、COは不完全燃焼とヘモグロビンの結合。COによるOH消費はメタン除去に影響する。

\Needspace{4\baselineskip}\paragraph{選択肢の検討（査読待ち初稿）}\mbox{}\par

\noindent\textbf{1：不適切・選択対象} PM2.5は2.5μm付近だけ。判定と一致する。\par

\noindent\textbf{2：条件内で適切} SO2は燃料中S由来。原典の条件内では適切な記述であり、不適切なものを選ぶ問の選択対象ではない。\par

\noindent\textbf{3：条件内で適切} NO2は燃焼と光化学前駆物質。原典の条件内では適切な記述であり、不適切なものを選ぶ問の選択対象ではない。\par

\noindent\textbf{4：条件内で適切} 光化学オキシダントはNOxとVOCの光反応。原典の条件内では適切な記述であり、不適切なものを選ぶ問の選択対象ではない。\par

\noindent\textbf{5：条件内で適切} COは不完全燃焼でメタン寿命にも影響。原典の条件内では適切な記述であり、不適切なものを選ぶ問の選択対象ではない。\par

\paragraph{再利用する考え方} 大きさの範囲と測定、一次・二次汚染物質、生成と除去。 指定年度・対象範囲・単位・否定を先に固定し、他年度の番号や現行値を転用しない。

\paragraph{共通理論} \hyperref[theory:env:air]{第\ref*{theory:env:air}節}

\par 参考：\cite{env58_pmstandard,env58_pm25}。

\par 難易度 3（校準初稿）、優先度未評価。

\subsection{令和5年度 問5-2}
\label{q:r05-reg-basic-5-2}

\noindent\textbf{原典：} sources/R05/basic.pdf、PDF物理16頁。

\textbf{公式正答：} 5（single）。

\paragraph{問題の要旨} PM2.5は2.5μm付近だけではない。

\Needspace{4\baselineskip}\paragraph{確認した条件・記述（独自要約）}\mbox{}\par

\noindent\textbf{原典：} 全五肢と否定方向・組合せ表を画像確認。\par

\noindent\textbf{資料時点：} 令和5年度の指定資料・統計年を保持し、後年値に置換しない。\par

\paragraph{問われる核心} PM2.5は2.5μm付近だけではない。

\paragraph{必要知識} PM分粒、一次排出と二次生成、SO2・NO2・CO・Ox。

\paragraph{考え方・正答を選ぶ理由（査読待ち初稿）} 5が不適切。2.5という値は代表粒径一種類を指すものではなく粒径区分である。各汚染物質は燃料由来・不完全燃焼・大気中の二次生成を分ける。

\Needspace{4\baselineskip}\paragraph{選択肢の検討（査読待ち初稿）}\mbox{}\par

\noindent\textbf{1：条件内で適切} SO2は硫黄燃料燃焼と酸性雨。発生源の硫黄と反応生成物を対応。\par

\noindent\textbf{2：条件内で適切} NOが酸化しNO2となる。燃焼で生じたNOの大気酸化はNO2の生成経路。\par

\noindent\textbf{3：条件内で適切} COは不完全燃焼・酸素運搬阻害。COがヘモグロビンと結合して酸素運搬を妨げる。\par

\noindent\textbf{4：条件内で適切} 光化学オキシダントはNOx・VOC等から二次生成。日射下の化学反応で生じる酸化性物質の総称。\par

\noindent\textbf{5：不適切・選択対象} PM2.5は粒径2.5μm付近のみ。微小粒子の区分は2.5μm以下側の粒径範囲であり、その近傍だけではない。実測は分粒特性で定義される。\par

\paragraph{再利用する考え方} PM分粒、一次排出と二次生成、SO2・NO2・CO・Ox。 否定方向と資料年、分母と単位を確認して五肢へ戻る。

\paragraph{共通理論} \hyperref[theory:env:air]{第\ref*{theory:env:air}節}

\par 参考：\cite{env58_pm25}。

\par 難易度 2（校準初稿）、優先度未評価。

\section{産業廃棄物フローを物質収支で検算}

\subsection{平成28年度 問5-1}
\label{q:h28-reg-basic-5-1}

\noindent\textbf{原典：} sources/H28/basic.pdf、PDF物理21頁。

\textbf{公式正答：} 5（single）。

\paragraph{問題の要旨} 産業廃棄物フローを物質収支で検算。

\Needspace{4\baselineskip}\paragraph{確認した条件・記述（独自要約）}\mbox{}\par

\noindent\textbf{排出：} 379百万t。\par

\noindent\textbf{中間処理：} 290。\par

\noindent\textbf{直接再生：} 83。\par

\noindent\textbf{処理後再生：} 124。\par

\noindent\textbf{処理後最終処分：} 7。\par

\noindent\textbf{全再生：} 207。\par

\noindent\textbf{全最終処分：} 13。\par

\paragraph{問われる核心} 産業廃棄物フローを物質収支で検算。

\paragraph{必要知識} 分岐の収支、直接と処理後、百分率と比率。

\paragraph{考え方・正答を選ぶ理由（査読待ち初稿）} 単位百万t。中間処理290は再利用124、処理後処分7、減量159へ分かれる。総処分13から直接処分は13−7=\allowbreak{}6で、60ではないので5。総排出も83+290+6=\allowbreak{}379、最終の207+159+13=\allowbreak{}379で一致。83/207≈.40、207/379≈.55、207/13≈16、124/83≈1.5、159/379≈.42。

\Needspace{4\baselineskip}\paragraph{選択肢の検討（査読待ち初稿）}\mbox{}\par

\noindent\textbf{1：条件内で適切} 直接再利用83で約40\%。原典の条件内では適切な記述であり、不適切なものを選ぶ問の選択対象ではない。\par

\noindent\textbf{2：条件内で適切} 再利用207で55\%・処分の16倍。原典の条件内では適切な記述であり、不適切なものを選ぶ問の選択対象ではない。\par

\noindent\textbf{3：条件内で適切} 処理後124で直接の1.5倍。原典の条件内では適切な記述であり、不適切なものを選ぶ問の選択対象ではない。\par

\noindent\textbf{4：条件内で適切} 減量159で42\%。原典の条件内では適切な記述であり、不適切なものを選ぶ問の選択対象ではない。\par

\noindent\textbf{5：不適切・選択対象} 直接処分60で16\%。判定と一致する。\par

\paragraph{再利用する考え方} 分岐の収支、直接と処理後、百分率と比率。 指定年度・対象範囲・単位・否定を先に固定し、他年度の番号や現行値を転用しない。

\paragraph{共通理論} \hyperref[theory:env:balance]{第\ref*{theory:env:balance}節}

\paragraph{関連問題} \hyperref[q:r06-reg-basic-5-2]{R06-REG-BASIC-5-2}

\par 難易度 2（校準初稿）、優先度未評価。

\subsection{令和6年度 問5-2}
\label{q:r06-reg-basic-5-2}

\noindent\textbf{原典：} sources/R06/basic.pdf、PDF物理18頁。

\textbf{公式正答：} 5（single）。

\paragraph{問題の要旨} 2020産業廃棄物フローの再生比は約53\%。

\Needspace{4\baselineskip}\paragraph{確認した条件・記述（独自要約）}\mbox{}\par

\noindent\textbf{単位：} 百万トン。\par

\noindent\textbf{原図：} 2020年度の概算値。減量を再生や処分へ加えない。\par

\paragraph{問われる核心} 2020産業廃棄物フローの再生比は約53\%。

\paragraph{必要知識} フロー保存、枝の量と総量、比の分母、複合命題。

\paragraph{考え方・正答を選ぶ理由（査読待ち初稿）} 5が不適切。図は排出374、中間293、直接再生77、中間後再生122、中間後処分5、総再生199、総処分9（百万t）。減量166、直接処分4を保存則から求める。374=\allowbreak{}199+9+166。5の二命題を分け、比の後半だけで適切としない。

\Needspace{4\baselineskip}\paragraph{選択肢の検討（査読待ち初稿）}\mbox{}\par

\noindent\textbf{1：条件内で適切} 直接再生77は再生199の約40\%。77/199≈38.7\%。約40\%に整合。\par

\noindent\textbf{2：条件内で適切} 中間後再生122は処理293の約40\%。122/293≈41.6\%。概算40\%に整合。\par

\noindent\textbf{3：条件内で適切} 減量166は排出374の約45\%。293−122−5=\allowbreak{}166、166/374≈44.4\%。\par

\noindent\textbf{4：条件内で適切} 直接最終4は排出374の約1\%。9−5=\allowbreak{}4、4/374≈1.07\%。\par

\noindent\textbf{5：不適切・選択対象} 再生199は排出374の約20\%・最終9の約22倍。199/374≈53.2\%であり20\%ではない。199/9≈22.1倍の後半が正しくても一肢全体は不適切。\par

\paragraph{再利用する考え方} フロー保存、枝の量と総量、比の分母、複合命題。 否定方向と資料年、分母と単位を確認して五肢へ戻る。

\paragraph{共通理論} \hyperref[theory:env:balance]{第\ref*{theory:env:balance}節}

\paragraph{関連問題} \hyperref[q:h28-reg-basic-5-1]{H28-REG-BASIC-5-1}

\par 難易度 2（校準初稿）、優先度未評価。

\section{生物多様性条約と外来種防除}

\subsection{平成28年度 問5-2}
\label{q:h28-reg-basic-5-2}

\noindent\textbf{原典：} sources/H28/basic.pdf、PDF物理22頁。

\textbf{公式正答：} 5（single）。

\par\textbf{検討未確定：} 1992・1993と議定書、当時の外来生物法の防除条項の一次本文を追加照合。動物愛護上の配慮が不要という意味にはしない。

\paragraph{問題の要旨} 生物多様性条約と外来種防除。

\Needspace{4\baselineskip}\paragraph{確認した条件・記述（独自要約）}\mbox{}\par

\noindent\textbf{時点：} 平成28年度に指定された資料・制度。\par

\noindent\textbf{照合：} 原典21画像で否定・五肢・組合せ表を確認。後年値で置換しない。\par

\paragraph{問われる核心} 生物多様性条約と外来種防除。

\paragraph{必要知識} 条約目的、署名と発効、外来種と防除。

\paragraph{考え方・正答を選ぶ理由（査読待ち初稿）} 外来種による生態系被害の防止には防除制度があり、一律に駆除を禁止する5が誤り。CBDの三目的は保全、持続可能な利用、遺伝資源の利益の公正衡平配分。組換え生物の扱いはCartagena議定書であり、遺伝資源の利益配分のNagoya議定書と混同しない。

\Needspace{4\baselineskip}\paragraph{選択肢の検討（査読待ち初稿）}\mbox{}\par

\noindent\textbf{1：条件内で適切} 組換え生物の議定書。原典の条件内では適切な記述であり、不適切なものを選ぶ問の選択対象ではない。\par

\noindent\textbf{2：条件内で適切} 1992署名開放・1993発効。原典の条件内では適切な記述であり、不適切なものを選ぶ問の選択対象ではない。\par

\noindent\textbf{3：条件内で適切} 保全・持続利用・利益配分。原典の条件内では適切な記述であり、不適切なものを選ぶ問の選択対象ではない。\par

\noindent\textbf{4：条件内で適切} 外来種に重大影響の例。原典の条件内では適切な記述であり、不適切なものを選ぶ問の選択対象ではない。\par

\noindent\textbf{5：不適切・選択対象} 日本は動物愛護で駆除禁止。判定と一致する。\par

\paragraph{再利用する考え方} 条約目的、署名と発効、外来種と防除。 指定年度・対象範囲・単位・否定を先に固定し、他年度の番号や現行値を転用しない。

\paragraph{共通理論} \hyperref[theory:env:biodiversity]{第\ref*{theory:env:biodiversity}節}

\par 難易度 3（校準初稿）、優先度未評価。

\subsection{令和2年度 問5-2}
\label{q:r02-reg-basic-5-2}

\noindent\textbf{原典：} sources/R02/basic.pdf、PDF物理20頁。

\textbf{公式正答：} 3（single）。

\par\textbf{検討未確定：} 出題時の防除条項とCBD全三目的の原条文照合は追加確認。

\paragraph{問題の要旨} 生物多様性保全と外来種防除。

\Needspace{4\baselineskip}\paragraph{確認した条件・記述（独自要約）}\mbox{}\par

\noindent\textbf{時点：} 令和2年度に指定された資料・制度。\par

\noindent\textbf{照合：} 原典18画像で否定・五肢・組合せ表を確認。後年値で置換しない。\par

\paragraph{問われる核心} 生物多様性保全と外来種防除。

\paragraph{必要知識} 条約と議定書、外来種、防除、倫理配慮。

\paragraph{考え方・正答を選ぶ理由（査読待ち初稿）} 3の一律駆除禁止が不適切。防除制度と扱いへの配慮を併せて考える。CBDの三目的とCartagenaの組換え生物の扱いは対象が異なり、Nagoyaの利益配分とも区別する。

\Needspace{4\baselineskip}\paragraph{選択肢の検討（査読待ち初稿）}\mbox{}\par

\noindent\textbf{1：条件内で適切} 組換え生物の国際枠組。条件内で適切な記述であり、不適切な記述を選ぶ対象ではない。\par

\noindent\textbf{2：条件内で適切} 外来種が在来種に影響。条件内で適切な記述であり、不適切な記述を選ぶ対象ではない。\par

\noindent\textbf{3：不適切・選択対象} 駆除は動物愛護で禁止。判定と一致する。\par

\noindent\textbf{4：条件内で適切} 1992署名開放1993発効。条件内で適切な記述であり、不適切な記述を選ぶ対象ではない。\par

\noindent\textbf{5：条件内で適切} 保全持続利用利益配分。条件内で適切な記述であり、不適切な記述を選ぶ対象ではない。\par

\paragraph{再利用する考え方} 条約と議定書、外来種、防除、倫理配慮。 集計対象・時点・法の主体・選択方向を固定し、同型でも元の五肢へ戻す。

\paragraph{共通理論} \hyperref[theory:env:biodiversity]{第\ref*{theory:env:biodiversity}節}

\par 難易度 3（校準初稿）、優先度未評価。

\section{2008年度の温室効果ガスと国際枠組}

\subsection{平成23年度 問5-1}
\label{q:h23-reg-basic-5-1}

\noindent\textbf{原典：} sources/H23/basic.pdf、PDF物理15頁。

\textbf{公式正答：} 1（single）。

\par\textbf{検討未確定：} 当時白書の全ガス内訳・部門図と2009合意の条文の追加照合が必要。CO2九割以上は当時確定値公表表の検索本文でも確認。

\paragraph{問題の要旨} 2008年度の温室効果ガスと国際枠組。

\Needspace{4\baselineskip}\paragraph{確認した条件・記述（独自要約）}\mbox{}\par

\noindent\textbf{時点：} 平成23年度に指定された資料・制度。\par

\noindent\textbf{照合：} 原典21画像で否定・五肢・組合せ表を確認。後年値で置換しない。\par

\paragraph{問われる核心} 2008年度の温室効果ガスと国際枠組。

\paragraph{必要知識} CO2換算、部門の定義、条約と議定書、採択年。

\paragraph{考え方・正答を選ぶ理由（査読待ち初稿）} CO2換算の2008年度総量の九割以上がCO2で、1。部門別では家庭が第二位という組合せが違う。UNFCCCの究極目的に2050という一律期限はなく、1997採択は京都議定書。2009コペンハーゲンの合意を全途上国に法的数量削減義務を課す議定書と扱わない。

\Needspace{4\baselineskip}\paragraph{選択肢の検討（査読待ち初稿）}\mbox{}\par

\noindent\textbf{1：適切・選択対象} CO2が九割以上。判定と一致する。\par

\noindent\textbf{2：不適切} 産業の次は家庭。式・定義・対応の独立判定と一致しない。\par

\noindent\textbf{3：不適切} 条約の濃度安定は2050期限。式・定義・対応の独立判定と一致しない。\par

\noindent\textbf{4：不適切} 1997パリ議定書。式・定義・対応の独立判定と一致しない。\par

\noindent\textbf{5：不適切} 2009途上国に基準年削減数値義務。式・定義・対応の独立判定と一致しない。\par

\paragraph{再利用する考え方} CO2換算、部門の定義、条約と議定書、採択年。 指定年度・対象範囲・単位・否定を先に固定し、他年度の番号や現行値を転用しない。

\paragraph{共通理論} \hyperref[theory:env:climate]{第\ref*{theory:env:climate}節}

\par 参考：\cite{env58_ghg2008}。

\par 難易度 3（校準初稿）、優先度未評価。

\section{暖房と自動車のCO2排出量を比較}

\subsection{平成25年度 問5-2}
\label{q:h25-reg-basic-5-2}

\noindent\textbf{原典：} sources/H25/basic.pdf、PDF物理19頁。

\textbf{公式正答：} 2（single）。

\paragraph{問題の要旨} 暖房と自動車のCO2排出量を比較。

\Needspace{4\baselineskip}\paragraph{確認した条件・記述（独自要約）}\mbox{}\par

\noindent\textbf{電力：} 500W、3h/日×50日、.42kg-CO2/kWh。\par

\noindent\textbf{車：} 10km/L、200km、2.32kg-CO2/L。\par

\paragraph{問われる核心} 暖房と自動車のCO2排出量を比較。

\paragraph{必要知識} エネルギー=電力×時間、燃料使用量、排出係数。

\paragraph{考え方・正答を選ぶ理由（査読待ち初稿）} 暖房は.5kW×3h×50=\allowbreak{}75kWh、排出31.5kg。自動車は200/10=\allowbreak{}20L、排出46.4kg。比31.5/46.4≈.679で2:3に最も近い。WをkWへ直してから時間を掛ける。

\Needspace{4\baselineskip}\paragraph{選択肢の検討（査読待ち初稿）}\mbox{}\par

\noindent\textbf{1：不適切} 1:3。式・定義・対応の独立判定と一致しない。\par

\noindent\textbf{2：適切・選択対象} 2:3。判定と一致する。\par

\noindent\textbf{3：不適切} 1:1。式・定義・対応の独立判定と一致しない。\par

\noindent\textbf{4：不適切} 3:2。式・定義・対応の独立判定と一致しない。\par

\noindent\textbf{5：不適切} 3:1。式・定義・対応の独立判定と一致しない。\par

\paragraph{再利用する考え方} エネルギー=電力×時間、燃料使用量、排出係数。 指定年度・対象範囲・単位・否定を先に固定し、他年度の番号や現行値を転用しない。

\paragraph{共通理論} \hyperref[theory:env:co2]{第\ref*{theory:env:co2}節}

\par 難易度 2（校準初稿）、優先度未評価。

\subsection{平成27年度 問5-3}
\label{q:h27-reg-basic-5-3}

\noindent\textbf{原典：} sources/H27/basic.pdf、PDF物理16頁。

\textbf{公式正答：} 5（single）。

\paragraph{問題の要旨} 利用人数と時間を含めた一人当たり排出。

\Needspace{4\baselineskip}\paragraph{確認した条件・記述（独自要約）}\mbox{}\par

\noindent\textbf{a：} 10km/L、100km、3人、2.32kg/L。\par

\noindent\textbf{b：} 8000kW、200km/h、500km、1000人、.47kg/kWh。\par

\noindent\textbf{c：} 500W、4h/日、50日、4人、.47kg/kWh。\par

\paragraph{問われる核心} 利用人数と時間を含めた一人当たり排出。

\paragraph{必要知識} 排出係数、乗車人数、走行時間、総量と原単位。

\paragraph{考え方・正答を選ぶ理由（査読待ち初稿）} aは100/10×2.32/3=\allowbreak{}7.73kg/人。bは500/200=\allowbreak{}2.5h、8000×2.5×.47/1000=\allowbreak{}9.4kg/人。cは.5×4×50×.47/4=\allowbreak{}11.75kg/人。c>b>aの5。総量の大小ではなく人数で割った量を比べる。

\Needspace{4\baselineskip}\paragraph{選択肢の検討（査読待ち初稿）}\mbox{}\par

\noindent\textbf{1：不適切} a>b>c。式・定義・対応の独立判定と一致しない。\par

\noindent\textbf{2：不適切} a>c>b。式・定義・対応の独立判定と一致しない。\par

\noindent\textbf{3：不適切} b>a>c。式・定義・対応の独立判定と一致しない。\par

\noindent\textbf{4：不適切} c>a>b。式・定義・対応の独立判定と一致しない。\par

\noindent\textbf{5：適切・選択対象} c>b>a。判定と一致する。\par

\paragraph{再利用する考え方} 排出係数、乗車人数、走行時間、総量と原単位。 指定年度・対象範囲・単位・否定を先に固定し、他年度の番号や現行値を転用しない。

\paragraph{共通理論} \hyperref[theory:env:co2]{第\ref*{theory:env:co2}節}

\par 難易度 2（校準初稿）、優先度未評価。

\section{科学と社会をつなぐ双方向の対話}

\subsection{平成25年度 問5-5}
\label{q:h25-reg-basic-5-5}

\noindent\textbf{原典：} sources/H25/basic.pdf、PDF物理22頁。

\textbf{公式正答：} 2（single）。

\paragraph{問題の要旨} 科学と社会をつなぐ双方向の対話。

\Needspace{4\baselineskip}\paragraph{確認した条件・記述（独自要約）}\mbox{}\par

\noindent\textbf{四文の順序：} ア=社会理解、イ=科学技術間だけ、ウ=双方向、エ=媒体へ技術使用\par

\noindent\textbf{問：} 適切な二文。\par

\paragraph{問われる核心} 科学と社会をつなぐ双方向の対話。

\paragraph{必要知識} 専門家と社会、理解と対話、情報発信と相互学習。

\paragraph{考え方・正答を選ぶ理由（査読待ち初稿）} 社会・国民の期待や考えを専門家自身が理解することと、専門家・非専門家間の双方向性が適切。基礎科学と応用技術の内部連携だけをこの領域全体の定義にはできず、媒体制作へ技術を使うことや科学者へ任せることも、社会と対話するという目的の代替ではない。原典の列順で組合せを選ぶ。

\Needspace{4\baselineskip}\paragraph{選択肢の検討（査読待ち初稿）}\mbox{}\par

\noindent\textbf{1：不適切} ア イ。式・定義・対応の独立判定と一致しない。\par

\noindent\textbf{2：適切・選択対象} ア ウ。判定と一致する。\par

\noindent\textbf{3：不適切} イ ウ。式・定義・対応の独立判定と一致しない。\par

\noindent\textbf{4：不適切} イ エ。式・定義・対応の独立判定と一致しない。\par

\noindent\textbf{5：不適切} ウ エ。式・定義・対応の独立判定と一致しない。\par

\paragraph{再利用する考え方} 専門家と社会、理解と対話、情報発信と相互学習。 指定年度・対象範囲・単位・否定を先に固定し、他年度の番号や現行値を転用しない。

\paragraph{共通理論} \hyperref[theory:env:communication]{第\ref*{theory:env:communication}節}

\par 参考：\cite{env58_communication}。

\par 難易度 2（校準初稿）、優先度未評価。

\subsection{平成28年度 問5-5}
\label{q:h28-reg-basic-5-5}

\noindent\textbf{原典：} sources/H28/basic.pdf、PDF物理24頁。

\textbf{公式正答：} 5（single）。

\paragraph{問題の要旨} 科学と社会をつなぐ双方向の対話。

\Needspace{4\baselineskip}\paragraph{確認した条件・記述（独自要約）}\mbox{}\par

\noindent\textbf{四文の順序：} ア=科学技術間だけ、イ=科学者に制作委任、ウ=社会理解、エ=双方向\par

\noindent\textbf{問：} 適切な二文。\par

\paragraph{問われる核心} 科学と社会をつなぐ双方向の対話。

\paragraph{必要知識} 専門家と社会、理解と対話、情報発信と相互学習。

\paragraph{考え方・正答を選ぶ理由（査読待ち初稿）} 社会・国民の期待や考えを専門家自身が理解することと、専門家・非専門家間の双方向性が適切。基礎科学と応用技術の内部連携だけをこの領域全体の定義にはできず、媒体制作へ技術を使うことや科学者へ任せることも、社会と対話するという目的の代替ではない。原典の列順で組合せを選ぶ。

\Needspace{4\baselineskip}\paragraph{選択肢の検討（査読待ち初稿）}\mbox{}\par

\noindent\textbf{1：不適切} ア イ。式・定義・対応の独立判定と一致しない。\par

\noindent\textbf{2：不適切} ア ウ。式・定義・対応の独立判定と一致しない。\par

\noindent\textbf{3：不適切} イ ウ。式・定義・対応の独立判定と一致しない。\par

\noindent\textbf{4：不適切} イ エ。式・定義・対応の独立判定と一致しない。\par

\noindent\textbf{5：適切・選択対象} ウ エ。判定と一致する。\par

\paragraph{再利用する考え方} 専門家と社会、理解と対話、情報発信と相互学習。 指定年度・対象範囲・単位・否定を先に固定し、他年度の番号や現行値を転用しない。

\paragraph{共通理論} \hyperref[theory:env:communication]{第\ref*{theory:env:communication}節}

\par 参考：\cite{env58_communication}。

\par 難易度 2（校準初稿）、優先度未評価。

\section{2010年の電源比とガス技術を分ける}

\subsection{平成24年度 問5-3}
\label{q:h24-reg-basic-5-3}

\noindent\textbf{原典：} sources/H24/basic.pdf、PDF物理16頁。

\textbf{公式正答：} 2（single）。

\par\textbf{検討未確定：} 一般電気事業用という分母の2010年度原表と2005以降の米国生産統計は追加取得。現行の全発電量比率と混同しない。

\paragraph{問題の要旨} 2010年の電源比とガス技術を分ける。

\Needspace{4\baselineskip}\paragraph{確認した条件・記述（独自要約）}\mbox{}\par

\noindent\textbf{時点：} 平成24年度に指定された資料・制度。\par

\noindent\textbf{照合：} 原典21画像で否定・五肢・組合せ表を確認。後年値で置換しない。\par

\paragraph{問われる核心} 2010年の電源比とガス技術を分ける。

\paragraph{必要知識} 歴史的電源構成、複合サイクル、非在来型天然ガス。

\paragraph{考え方・正答を選ぶ理由（査読待ち初稿）} 2010年度の原子力比率は約30\%。天然ガスのコンバインドサイクルはガスタービンと排熱による蒸気タービンを組み合わせる。米国で増産が注目された非在来ガスはシェールガス。三語の組合せ2。

\Needspace{4\baselineskip}\paragraph{選択肢の検討（査読待ち初稿）}\mbox{}\par

\noindent\textbf{1：不適切} 30/再熱再生/タイトサンド。式・定義・対応の独立判定と一致しない。\par

\noindent\textbf{2：適切・選択対象} 30/コンバインド/シェール。判定と一致する。\par

\noindent\textbf{3：不適切} 30/コンバインド/タイトサンド。式・定義・対応の独立判定と一致しない。\par

\noindent\textbf{4：不適切} 20/再熱再生/シェール。式・定義・対応の独立判定と一致しない。\par

\noindent\textbf{5：不適切} 20/コンバインド/シェール。式・定義・対応の独立判定と一致しない。\par

\paragraph{再利用する考え方} 歴史的電源構成、複合サイクル、非在来型天然ガス。 指定年度・対象範囲・単位・否定を先に固定し、他年度の番号や現行値を転用しない。

\paragraph{共通理論} \hyperref[theory:env:energy2010]{第\ref*{theory:env:energy2010}節}

\par 難易度 3（校準初稿）、優先度未評価。

\subsection{令和2年度 問5-4}
\label{q:r02-reg-basic-5-4}

\noindent\textbf{原典：} sources/R02/basic.pdf、PDF物理21頁。

\textbf{公式正答：} 1（single）。

\par\textbf{検討未確定：} 2010の指定分母の原表と米国由来輸入の開始資料は未取得。H24と同じ番号でないことを確認済み。

\paragraph{問題の要旨} 2010原子力比と複合サイクル・シェール。

\Needspace{4\baselineskip}\paragraph{確認した条件・記述（独自要約）}\mbox{}\par

\noindent\textbf{時点：} 令和2年度に指定された資料・制度。\par

\noindent\textbf{照合：} 原典18画像で否定・五肢・組合せ表を確認。後年値で置換しない。\par

\paragraph{問われる核心} 2010原子力比と複合サイクル・シェール。

\paragraph{必要知識} 歴史的電源比、複合サイクル、ガスの供給。

\paragraph{考え方・正答を選ぶ理由（査読待ち初稿）} 2010一般電気事業用原子力約30\%、ガスタービンと排熱蒸気のコンバインド、米国非在来型シェールで1。原子力の指定時点と、輸出への期待から実際のLNG輸入へ変わった本文の時点を別々に読む。

\Needspace{4\baselineskip}\paragraph{選択肢の検討（査読待ち初稿）}\mbox{}\par

\noindent\textbf{1：適切・選択対象} 30/コンバインド/シェール。判定と一致する。\par

\noindent\textbf{2：不適切} 20/コンバインド/シェール。式・定義・対応の独立判定と一致しない。\par

\noindent\textbf{3：不適切} 20/再熱再生/シェール。式・定義・対応の独立判定と一致しない。\par

\noindent\textbf{4：不適切} 30/コンバインド/タイトサンド。式・定義・対応の独立判定と一致しない。\par

\noindent\textbf{5：不適切} 30/再熱再生/タイトサンド。式・定義・対応の独立判定と一致しない。\par

\paragraph{再利用する考え方} 歴史的電源比、複合サイクル、ガスの供給。 集計対象・時点・法の主体・選択方向を固定し、同型でも元の五肢へ戻す。

\paragraph{共通理論} \hyperref[theory:env:energy2010]{第\ref*{theory:env:energy2010}節}

\par 難易度 3（校準初稿）、優先度未評価。

\section{2011年一人当たり一次エネルギー}

\subsection{平成27年度 問5-4}
\label{q:h27-reg-basic-5-4}

\noindent\textbf{原典：} sources/H27/basic.pdf、PDF物理17頁。

\textbf{公式正答：} 5（single）。

\par\textbf{検討未確定：} IEAの指定2011年国別原表は未取得。現行データページを当時値の照合証跡にはせず、Astraの追加確認対象に残す。

\paragraph{問題の要旨} 2011年一人当たり一次エネルギー。

\Needspace{4\baselineskip}\paragraph{確認した条件・記述（独自要約）}\mbox{}\par

\noindent\textbf{時点：} 平成27年度に指定された資料・制度。\par

\noindent\textbf{照合：} 原典21画像で否定・五肢・組合せ表を確認。後年値で置換しない。\par

\paragraph{問われる核心} 2011年一人当たり一次エネルギー。

\paragraph{必要知識} 一次エネルギー、石油換算、人口当たり、年次比較。

\paragraph{考え方・正答を選ぶ理由（査読待ち初稿）} 原典の2011年指標では米加が7toe以上、韓露約5toe、独日約4toeで5。1toe≈42GJを使い、国全体の総量と人口当たりを混同しない。世界1.9、中国2.0という同じ年・同じ指標で比較する。

\Needspace{4\baselineskip}\paragraph{選択肢の検討（査読待ち初稿）}\mbox{}\par

\noindent\textbf{1：不適切} 独日/米加/韓露。式・定義・対応の独立判定と一致しない。\par

\noindent\textbf{2：不適切} 韓露/独日/米加。式・定義・対応の独立判定と一致しない。\par

\noindent\textbf{3：不適切} 韓露/米加/独日。式・定義・対応の独立判定と一致しない。\par

\noindent\textbf{4：不適切} 米加/独日/韓露。式・定義・対応の独立判定と一致しない。\par

\noindent\textbf{5：適切・選択対象} 米加/韓露/独日。判定と一致する。\par

\paragraph{再利用する考え方} 一次エネルギー、石油換算、人口当たり、年次比較。 指定年度・対象範囲・単位・否定を先に固定し、他年度の番号や現行値を転用しない。

\paragraph{共通理論} \hyperref[theory:env:energycapita]{第\ref*{theory:env:energycapita}節}

\par 難易度 3（校準初稿）、優先度未評価。

\subsection{令和3年度 問5-4}
\label{q:r03-reg-basic-5-4}

\noindent\textbf{原典：} sources/R03/basic.pdf、PDF物理24頁。

\textbf{公式正答：} 2（single）。

\par\textbf{検討未確定：} IEA2018国別値の元統計表は追加確認対象。既知の大国総量から値を逆算しない。

\paragraph{問題の要旨} 2018一人当たり一次エネルギーの国群。

\Needspace{4\baselineskip}\paragraph{確認した条件・記述（独自要約）}\mbox{}\par

\noindent\textbf{原典：} 全五肢と否定方向・組合せ表を画像確認。\par

\noindent\textbf{資料時点：} 令和3年度の指定資料・統計年を保持し、後年値に置換しない。\par

\paragraph{問われる核心} 2018一人当たり一次エネルギーの国群。

\paragraph{必要知識} 石油換算トン約42GJ、一人当たりと総量、同年比較。

\paragraph{考え方・正答を選ぶ理由（査読待ち初稿）} 指定されたIEA2018の一人当たり一次供給では、6石油換算トン以上が米国・カナダ、5〜6が韓国・ロシア、3〜4がドイツ・日本。世界1.9、中国2.3と同じ定義・年で比較する。国全体の総量順位や電力だけの順位ではなく、人口で割った一次供給であり2。

\Needspace{4\baselineskip}\paragraph{選択肢の検討（査読待ち初稿）}\mbox{}\par

\noindent\textbf{1：不適切} 米加/独日/韓露。独立に求めた組合せ 米加/韓露/独日 と異なる。\par

\noindent\textbf{2：適切・選択対象} 米加/韓露/独日。独立に求めた組合せ 米加/韓露/独日 と一致する。\par

\noindent\textbf{3：不適切} 独日/米加/韓露。独立に求めた組合せ 米加/韓露/独日 と異なる。\par

\noindent\textbf{4：不適切} 韓露/独日/米加。独立に求めた組合せ 米加/韓露/独日 と異なる。\par

\noindent\textbf{5：不適切} 韓露/米加/独日。独立に求めた組合せ 米加/韓露/独日 と異なる。\par

\paragraph{再利用する考え方} 石油換算トン約42GJ、一人当たりと総量、同年比較。 否定方向と資料年、分母と単位を確認して五肢へ戻る。

\paragraph{共通理論} \hyperref[theory:env:energycapita]{第\ref*{theory:env:energycapita}節}

\par 難易度 3（校準初稿）、優先度未評価。

\section{2015年策定の2030電源見通しを読む}

\subsection{平成28年度 問5-4}
\label{q:h28-reg-basic-5-4}

\noindent\textbf{原典：} sources/H28/basic.pdf、PDF物理23頁。

\textbf{公式正答：} 1（single）。

\par\textbf{検討未確定：} 2015公開関連資料の検索と白書骨子を参照。策定7月の最終PDFの全図表照合は未了で後工程に残す。

\paragraph{問題の要旨} 2015年策定の2030電源見通しを読む。

\Needspace{4\baselineskip}\paragraph{確認した条件・記述（独自要約）}\mbox{}\par

\noindent\textbf{時点：} 平成28年度に指定された資料・制度。\par

\noindent\textbf{照合：} 原典21画像で否定・五肢・組合せ表を確認。後年値で置換しない。\par

\paragraph{問われる核心} 2015年策定の2030電源見通しを読む。

\paragraph{必要知識} 策定年と目標年、電源構成と自給率、見通し。

\paragraph{考え方・正答を選ぶ理由（査読待ち初稿）} 2015の見通しの2030構成は、25–27\%がLNGで石油は約3\%。従って1が不適切。再エネ22–24\%、原子力20–22\%と、需給の省エネ・自給率目標は当時見通しの範囲。2030を対象とする後年改訂計画の比率をこの問題へ遡及しない。

\Needspace{4\baselineskip}\paragraph{選択肢の検討（査読待ち初稿）}\mbox{}\par

\noindent\textbf{1：不適切・選択対象} 石油25–27\%。判定と一致する。\par

\noindent\textbf{2：条件内で適切} 再エネ22–24\%。原典の条件内では適切な記述であり、不適切なものを選ぶ問の選択対象ではない。\par

\noindent\textbf{3：条件内で適切} 原子力20–22\%。原典の条件内では適切な記述であり、不適切なものを選ぶ問の選択対象ではない。\par

\noindent\textbf{4：条件内で適切} 自給率約25\%と準国産。原典の条件内では適切な記述であり、不適切なものを選ぶ問の選択対象ではない。\par

\noindent\textbf{5：条件内で適切} 省エネで最終消費低下。原典の条件内では適切な記述であり、不適切なものを選ぶ問の選択対象ではない。\par

\paragraph{再利用する考え方} 策定年と目標年、電源構成と自給率、見通し。 指定年度・対象範囲・単位・否定を先に固定し、他年度の番号や現行値を転用しない。

\paragraph{共通理論} \hyperref[theory:env:energymix2015]{第\ref*{theory:env:energymix2015}節}

\par 参考：\cite{env58_mix2015}。

\par 難易度 3（校準初稿）、優先度未評価。

\subsection{令和元年度 問5-3}
\label{q:r01-reg-basic-5-3}

\noindent\textbf{原典：} sources/R01/basic.pdf、PDF物理20頁。

\textbf{公式正答：} 3（single）。

\par\textbf{検討未確定：} 2015の7月最終PDF全図表をまだ取得していない。原典と公開骨子に基づく初稿として追加照合。

\paragraph{問題の要旨} 2015策定の2030比率で燃料名を点検。

\Needspace{4\baselineskip}\paragraph{確認した条件・記述（独自要約）}\mbox{}\par

\noindent\textbf{時点：} 令和元年度に指定された資料・制度。\par

\noindent\textbf{照合：} 原典18画像で否定・五肢・組合せ表を確認。後年値で置換しない。\par

\paragraph{問われる核心} 2015策定の2030比率で燃料名を点検。

\paragraph{必要知識} 策定年・目標年・実績と見通し・電源分類。

\paragraph{考え方・正答を選ぶ理由（査読待ち初稿）} 25–27\%はLNG火力の目標であり、石油火力は約3\%。燃料名の誤りの3を選ぶ。策定時点の他の比率・省エネ・自給率を保持し、後年の目標改訂値へ置換しない。

\Needspace{4\baselineskip}\paragraph{選択肢の検討（査読待ち初稿）}\mbox{}\par

\noindent\textbf{1：条件内で適切} 原子力20–22。条件内で適切な記述であり、不適切な記述を選ぶ対象ではない。\par

\noindent\textbf{2：条件内で適切} 再エネ22–24。条件内で適切な記述であり、不適切な記述を選ぶ対象ではない。\par

\noindent\textbf{3：不適切・選択対象} 石油25–27。判定と一致する。\par

\noindent\textbf{4：条件内で適切} 2013より最終消費減。条件内で適切な記述であり、不適切な記述を選ぶ対象ではない。\par

\noindent\textbf{5：条件内で適切} 自給率25程度と準国産。条件内で適切な記述であり、不適切な記述を選ぶ対象ではない。\par

\paragraph{再利用する考え方} 策定年・目標年・実績と見通し・電源分類。 集計対象・時点・法の主体・選択方向を固定し、同型でも元の五肢へ戻す。

\paragraph{共通理論} \hyperref[theory:env:energymix2015]{第\ref*{theory:env:energymix2015}節}

\par 参考：\cite{env58_mix2015}。

\par 難易度 3（校準初稿）、優先度未評価。

\section{専門職責任は職務規定だけでは尽くせない}

\subsection{平成26年度 問5-5}
\label{q:h26-reg-basic-5-5}

\noindent\textbf{原典：} sources/H26/basic.pdf、PDF物理17頁。

\textbf{公式正答：} 3（single）。

\paragraph{問題の要旨} 専門職責任は職務規定だけでは尽くせない。

\Needspace{4\baselineskip}\paragraph{確認した条件・記述（独自要約）}\mbox{}\par

\noindent\textbf{時点：} 平成26年度に指定された資料・制度。\par

\noindent\textbf{照合：} 原典21画像で否定・五肢・組合せ表を確認。後年値で置換しない。\par

\paragraph{問われる核心} 専門職責任は職務規定だけでは尽くせない。

\paragraph{必要知識} 力量、公益、説明責任、法と倫理の範囲。

\paragraph{考え方・正答を選ぶ理由（査読待ち初稿）} 規定に明示がないことを責任免除の根拠とする3が不適切。専門職は力量の範囲、安全・公益、説明と結果への責任を考える。依頼業務が力量を超える場合、能力ある者との協働や適切な引継ぎを含め、無能力のまま責任を引き受けることを避ける。

\Needspace{4\baselineskip}\paragraph{選択肢の検討（査読待ち初稿）}\mbox{}\par

\noindent\textbf{1：条件内で適切} 職業で規範が異なる場合。原典の条件内では適切な記述であり、不適切なものを選ぶ問の選択対象ではない。\par

\noindent\textbf{2：条件内で適切} 力量外を引受けない。原典の条件内では適切な記述であり、不適切なものを選ぶ問の選択対象ではない。\par

\noindent\textbf{3：不適切・選択対象} 明文規定なしなら責任なし。判定と一致する。\par

\noindent\textbf{4：条件内で適切} 行動規範は変化。原典の条件内では適切な記述であり、不適切なものを選ぶ問の選択対象ではない。\par

\noindent\textbf{5：条件内で適切} 専門知識は高い基準。原典の条件内では適切な記述であり、不適切なものを選ぶ問の選択対象ではない。\par

\paragraph{再利用する考え方} 力量、公益、説明責任、法と倫理の範囲。 指定年度・対象範囲・単位・否定を先に固定し、他年度の番号や現行値を転用しない。

\paragraph{共通理論} \hyperref[theory:env:ethics]{第\ref*{theory:env:ethics}節}

\par 参考：\cite{ipej_ethics2011}。

\par 難易度 2（校準初稿）、優先度未評価。

\subsection{平成30年度 問5-6}
\label{q:h30-reg-basic-5-6}

\noindent\textbf{原典：} sources/H30/basic.pdf、PDF物理20頁。

\textbf{公式正答：} 5（single）。

\paragraph{問題の要旨} 倫理責任は明文化の有無だけでは決まらない。

\Needspace{4\baselineskip}\paragraph{確認した条件・記述（独自要約）}\mbox{}\par

\noindent\textbf{時点：} 平成30年度に指定された資料・制度。\par

\noindent\textbf{照合：} 原典18画像で否定・五肢・組合せ表を確認。後年値で置換しない。\par

\paragraph{問われる核心} 倫理責任は明文化の有無だけでは決まらない。

\paragraph{必要知識} 公益、力量、責任、明文と倫理。

\paragraph{考え方・正答を選ぶ理由（査読待ち初稿）} 明文規定がなくても専門職の公益・安全・説明に関する責任は残るので5。力量の範囲を見極め、能力ある者との協働や引継ぎで適切に業務を処理する。

\Needspace{4\baselineskip}\paragraph{選択肢の検討（査読待ち初稿）}\mbox{}\par

\noindent\textbf{1：条件内で適切} 専門家は高い基準。条件内で適切な記述であり、不適切な記述を選ぶ対象ではない。\par

\noindent\textbf{2：条件内で適切} 職業で規範相違。条件内で適切な記述であり、不適切な記述を選ぶ対象ではない。\par

\noindent\textbf{3：条件内で適切} 力量超え仕事を引受けない。条件内で適切な記述であり、不適切な記述を選ぶ対象ではない。\par

\noindent\textbf{4：条件内で適切} 規範は変わる。条件内で適切な記述であり、不適切な記述を選ぶ対象ではない。\par

\noindent\textbf{5：不適切・選択対象} 規定なしなら責任なし。判定と一致する。\par

\paragraph{再利用する考え方} 公益、力量、責任、明文と倫理。 集計対象・時点・法の主体・選択方向を固定し、同型でも元の五肢へ戻す。

\paragraph{共通理論} \hyperref[theory:env:ethics]{第\ref*{theory:env:ethics}節}

\par 参考：\cite{ipej_ethics2011}。

\par 難易度 2（校準初稿）、優先度未評価。

\section{2012開始のFITの原則と拒否事由}

\subsection{平成26年度 問5-3}
\label{q:h26-reg-basic-5-3}

\noindent\textbf{原典：} sources/H26/basic.pdf、PDF物理17頁。

\textbf{公式正答：} 5（single）。

\par\textbf{検討未確定：} 2013METI資料の検索本文で正当拒否事由を確認したがPDF全文取得は失敗。出題時価格決定と対象水力条件の詳細、法令履歴は追加照合。

\paragraph{問題の要旨} 2012開始のFITの原則と拒否事由。

\Needspace{4\baselineskip}\paragraph{確認した条件・記述（独自要約）}\mbox{}\par

\noindent\textbf{時点：} 平成26年度に指定された資料・制度。\par

\noindent\textbf{照合：} 原典21画像で否定・五肢・組合せ表を確認。後年値で置換しない。\par

\paragraph{問われる核心} 2012開始のFITの原則と拒否事由。

\paragraph{必要知識} 固定価格、買取と接続、賦課金、例外。

\paragraph{考え方・正答を選ぶ理由（査読待ち初稿）} 正当な理由がある場合まで接続拒否を一切できないとする5は誤り。2012開始の対象は太陽光・風力・水力・地熱・バイオマスで、調達価格の設定、使用量に応じた賦課、費用調整という仕組みを区別する。例外を削って無条件義務にしない。

\Needspace{4\baselineskip}\paragraph{選択肢の検討（査読待ち初稿）}\mbox{}\par

\noindent\textbf{1：条件内で適切} 五種類の再エネ。原典の条件内では適切な記述であり、不適切なものを選ぶ問の選択対象ではない。\par

\noindent\textbf{2：条件内で適切} 大臣が年度価格設定。原典の条件内では適切な記述であり、不適切なものを選ぶ問の選択対象ではない。\par

\noindent\textbf{3：条件内で適切} 電力量に応じ賦課。原典の条件内では適切な記述であり、不適切なものを選ぶ問の選択対象ではない。\par

\noindent\textbf{4：条件内で適切} 地域負担差を調整。原典の条件内では適切な記述であり、不適切なものを選ぶ問の選択対象ではない。\par

\noindent\textbf{5：不適切・選択対象} 接続は一切拒否不可。判定と一致する。\par

\paragraph{再利用する考え方} 固定価格、買取と接続、賦課金、例外。 指定年度・対象範囲・単位・否定を先に固定し、他年度の番号や現行値を転用しない。

\paragraph{共通理論} \hyperref[theory:env:fit]{第\ref*{theory:env:fit}節}

\par 参考：\cite{env58_fit2013}。

\par 難易度 3（校準初稿）、優先度未評価。

\section{複合発電の高温側から低温側への順}

\subsection{平成28年度 問5-3}
\label{q:h28-reg-basic-5-3}

\noindent\textbf{原典：} sources/H28/basic.pdf、PDF物理22頁。

\textbf{公式正答：} 3（single）。

\paragraph{問題の要旨} 複合発電の高温側から低温側への順。

\Needspace{4\baselineskip}\paragraph{確認した条件・記述（独自要約）}\mbox{}\par

\noindent\textbf{時点：} 平成28年度に指定された資料・制度。\par

\noindent\textbf{照合：} 原典21画像で否定・五肢・組合せ表を確認。後年値で置換しない。\par

\paragraph{問われる核心} 複合発電の高温側から低温側への順。

\paragraph{必要知識} 温度段階、電気化学、静電蓄積、エネルギー変換。

\paragraph{考え方・正答を選ぶ理由（査読待ち初稿）} コンバインドは燃焼ガスでガスタービンを回し、その高温排熱から蒸気を作って蒸気タービンへ進む。逆順の3が不適切。二次電池は再充電、EDLCは界面で蓄える。揚水は発電を消費側から切り離して電気を水の位置エネルギーへ変換する貯蔵で、発電時だけの効率を往復効率としない。

\Needspace{4\baselineskip}\paragraph{選択肢の検討（査読待ち初稿）}\mbox{}\par

\noindent\textbf{1：条件内で適切} 一次使切・二次充電。原典の条件内では適切な記述であり、不適切なものを選ぶ問の選択対象ではない。\par

\noindent\textbf{2：条件内で適切} EDLC急速・長サイクル。原典の条件内では適切な記述であり、不適切なものを選ぶ問の選択対象ではない。\par

\noindent\textbf{3：不適切・選択対象} 先に蒸気タービン次に廃熱ガスタービン。判定と一致する。\par

\noindent\textbf{4：条件内で適切} 燃料電池は水電解の逆。原典の条件内では適切な記述であり、不適切なものを選ぶ問の選択対象ではない。\par

\noindent\textbf{5：条件内で適切} 揚水は位置エネルギーに貯蔵。原典の条件内では適切な記述であり、不適切なものを選ぶ問の選択対象ではない。\par

\paragraph{再利用する考え方} 温度段階、電気化学、静電蓄積、エネルギー変換。 指定年度・対象範囲・単位・否定を先に固定し、他年度の番号や現行値を転用しない。

\paragraph{共通理論} \hyperref[theory:env:generation]{第\ref*{theory:env:generation}節}

\par 難易度 2（校準初稿）、優先度未評価。

\section{発熱量は指定表と単位質量で比較}

\subsection{平成24年度 問5-2}
\label{q:h24-reg-basic-5-2}

\noindent\textbf{原典：} sources/H24/basic.pdf、PDF物理15頁。

\textbf{公式正答：} 4（single）。

\par\textbf{検討未確定：} 2005年度指定表の全四値を同じ定義で取得する作業は未完了。2023改訂資料の歴史表を補助参照し、後年値へ置換しない。

\paragraph{問題の要旨} 発熱量は指定表と単位質量で比較。

\Needspace{4\baselineskip}\paragraph{確認した条件・記述（独自要約）}\mbox{}\par

\noindent\textbf{時点：} 平成24年度に指定された資料・制度。\par

\noindent\textbf{照合：} 原典21画像で否定・五肢・組合せ表を確認。後年値で置換しない。\par

\paragraph{問われる核心} 発熱量は指定表と単位質量で比較。

\paragraph{必要知識} 高位・低位発熱量、単位、燃料品種、密度換算。

\paragraph{考え方・正答を選ぶ理由（査読待ち初稿）} 指定された燃料種の質量当たり発熱量は天然ガスCが最大、原油A、輸入原料炭B、廃材Dの順なので4。天然ガスの表が体積単位なら密度でMJ/kgへ換算する。液体のMJ/Lを質量当たり値へ直接並べない。

\Needspace{4\baselineskip}\paragraph{選択肢の検討（査読待ち初稿）}\mbox{}\par

\noindent\textbf{1：不適切} A>B>C>D。式・定義・対応の独立判定と一致しない。\par

\noindent\textbf{2：不適切} A>B>D>C。式・定義・対応の独立判定と一致しない。\par

\noindent\textbf{3：不適切} B>C>D>A。式・定義・対応の独立判定と一致しない。\par

\noindent\textbf{4：適切・選択対象} C>A>B>D。判定と一致する。\par

\noindent\textbf{5：不適切} C>A>D>B。式・定義・対応の独立判定と一致しない。\par

\paragraph{再利用する考え方} 高位・低位発熱量、単位、燃料品種、密度換算。 指定年度・対象範囲・単位・否定を先に固定し、他年度の番号や現行値を転用しない。

\paragraph{共通理論} \hyperref[theory:env:heating]{第\ref*{theory:env:heating}節}

\par 参考：\cite{env58_heating}。

\par 難易度 3（校準初稿）、優先度未評価。

\subsection{令和元年度再試験 問5-3}
\label{q:r01-ret-basic-5-3}

\noindent\textbf{原典：} sources/R01\_retest/basic.pdf、PDF物理18頁。

\textbf{公式正答：} 3（single）。

\par\textbf{検討未確定：} 指定時版の標準発熱量の全四値は未取得。後年改訂表を指定時表の代替確定にはしない。

\paragraph{問題の要旨} LNG・原油・一般炭・絶乾材の質量発熱量。

\Needspace{4\baselineskip}\paragraph{確認した条件・記述（独自要約）}\mbox{}\par

\noindent\textbf{時点：} 令和元年度再試験に指定された資料・制度。\par

\noindent\textbf{照合：} 原典18画像で否定・五肢・組合せ表を確認。後年値で置換しない。\par

\paragraph{問われる核心} LNG・原油・一般炭・絶乾材の質量発熱量。

\paragraph{必要知識} 燃料指定、乾燥状態、熱量単位と密度換算。

\paragraph{考え方・正答を選ぶ理由（査読待ち初稿）} 指定表の単位質量ではLNG、原油、輸入一般炭、絶乾廃材の順で3。平成24年度は気化ガスと原料炭という異なる指定なので同じ数値を転用しない。ここでも体積単位は密度換算が必要。

\Needspace{4\baselineskip}\paragraph{選択肢の検討（査読待ち初稿）}\mbox{}\par

\noindent\textbf{1：不適切} A>B>C>D。式・定義・対応の独立判定と一致しない。\par

\noindent\textbf{2：不適切} B>A>D>C。式・定義・対応の独立判定と一致しない。\par

\noindent\textbf{3：適切・選択対象} C>A>B>D。判定と一致する。\par

\noindent\textbf{4：不適切} C>B>D>A。式・定義・対応の独立判定と一致しない。\par

\noindent\textbf{5：不適切} D>C>B>A。式・定義・対応の独立判定と一致しない。\par

\paragraph{再利用する考え方} 燃料指定、乾燥状態、熱量単位と密度換算。 集計対象・時点・法の主体・選択方向を固定し、同型でも元の五肢へ戻す。

\paragraph{共通理論} \hyperref[theory:env:heating]{第\ref*{theory:env:heating}節}

\par 参考：\cite{env58_heating}。

\par 難易度 3（校準初稿）、優先度未評価。

\section{発表・発明を古い順へ並べる}

\subsection{平成25年度 問5-6}
\label{q:h25-reg-basic-5-6}

\noindent\textbf{原典：} sources/H25/basic.pdf、PDF物理22頁。

\textbf{公式正答：} 1（single）。

\par\textbf{検討未確定：} 全五業績の初出一次史料は未取得。特にハレー彗星の「発見」は周期性の同定を指す出題表現として留保。

\paragraph{問題の要旨} 発表・発明を古い順へ並べる。

\Needspace{4\baselineskip}\paragraph{確認した条件・記述（独自要約）}\mbox{}\par

\noindent\textbf{時点：} 平成25年度に指定された資料・制度。\par

\noindent\textbf{照合：} 原典21画像で否定・五肢・組合せ表を確認。後年値で置換しない。\par

\paragraph{問われる核心} 発表・発明を古い順へ並べる。

\paragraph{必要知識} 科学史の世紀区分、業績の段階、順序制約。

\paragraph{考え方・正答を選ぶ理由（査読待ち初稿）} ハレーの周期彗星の予測・同定が18世紀初頭、Wattの蒸気機関改良が18世紀後半、Darwin–Wallaceの自然選択発表が1858、一般相対論が1915、Fukuiのフロンティア理論が1952。エ→ア→イ→オ→ウで1。発見と理論発表、特許と実用化のどの段階をいうかも確認する。

\Needspace{4\baselineskip}\paragraph{選択肢の検討（査読待ち初稿）}\mbox{}\par

\noindent\textbf{1：適切・選択対象} エア イオウ。判定と一致する。\par

\noindent\textbf{2：不適切} アイエウオ。式・定義・対応の独立判定と一致しない。\par

\noindent\textbf{3：不適切} エアイウオ。式・定義・対応の独立判定と一致しない。\par

\noindent\textbf{4：不適切} アエオイウ。式・定義・対応の独立判定と一致しない。\par

\noindent\textbf{5：不適切} アイエオウ。式・定義・対応の独立判定と一致しない。\par

\paragraph{再利用する考え方} 科学史の世紀区分、業績の段階、順序制約。 指定年度・対象範囲・単位・否定を先に固定し、他年度の番号や現行値を転用しない。

\paragraph{共通理論} \hyperref[theory:env:historyorder]{第\ref*{theory:env:historyorder}節}

\par 参考：\cite{env58_steam}。

\par 難易度 3（校準初稿）、優先度未評価。

\subsection{平成26年度 問5-6}
\label{q:h26-reg-basic-5-6}

\noindent\textbf{原典：} sources/H26/basic.pdf、PDF物理18頁。

\textbf{公式正答：} 4（single）。

\par\textbf{検討未確定：} Galileo・Wilsonの初出史料は追加照合。ACSの1935生産選定と1939商業化を分ける。

\paragraph{問題の要旨} 観測・動力・検出・高分子の技術史。

\Needspace{4\baselineskip}\paragraph{確認した条件・記述（独自要約）}\mbox{}\par

\noindent\textbf{時点：} 平成26年度に指定された資料・制度。\par

\noindent\textbf{照合：} 原典21画像で否定・五肢・組合せ表を確認。後年値で置換しない。\par

\paragraph{問われる核心} 観測・動力・検出・高分子の技術史。

\paragraph{必要知識} 改良と新機構、検出技術、合成と商業化。

\paragraph{考え方・正答を選ぶ理由（査読待ち初稿）} Galileoの望遠鏡観測が1609年頃、Newcomen機関1712、Wattの改良特許1769、Wilson霧箱は20世紀初頭、Carothersのナイロンは1930年代。ウ→ア→イ→オ→エの4。ナイロンの合成・生産選定と1939の商業化を同一年にせず、霧箱より後という比較に必要な範囲で扱う。

\Needspace{4\baselineskip}\paragraph{選択肢の検討（査読待ち初稿）}\mbox{}\par

\noindent\textbf{1：不適切} イアウオエ。式・定義・対応の独立判定と一致しない。\par

\noindent\textbf{2：不適切} イウアオエ。式・定義・対応の独立判定と一致しない。\par

\noindent\textbf{3：不適切} イアウエオ。式・定義・対応の独立判定と一致しない。\par

\noindent\textbf{4：適切・選択対象} ウアイオエ。判定と一致する。\par

\noindent\textbf{5：不適切} ウイアエオ。式・定義・対応の独立判定と一致しない。\par

\paragraph{再利用する考え方} 改良と新機構、検出技術、合成と商業化。 指定年度・対象範囲・単位・否定を先に固定し、他年度の番号や現行値を転用しない。

\paragraph{共通理論} \hyperref[theory:env:historyorder]{第\ref*{theory:env:historyorder}節}

\par 参考：\cite{env58_steam,env58_nylon}。

\par 難易度 3（校準初稿）、優先度未評価。

\subsection{平成28年度 問5-6}
\label{q:h28-reg-basic-5-6}

\noindent\textbf{原典：} sources/H28/basic.pdf、PDF物理24頁。

\textbf{公式正答：} 4（single）。

\par\textbf{検討未確定：} 雷の実験・三極管・トランジスタ各原報と特許の一次史料は追加照合。三人の貢献を同一作業に単純化しない。

\paragraph{問題の要旨} 電気・化学・電子の代表的業績の順序。

\Needspace{4\baselineskip}\paragraph{確認した条件・記述（独自要約）}\mbox{}\par

\noindent\textbf{時点：} 平成28年度に指定された資料・制度。\par

\noindent\textbf{照合：} 原典21画像で否定・五肢・組合せ表を確認。後年値で置換しない。\par

\paragraph{問われる核心} 電気・化学・電子の代表的業績の順序。

\paragraph{必要知識} 周期律、放射性元素、電子管と半導体、年代順。

\paragraph{考え方・正答を選ぶ理由（査読待ち初稿）} 雷の電気的性質のFranklinが18世紀、Mendeleevの周期律1869、Curie夫妻の新元素1898、De Forestの三極管1906頃、トランジスタ1947。ウ→ア→イ→オ→エの4。電子管から半導体素子への世紀をまたぐ発展と、発見・発明・商用普及の段階を分ける。

\Needspace{4\baselineskip}\paragraph{選択肢の検討（査読待ち初稿）}\mbox{}\par

\noindent\textbf{1：不適切} アイウエオ。式・定義・対応の独立判定と一致しない。\par

\noindent\textbf{2：不適切} アイウオエ。式・定義・対応の独立判定と一致しない。\par

\noindent\textbf{3：不適切} ウイアオエ。式・定義・対応の独立判定と一致しない。\par

\noindent\textbf{4：適切・選択対象} ウアイオエ。判定と一致する。\par

\noindent\textbf{5：不適切} ウイアエオ。式・定義・対応の独立判定と一致しない。\par

\paragraph{再利用する考え方} 周期律、放射性元素、電子管と半導体、年代順。 指定年度・対象範囲・単位・否定を先に固定し、他年度の番号や現行値を転用しない。

\paragraph{共通理論} \hyperref[theory:env:historyorder]{第\ref*{theory:env:historyorder}節}

\par 参考：\cite{env58_curie}。

\par 難易度 3（校準初稿）、優先度未評価。

\subsection{平成30年度 問5-5}
\label{q:h30-reg-basic-5-5}

\noindent\textbf{原典：} sources/H30/basic.pdf、PDF物理20頁。

\textbf{公式正答：} 3（single）。

\par\textbf{検討未確定：} 五項目の初出全原報は未取得。HaberとHahnの単独人物表記は共同の貢献の全体を示すとはしない。

\paragraph{問題の要旨} 動力・通信・化学・核の技術成果の順。

\Needspace{4\baselineskip}\paragraph{確認した条件・記述（独自要約）}\mbox{}\par

\noindent\textbf{時点：} 平成30年度に指定された資料・制度。\par

\noindent\textbf{照合：} 原典18画像で否定・五肢・組合せ表を確認。後年値で置換しない。\par

\paragraph{問われる核心} 動力・通信・化学・核の技術成果の順。

\paragraph{必要知識} 科学的確認、技術発明、工業化の段階。

\paragraph{考え方・正答を選ぶ理由（査読待ち初稿）} Wattの改良が18世紀後半、Bellの電話1876、Hertzの電磁波実験1880年代、Haberのアンモニア合成基礎1909頃、Hahnらの核分裂1938。オ→ウ→エ→ア→イの3。Haberの基礎とBoschによる大規模工業化、核分裂発見と発電開始を分ける。

\Needspace{4\baselineskip}\paragraph{選択肢の検討（査読待ち初稿）}\mbox{}\par

\noindent\textbf{1：不適切} ウエオイア。式・定義・対応の独立判定と一致しない。\par

\noindent\textbf{2：不適切} ウオアエイ。式・定義・対応の独立判定と一致しない。\par

\noindent\textbf{3：適切・選択対象} オウエアイ。判定と一致する。\par

\noindent\textbf{4：不適切} オエウイア。式・定義・対応の独立判定と一致しない。\par

\noindent\textbf{5：不適切} アオウエイ。式・定義・対応の独立判定と一致しない。\par

\paragraph{再利用する考え方} 科学的確認、技術発明、工業化の段階。 集計対象・時点・法の主体・選択方向を固定し、同型でも元の五肢へ戻す。

\paragraph{共通理論} \hyperref[theory:env:historyorder]{第\ref*{theory:env:historyorder}節}

\par 参考：\cite{env58_steam}。

\paragraph{関連問題} \hyperref[q:r03-reg-basic-5-5]{R03-REG-BASIC-5-5}

\par 難易度 3（校準初稿）、優先度未評価。

\subsection{令和元年度再試験 問5-5}
\label{q:r01-ret-basic-5-5}

\noindent\textbf{原典：} sources/R01\_retest/basic.pdf、PDF物理19頁。

\textbf{公式正答：} 4（single）。

\par\textbf{検討未確定：} 初出史料の全五項目照合は未完了。原典の彗星「発見」は周期性の同定として留保。

\paragraph{問題の要旨} 同じ年表でも元の選択順へ戻す。

\Needspace{4\baselineskip}\paragraph{確認した条件・記述（独自要約）}\mbox{}\par

\noindent\textbf{時点：} 令和元年度再試験に指定された資料・制度。\par

\noindent\textbf{照合：} 原典18画像で否定・五肢・組合せ表を確認。後年値で置換しない。\par

\paragraph{問われる核心} 同じ年表でも元の選択順へ戻す。

\paragraph{必要知識} 初出の段階、時代順、候補の並べ替え。

\paragraph{考え方・正答を選ぶ理由（査読待ち初稿）} Halleyの周期性同定が18世紀初頭、Watt18世紀後半、Darwin–Wallace1858、Einstein1915、Fukui1952。エ→ア→イ→オ→ウの4。平成25年度の同型問題の正答1を番号だけ転用しない。

\Needspace{4\baselineskip}\paragraph{選択肢の検討（査読待ち初稿）}\mbox{}\par

\noindent\textbf{1：不適切} アイエウオ。式・定義・対応の独立判定と一致しない。\par

\noindent\textbf{2：不適切} エアイウオ。式・定義・対応の独立判定と一致しない。\par

\noindent\textbf{3：不適切} アエオイウ。式・定義・対応の独立判定と一致しない。\par

\noindent\textbf{4：適切・選択対象} エアイオウ。判定と一致する。\par

\noindent\textbf{5：不適切} アイエオウ。式・定義・対応の独立判定と一致しない。\par

\paragraph{再利用する考え方} 初出の段階、時代順、候補の並べ替え。 集計対象・時点・法の主体・選択方向を固定し、同型でも元の五肢へ戻す。

\paragraph{共通理論} \hyperref[theory:env:historyorder]{第\ref*{theory:env:historyorder}節}

\par 参考：\cite{env58_steam}。

\par 難易度 3（校準初稿）、優先度未評価。

\subsection{令和2年度 問5-6}
\label{q:r02-reg-basic-5-6}

\noindent\textbf{原典：} sources/R02/basic.pdf、PDF物理22頁。

\textbf{公式正答：} 1（single）。

\par\textbf{検討未確定：} 全初出史料の取得は未了。電子管・トランジスタの共同貢献と特許の段階は追加照合。

\paragraph{問題の要旨} 種痘・周期律・新元素・電子素子の順。

\Needspace{4\baselineskip}\paragraph{確認した条件・記述（独自要約）}\mbox{}\par

\noindent\textbf{時点：} 令和2年度に指定された資料・制度。\par

\noindent\textbf{照合：} 原典18画像で否定・五肢・組合せ表を確認。後年値で置換しない。\par

\paragraph{問われる核心} 種痘・周期律・新元素・電子素子の順。

\paragraph{必要知識} 医療経験、元素周期、電子管から半導体。

\paragraph{考え方・正答を選ぶ理由（査読待ち初稿）} Jenner1796、Mendeleev1869、Curie夫妻1898、De Forest三極管1906頃、トランジスタ1947。イ→エ→ア→オ→ウの1。観測された事実、法則化、装置の発明を同じ授賞年へまとめない。

\Needspace{4\baselineskip}\paragraph{選択肢の検討（査読待ち初稿）}\mbox{}\par

\noindent\textbf{1：適切・選択対象} イエアオウ。判定と一致する。\par

\noindent\textbf{2：不適切} イエオウア。式・定義・対応の独立判定と一致しない。\par

\noindent\textbf{3：不適切} イオエアウ。式・定義・対応の独立判定と一致しない。\par

\noindent\textbf{4：不適切} エイオアウ。式・定義・対応の独立判定と一致しない。\par

\noindent\textbf{5：不適切} エオイアウ。式・定義・対応の独立判定と一致しない。\par

\paragraph{再利用する考え方} 医療経験、元素周期、電子管から半導体。 集計対象・時点・法の主体・選択方向を固定し、同型でも元の五肢へ戻す。

\paragraph{共通理論} \hyperref[theory:env:historyorder]{第\ref*{theory:env:historyorder}節}

\par 参考：\cite{env59_jenner,env58_curie}。

\par 難易度 3（校準初稿）、優先度未評価。

\subsection{令和3年度 問5-5}
\label{q:r03-reg-basic-5-5}

\noindent\textbf{原典：} sources/R03/basic.pdf、PDF物理25頁。

\textbf{公式正答：} 4（single）。

\par\textbf{検討未確定：} 業績の発表・特許・工業化の区分と各史料の年の追加照合が必要。

\paragraph{問題の要旨} 蒸気機関から核分裂までの順序。

\Needspace{4\baselineskip}\paragraph{確認した条件・記述（独自要約）}\mbox{}\par

\noindent\textbf{原典：} 全五肢と否定方向・組合せ表を画像確認。\par

\noindent\textbf{資料時点：} 令和3年度の指定資料・統計年を保持し、後年値に置換しない。\par

\paragraph{問われる核心} 蒸気機関から核分裂までの順序。

\paragraph{必要知識} 人物・対象業績・年代の三点照合。

\paragraph{考え方・正答を選ぶ理由（査読待ち初稿）} オWatt蒸気機関改良は18世紀後半、ウBell電話1876、エHertz電磁波確認1880年代後半、アHaberアンモニア合成20世紀初め、イHahn核分裂1938の順。オウエアイの4。発明・実験確認・工業化・受賞の年を一つの年として扱わない。

\Needspace{4\baselineskip}\paragraph{選択肢の検討（査読待ち初稿）}\mbox{}\par

\noindent\textbf{1：不適切} アオウエイ。独立に求めた組合せ オウエアイ と異なる。\par

\noindent\textbf{2：不適切} ウエオイア。独立に求めた組合せ オウエアイ と異なる。\par

\noindent\textbf{3：不適切} ウオアエイ。独立に求めた組合せ オウエアイ と異なる。\par

\noindent\textbf{4：適切・選択対象} オウエアイ。独立に求めた組合せ オウエアイ と一致する。\par

\noindent\textbf{5：不適切} オエウイア。独立に求めた組合せ オウエアイ と異なる。\par

\paragraph{再利用する考え方} 人物・対象業績・年代の三点照合。 否定方向と資料年、分母と単位を確認して五肢へ戻る。

\paragraph{共通理論} \hyperref[theory:env:historyorder]{第\ref*{theory:env:historyorder}節}

\paragraph{関連問題} \hyperref[q:h30-reg-basic-5-5]{H30-REG-BASIC-5-5}

\par 難易度 3（校準初稿）、優先度未評価。

\subsection{令和4年度 問5-6}
\label{q:r04-reg-basic-5-6}

\noindent\textbf{原典：} sources/R04/basic.pdf、PDF物理20頁。

\textbf{公式正答：} 3（single）。

\par\textbf{検討未確定：} 各業績の公刊・特許・開発年の史料は追加照合。ナイロンの開発と販売開始を同一視しない。

\paragraph{問題の要旨} 尿素合成・転炉・赤痢菌・KS鋼・ナイロン。

\Needspace{4\baselineskip}\paragraph{確認した条件・記述（独自要約）}\mbox{}\par

\noindent\textbf{原典：} 全五肢と否定方向・組合せ表を画像確認。\par

\noindent\textbf{資料時点：} 令和4年度の指定資料・統計年を保持し、後年値に置換しない。\par

\paragraph{問われる核心} 尿素合成・転炉・赤痢菌・KS鋼・ナイロン。

\paragraph{必要知識} 人物と対象物、発明・発見・商品化の年代。

\paragraph{考え方・正答を選ぶ理由（査読待ち初稿）} エWöhler尿素合成1828、アBessemer転炉1856、オ志賀赤痢菌1897、イ本多KS鋼1917、ウCarothersナイロン1935の順で3。材料の発明年と量産・販売年は分ける。化学合成から製鋼、細菌、磁性材料、高分子へ時代が進む。

\Needspace{4\baselineskip}\paragraph{選択肢の検討（査読待ち初稿）}\mbox{}\par

\noindent\textbf{1：不適切} アエイオウ。独立に求めた組合せ エアオイウ と異なる。\par

\noindent\textbf{2：不適切} アエオイウ。独立に求めた組合せ エアオイウ と異なる。\par

\noindent\textbf{3：適切・選択対象} エアオイウ。独立に求めた組合せ エアオイウ と一致する。\par

\noindent\textbf{4：不適切} エオアウイ。独立に求めた組合せ エアオイウ と異なる。\par

\noindent\textbf{5：不適切} オエアウイ。独立に求めた組合せ エアオイウ と異なる。\par

\paragraph{再利用する考え方} 人物と対象物、発明・発見・商品化の年代。 否定方向と資料年、分母と単位を確認して五肢へ戻る。

\paragraph{共通理論} \hyperref[theory:env:historyorder]{第\ref*{theory:env:historyorder}節}

\par 参考：\cite{env58_nylon}。

\par 難易度 3（校準初稿）、優先度未評価。

\subsection{令和6年度 問5-6}
\label{q:r06-reg-basic-5-6}

\noindent\textbf{原典：} sources/R06/basic.pdf、PDF物理20頁。

\textbf{公式正答：} 2（single）。

\par\textbf{検討未確定：} 各業績の指定史料と発表・出版年の照合を追加。

\paragraph{問題の要旨} BernoulliからWegenerまでの順序。

\Needspace{4\baselineskip}\paragraph{確認した条件・記述（独自要約）}\mbox{}\par

\noindent\textbf{原典：} 全五肢と否定方向・組合せ表を画像確認。\par

\noindent\textbf{資料時点：} 令和6年度の指定資料・統計年を保持し、後年値に置換しない。\par

\paragraph{問われる核心} BernoulliからWegenerまでの順序。

\paragraph{必要知識} 業績の年代区分と科学・技術の並行発展。

\paragraph{考え方・正答を選ぶ理由（査読待ち初稿）} アBernoulli1738の流体力学、イLavoisier18世紀後半の酸素燃焼理論、オRocket1829、ウDarwin種の起源1859、エWegener大陸移動説1912の順で2。実験、出版、車両製作、学説発表の基準を分ける。

\Needspace{4\baselineskip}\paragraph{選択肢の検討（査読待ち初稿）}\mbox{}\par

\noindent\textbf{1：不適切} アイウエオ。独立に求めた組合せ アイオウエ と異なる。\par

\noindent\textbf{2：適切・選択対象} アイオウエ。独立に求めた組合せ アイオウエ と一致する。\par

\noindent\textbf{3：不適切} イアウオエ。独立に求めた組合せ アイオウエ と異なる。\par

\noindent\textbf{4：不適切} エアイオウ。独立に求めた組合せ アイオウエ と異なる。\par

\noindent\textbf{5：不適切} エイアウオ。独立に求めた組合せ アイオウエ と異なる。\par

\paragraph{再利用する考え方} 業績の年代区分と科学・技術の並行発展。 否定方向と資料年、分母と単位を確認して五肢へ戻る。

\paragraph{共通理論} \hyperref[theory:env:historyorder]{第\ref*{theory:env:historyorder}節}

\par 難易度 3（校準初稿）、優先度未評価。

\section{人物と業績：放射能とラジウムを区別}

\subsection{平成23年度 問5-5}
\label{q:h23-reg-basic-5-5}

\noindent\textbf{原典：} sources/H23/basic.pdf、PDF物理17頁。

\textbf{公式正答：} 4（single）。

\par\textbf{検討未確定：} 原典は「アントワーヌ・ベクレル」と記載し人名が曖昧。放射能発見者Henri Becquerelとの対応を区別。Eastmanの「フィルム乾板」も乾板とロールフィルムをまとめた表現であり、原典を厳密な発明年表に一般化しない。

\paragraph{問題の要旨} 人物と業績：放射能とラジウムを区別。

\Needspace{4\baselineskip}\paragraph{確認した条件・記述（独自要約）}\mbox{}\par

\noindent\textbf{時点：} 平成23年度に指定された資料・制度。\par

\noindent\textbf{照合：} 原典21画像で否定・五肢・組合せ表を確認。後年値で置換しない。\par

\paragraph{問われる核心} 人物と業績：放射能とラジウムを区別。

\paragraph{必要知識} 放射能現象と元素発見、発明と改良・普及の区別。

\paragraph{考え方・正答を選ぶ理由（査読待ち初稿）} ラジウムとポロニウムの発見は1898年のキュリー夫妻の仕事。Henri Becquerelは1896年ウランの自然放射を見いだしたので4が誤り。Galileoは望遠鏡を改良して天体観測へ用いた。Huygensの振子時計、Arkwrightの水力紡績、Eastmanの写真材料・乾板の工業化という他の対応は、出題の代表的業績の範囲で扱う。

\Needspace{4\baselineskip}\paragraph{選択肢の検討（査読待ち初稿）}\mbox{}\par

\noindent\textbf{1：条件内で適切} Galileo/天体望遠鏡で観測。望遠鏡の発明者そのものとする記述ではなく観測への利用。\par

\noindent\textbf{2：条件内で適切} Huygens/振子時計。振子を時間標準にする時計。\par

\noindent\textbf{3：条件内で適切} Arkwright/水力紡績機。産業革命の水力紡績機。\par

\noindent\textbf{4：不適切・選択対象} Becquerel/ラジウム発見。放射能の発見とラジウム発見を取り違える。\par

\noindent\textbf{5：条件内で適切} Eastman/写真用フィルム乾板。乾板と柔軟な写真フィルムの開発・商業化という原典の代表表現。\par

\paragraph{再利用する考え方} 放射能現象と元素発見、発明と改良・普及の区別。 指定年度・対象範囲・単位・否定を先に固定し、他年度の番号や現行値を転用しない。

\paragraph{共通理論} \hyperref[theory:env:historypairs]{第\ref*{theory:env:historypairs}節}

\par 参考：\cite{env58_curie}。

\par 難易度 3（校準初稿）、優先度未評価。

\subsection{平成27年度 問5-6}
\label{q:h27-reg-basic-5-6}

\noindent\textbf{原典：} sources/H27/basic.pdf、PDF物理18頁。

\textbf{公式正答：} 4（single）。

\par\textbf{検討未確定：} 原典は「アントワーヌ・ベクレル」と記載し人名が曖昧。放射能発見者Henri Becquerelとの対応を区別。Eastmanの「フィルム乾板」も乾板とロールフィルムをまとめた表現であり、原典を厳密な発明年表に一般化しない。

\paragraph{問題の要旨} 人物と業績：放射能とラジウムを区別。

\Needspace{4\baselineskip}\paragraph{確認した条件・記述（独自要約）}\mbox{}\par

\noindent\textbf{時点：} 平成27年度に指定された資料・制度。\par

\noindent\textbf{照合：} 原典21画像で否定・五肢・組合せ表を確認。後年値で置換しない。\par

\paragraph{問われる核心} 人物と業績：放射能とラジウムを区別。

\paragraph{必要知識} 放射能現象と元素発見、発明と改良・普及の区別。

\paragraph{考え方・正答を選ぶ理由（査読待ち初稿）} ラジウムとポロニウムの発見は1898年のキュリー夫妻の仕事。Henri Becquerelは1896年ウランの自然放射を見いだしたので4が誤り。Galileoは望遠鏡を改良して天体観測へ用いた。Huygensの振子時計、Arkwrightの水力紡績、Eastmanの写真材料・乾板の工業化という他の対応は、出題の代表的業績の範囲で扱う。

\Needspace{4\baselineskip}\paragraph{選択肢の検討（査読待ち初稿）}\mbox{}\par

\noindent\textbf{1：条件内で適切} Galileo/天体望遠鏡で観測。望遠鏡の発明者そのものとする記述ではなく観測への利用。\par

\noindent\textbf{2：条件内で適切} Huygens/振子時計。振子を時間標準にする時計。\par

\noindent\textbf{3：条件内で適切} Arkwright/水力紡績機。産業革命の水力紡績機。\par

\noindent\textbf{4：不適切・選択対象} Becquerel/ラジウム発見。放射能の発見とラジウム発見を取り違える。\par

\noindent\textbf{5：条件内で適切} Eastman/写真用フィルム乾板。乾板と柔軟な写真フィルムの開発・商業化という原典の代表表現。\par

\paragraph{再利用する考え方} 放射能現象と元素発見、発明と改良・普及の区別。 指定年度・対象範囲・単位・否定を先に固定し、他年度の番号や現行値を転用しない。

\paragraph{共通理論} \hyperref[theory:env:historypairs]{第\ref*{theory:env:historypairs}節}

\par 参考：\cite{env58_curie}。

\par 難易度 3（校準初稿）、優先度未評価。

\section{産業革命の改良・工場化と担い手}

\subsection{平成24年度 問5-5}
\label{q:h24-reg-basic-5-5}

\noindent\textbf{原典：} sources/H24/basic.pdf、PDF物理17頁。

\textbf{公式正答：} 4（single）。

\par\textbf{検討未確定：} からくりからの転用という肢3の個別史料は追加照合。一般論を特定発明の確定年表に拡大しない。

\paragraph{問題の要旨} 産業革命の改良・工場化と担い手。

\Needspace{4\baselineskip}\paragraph{確認した条件・記述（独自要約）}\mbox{}\par

\noindent\textbf{時点：} 平成24年度に指定された資料・制度。\par

\noindent\textbf{照合：} 原典21画像で否定・五肢・組合せ表を確認。後年値で置換しない。\par

\paragraph{問われる核心} 産業革命の改良・工場化と担い手。

\paragraph{必要知識} 発明、改良、普及、研究組織と生産制度。

\paragraph{考え方・正答を選ぶ理由（査読待ち初稿）} 産業革命期の代表的機械は職人・発明家・事業者の改良と事業化で広まり、当時の欧州大学研究者が開発を主導した例が多いという4は不適切。新規性だけでなく既存機構の組合せ、工場制度、生産と社会への影響も見る。Wattと大学との接点のような個別事例は、全体を大学主導とする証明ではない。

\Needspace{4\baselineskip}\paragraph{選択肢の検討（査読待ち初稿）}\mbox{}\par

\noindent\textbf{1：条件内で適切} 既存技術の改良の積重ね。原典の条件内では適切な記述であり、不適切なものを選ぶ問の選択対象ではない。\par

\noindent\textbf{2：条件内で適切} 機械化で大工場へ。原典の条件内では適切な記述であり、不適切なものを選ぶ問の選択対象ではない。\par

\noindent\textbf{3：条件内で適切} 娯楽自動機構から転用もある。原典の条件内では適切な記述であり、不適切なものを選ぶ問の選択対象ではない。\par

\noindent\textbf{4：不適切・選択対象} 大学研究者主導の産学協同例が多い。判定と一致する。\par

\noindent\textbf{5：条件内で適切} 労働と児童労働が社会問題。原典の条件内では適切な記述であり、不適切なものを選ぶ問の選択対象ではない。\par

\paragraph{再利用する考え方} 発明、改良、普及、研究組織と生産制度。 指定年度・対象範囲・単位・否定を先に固定し、他年度の番号や現行値を転用しない。

\paragraph{共通理論} \hyperref[theory:env:industrial]{第\ref*{theory:env:industrial}節}

\par 参考：\cite{env58_steam}。

\par 難易度 3（校準初稿）、優先度未評価。

\subsection{平成29年度 問5-5}
\label{q:h29-reg-basic-5-5}

\noindent\textbf{原典：} sources/H29/basic.pdf、PDF物理19頁。

\textbf{公式正答：} 2（single）。

\par\textbf{検討未確定：} 個別の娯楽機構転用史料は未取得。大学との個別接点も無かったという全称にはしない。

\paragraph{問題の要旨} 産業革命の担い手を大学主導へ一般化しない。

\Needspace{4\baselineskip}\paragraph{確認した条件・記述（独自要約）}\mbox{}\par

\noindent\textbf{時点：} 平成29年度に指定された資料・制度。\par

\noindent\textbf{照合：} 原典18画像で否定・五肢・組合せ表を確認。後年値で置換しない。\par

\paragraph{問われる核心} 産業革命の担い手を大学主導へ一般化しない。

\paragraph{必要知識} 産業革命、改良と普及、制度と労働。

\paragraph{考え方・正答を選ぶ理由（査読待ち初稿）} 2が不適切。職人・発明家・企業家の改良と事業化を大学研究者主導へ一括置換できない。機械技術の利用、生産組織、労働制度という別の側面を合わせて整理する。

\Needspace{4\baselineskip}\paragraph{選択肢の検討（査読待ち初稿）}\mbox{}\par

\noindent\textbf{1：条件内で適切} 既存改良の積重ね。条件内で適切な記述であり、不適切な記述を選ぶ対象ではない。\par

\noindent\textbf{2：不適切・選択対象} 欧州大学研究者主導が多い。判定と一致する。\par

\noindent\textbf{3：条件内で適切} 大工場への機械化。条件内で適切な記述であり、不適切な記述を選ぶ対象ではない。\par

\noindent\textbf{4：条件内で適切} 娯楽機構から転用例。条件内で適切な記述であり、不適切な記述を選ぶ対象ではない。\par

\noindent\textbf{5：条件内で適切} 児童労働等の社会問題。条件内で適切な記述であり、不適切な記述を選ぶ対象ではない。\par

\paragraph{再利用する考え方} 産業革命、改良と普及、制度と労働。 集計対象・時点・法の主体・選択方向を固定し、同型でも元の五肢へ戻す。

\paragraph{共通理論} \hyperref[theory:env:industrial]{第\ref*{theory:env:industrial}節}

\par 参考：\cite{env58_steam}。

\par 難易度 3（校準初稿）、優先度未評価。

\subsection{令和6年度 問5-5}
\label{q:r06-reg-basic-5-5}

\noindent\textbf{原典：} sources/R06/basic.pdf、PDF物理20頁。

\textbf{公式正答：} 3（single）。

\par\textbf{検討未確定：} 娯楽機構の具体的転用例と地域ごとの担い手について原典史料追加照合。

\paragraph{問題の要旨} 産業革命は大学主導の開発がほとんどではない。

\Needspace{4\baselineskip}\paragraph{確認した条件・記述（独自要約）}\mbox{}\par

\noindent\textbf{原典：} 全五肢と否定方向・組合せ表を画像確認。\par

\noindent\textbf{資料時点：} 令和6年度の指定資料・統計年を保持し、後年値に置換しない。\par

\paragraph{問われる核心} 産業革命は大学主導の開発がほとんどではない。

\paragraph{必要知識} 技術改良の担い手、工場制、生産性と労働問題。

\paragraph{考え方・正答を選ぶ理由（査読待ち初稿）} 3が不適切。18〜19世紀の産業技術は職人の経験、発明家の工夫、資本・市場の需要等から発達し、ほぼ全てを現代的産学共同研究と定義できない。科学・技術・経済・労働の関係を区別する。

\Needspace{4\baselineskip}\paragraph{選択肢の検討（査読待ち初稿）}\mbox{}\par

\noindent\textbf{1：条件内で適切} 多くは既存技術の改良の積み重ね。技術革新は段階的な改良を含む。\par

\noindent\textbf{2：条件内で適切} 娯楽用機構のアイデアが転用。機構や技能の用途間移転がある。\par

\noindent\textbf{3：不適切・選択対象} ほとんどが大学研究者主導の産学協同。職人・発明家・事業家等の活動を落とす過度な一般化で不適切。\par

\noindent\textbf{4：条件内で適切} 機械化で大工場へ。動力と機械の集積が工場制へつながる。\par

\noindent\textbf{5：条件内で適切} 効率化と長時間・児童労働等の問題。生産性向上と労働者の環境改善は自動的には同時に実現しない。\par

\paragraph{再利用する考え方} 技術改良の担い手、工場制、生産性と労働問題。 否定方向と資料年、分母と単位を確認して五肢へ戻る。

\paragraph{共通理論} \hyperref[theory:env:industrial]{第\ref*{theory:env:industrial}節}

\par 難易度 3（校準初稿）、優先度未評価。

\section{知的財産の範囲は特許発明より広い}

\subsection{平成27年度 問5-5}
\label{q:h27-reg-basic-5-5}

\noindent\textbf{原典：} sources/H27/basic.pdf、PDF物理17頁。

\textbf{公式正答：} 3（single）。

\paragraph{問題の要旨} 知的財産の範囲は特許発明より広い。

\Needspace{4\baselineskip}\paragraph{確認した条件・記述（独自要約）}\mbox{}\par

\noindent\textbf{時点：} 平成27年度に指定された資料・制度。\par

\noindent\textbf{照合：} 原典21画像で否定・五肢・組合せ表を確認。後年値で置換しない。\par

\paragraph{問われる核心} 知的財産の範囲は特許発明より広い。

\paragraph{必要知識} 発明、知的財産、産業財産権、基本法の責務。

\paragraph{考え方・正答を選ぶ理由（査読待ち初稿）} 特許法の発明は自然法則を利用した高度な技術的思想。知的財産基本法の知財は著作物・商標・商号・営業秘密なども含むので、自然法則利用へ限定する3が不適切。法律名とその定義の範囲を別々に確認する。

\Needspace{4\baselineskip}\paragraph{選択肢の検討（査読待ち初稿）}\mbox{}\par

\noindent\textbf{1：条件内で適切} 特許は発明保護利用と産業発達。原典の条件内では適切な記述であり、不適切なものを選ぶ問の選択対象ではない。\par

\noindent\textbf{2：条件内で適切} 発明は自然法則利用の高度技術思想。原典の条件内では適切な記述であり、不適切なものを選ぶ問の選択対象ではない。\par

\noindent\textbf{3：不適切・選択対象} 知財は自然法則利用のものだけ。判定と一致する。\par

\noindent\textbf{4：条件内で適切} 基本法は創造保護活用の基本事項。原典の条件内では適切な記述であり、不適切なものを選ぶ問の選択対象ではない。\par

\noindent\textbf{5：条件内で適切} 国に施策策定実施責務。原典の条件内では適切な記述であり、不適切なものを選ぶ問の選択対象ではない。\par

\paragraph{再利用する考え方} 発明、知的財産、産業財産権、基本法の責務。 指定年度・対象範囲・単位・否定を先に固定し、他年度の番号や現行値を転用しない。

\paragraph{共通理論} \hyperref[theory:env:ip]{第\ref*{theory:env:ip}節}

\par 参考：\cite{ip_patent_statute,ip_broad_scope}。

\par 難易度 2（校準初稿）、優先度未評価。

\subsection{令和元年度 問5-6}
\label{q:r01-reg-basic-5-6}

\noindent\textbf{原典：} sources/R01/basic.pdf、PDF物理22頁。

\textbf{公式正答：} 5（single）。

\paragraph{問題の要旨} 特許と基本法の定義、国の責務。

\Needspace{4\baselineskip}\paragraph{確認した条件・記述（独自要約）}\mbox{}\par

\noindent\textbf{時点：} 令和元年度に指定された資料・制度。\par

\noindent\textbf{照合：} 原典18画像で否定・五肢・組合せ表を確認。後年値で置換しない。\par

\paragraph{問われる核心} 特許と基本法の定義、国の責務。

\paragraph{必要知識} 特許法1・2条、基本法2・5条、対象の広さ。

\paragraph{考え方・正答を選ぶ理由（査読待ち初稿）} 知的財産基本法は国の施策策定・実施責務を定めるので、無しとする5が誤り。発明の定義と知財全体の定義を区別し、商標や商号は特許発明の要件で排除しない。

\Needspace{4\baselineskip}\paragraph{選択肢の検討（査読待ち初稿）}\mbox{}\par

\noindent\textbf{1：条件内で適切} 発明は自然法則利用高度思想。条件内で適切な記述であり、不適切な記述を選ぶ対象ではない。\par

\noindent\textbf{2：条件内で適切} 特許保護利用と産業。条件内で適切な記述であり、不適切な記述を選ぶ対象ではない。\par

\noindent\textbf{3：条件内で適切} 知財には商標商号も。条件内で適切な記述であり、不適切な記述を選ぶ対象ではない。\par

\noindent\textbf{4：条件内で適切} 基本法は創造保護活用。条件内で適切な記述であり、不適切な記述を選ぶ対象ではない。\par

\noindent\textbf{5：不適切・選択対象} 国は施策責務なし。判定と一致する。\par

\paragraph{再利用する考え方} 特許法1・2条、基本法2・5条、対象の広さ。 集計対象・時点・法の主体・選択方向を固定し、同型でも元の五肢へ戻す。

\paragraph{共通理論} \hyperref[theory:env:ip]{第\ref*{theory:env:ip}節}

\par 参考：\cite{ip_patent_statute,ip_broad_scope}。

\par 難易度 2（校準初稿）、優先度未評価。

\section{気候緩和と適応、LCAの境界}

\subsection{平成25年度 問5-3}
\label{q:h25-reg-basic-5-3}

\noindent\textbf{原典：} sources/H25/basic.pdf、PDF物理20頁。

\textbf{公式正答：} 4（single）。

\par\textbf{検討未確定：} CDMとCFPの出題当時規則全文は未取得。定義の独立判断と条文・規格の確定は分ける。

\paragraph{問題の要旨} 気候緩和と適応、LCAの境界。

\Needspace{4\baselineskip}\paragraph{確認した条件・記述（独自要約）}\mbox{}\par

\noindent\textbf{時点：} 平成25年度に指定された資料・制度。\par

\noindent\textbf{照合：} 原典21画像で否定・五肢・組合せ表を確認。後年値で置換しない。\par

\paragraph{問われる核心} 気候緩和と適応、LCAの境界。

\paragraph{必要知識} 緩和・適応、機能単位とシステム境界、環境基準。

\paragraph{考え方・正答を選ぶ理由（査読待ち初稿）} 排出削減や吸収を通じて原因を減らすのが緩和。影響に対して被害を減らす対策が適応なので4が取り違え。CFPは温室効果ガスのCO2換算に着目し、LCAはより広い環境影響の評価。環境基準は環境の望ましい状態で、個別工場の排出許容値と同じではない。

\Needspace{4\baselineskip}\paragraph{選択肢の検討（査読待ち初稿）}\mbox{}\par

\noindent\textbf{1：条件内で適切} 環境基準は維持が望ましい。原典の条件内では適切な記述であり、不適切なものを選ぶ問の選択対象ではない。\par

\noindent\textbf{2：条件内で適切} CDMは京都の認証削減。原典の条件内では適切な記述であり、不適切なものを選ぶ問の選択対象ではない。\par

\noindent\textbf{3：条件内で適切} CFPは全過程CO2換算。原典の条件内では適切な記述であり、不適切なものを選ぶ問の選択対象ではない。\par

\noindent\textbf{4：不適切・選択対象} 適応は排出削減で原因抑制。判定と一致する。\par

\noindent\textbf{5：条件内で適切} LCAは資源から廃棄まで評価。原典の条件内では適切な記述であり、不適切なものを選ぶ問の選択対象ではない。\par

\paragraph{再利用する考え方} 緩和・適応、機能単位とシステム境界、環境基準。 指定年度・対象範囲・単位・否定を先に固定し、他年度の番号や現行値を転用しない。

\paragraph{共通理論} \hyperref[theory:env:lca]{第\ref*{theory:env:lca}節}

\par 参考：\cite{env58_pmstandard}。

\par 難易度 3（校準初稿）、優先度未評価。

\subsection{令和元年度 問5-2}
\label{q:r01-reg-basic-5-2}

\noindent\textbf{原典：} sources/R01/basic.pdf、PDF物理20頁。

\textbf{公式正答：} 2（single）。

\par\textbf{検討未確定：} JCMの2019運用ルール全文とCFP規格は追加照合。

\paragraph{問題の要旨} 排出原因を減らす緩和と影響への適応。

\Needspace{4\baselineskip}\paragraph{確認した条件・記述（独自要約）}\mbox{}\par

\noindent\textbf{時点：} 令和元年度に指定された資料・制度。\par

\noindent\textbf{照合：} 原典18画像で否定・五肢・組合せ表を確認。後年値で置換しない。\par

\paragraph{問われる核心} 排出原因を減らす緩和と影響への適応。

\paragraph{必要知識} 緩和適応、LCA、CFP、JCM、環境基準。

\paragraph{考え方・正答を選ぶ理由（査読待ち初稿）} 原因の削減を適応と呼ぶ2が誤り。排出削減・吸収は緩和で、影響の被害回避が適応。CFPとLCA、環境基準の定義は対象量と境界を固定して判定する。JCMは削減・吸収への貢献を定量化する枠組みで、効果の無条件な認定ではない。

\Needspace{4\baselineskip}\paragraph{選択肢の検討（査読待ち初稿）}\mbox{}\par

\noindent\textbf{1：条件内で適切} JCMは貢献定量化と目標活用。条件内で適切な記述であり、不適切な記述を選ぶ対象ではない。\par

\noindent\textbf{2：不適切・選択対象} 適応は原因のガス排出削減。判定と一致する。\par

\noindent\textbf{3：条件内で適切} CFP全過程CO2換算。条件内で適切な記述であり、不適切な記述を選ぶ対象ではない。\par

\noindent\textbf{4：条件内で適切} LCA全段階影響評価。条件内で適切な記述であり、不適切な記述を選ぶ対象ではない。\par

\noindent\textbf{5：条件内で適切} 環境基準は望ましい状態。条件内で適切な記述であり、不適切な記述を選ぶ対象ではない。\par

\paragraph{再利用する考え方} 緩和適応、LCA、CFP、JCM、環境基準。 集計対象・時点・法の主体・選択方向を固定し、同型でも元の五肢へ戻す。

\paragraph{共通理論} \hyperref[theory:env:lca]{第\ref*{theory:env:lca}節}

\par 参考：\cite{env59_ndc,env58_pmstandard}。

\par 難易度 3（校準初稿）、優先度未評価。

\section{気液の同質量比較からLNGの体積比}

\subsection{平成23年度 問5-3}
\label{q:h23-reg-basic-5-3}

\noindent\textbf{原典：} sources/H23/basic.pdf、PDF物理16頁。

\textbf{公式正答：} 2（single）。

\paragraph{問題の要旨} 気液の同質量比較からLNGの体積比。

\Needspace{4\baselineskip}\paragraph{確認した条件・記述（独自要約）}\mbox{}\par

\noindent\textbf{気体：} 0℃1気圧、純CH4の理想気体。\par

\noindent\textbf{液体：} 密度425kg/m\textsuperscript{3}一定とする原典モデル。\par

\paragraph{問われる核心} 気液の同質量比較からLNGの体積比。

\paragraph{必要知識} 理想気体モル体積、分子量、密度と体積、単位。

\paragraph{考え方・正答を選ぶ理由（査読待ち初稿）} 同じ物質量を比べる。0℃1気圧で1molの理想気体は約22.4L、メタン質量16g。液体体積は0.016kg/(425kg/m\textsuperscript{3})=\allowbreak{}3.76×10\textasciicircum{}-5m\textsuperscript{3}、気体は.0224m\textsuperscript{3}。比は.00168≈1/595なので最も近い1/600。密度比ρgas/ρliquidでも同じ。冷却の温度条件を気体標準状態と混同しない。

\Needspace{4\baselineskip}\paragraph{選択肢の検討（査読待ち初稿）}\mbox{}\par

\noindent\textbf{1：不適切} 1/400。式・定義・対応の独立判定と一致しない。\par

\noindent\textbf{2：適切・選択対象} 1/600。判定と一致する。\par

\noindent\textbf{3：不適切} 1/800。式・定義・対応の独立判定と一致しない。\par

\noindent\textbf{4：不適切} 1/1000。式・定義・対応の独立判定と一致しない。\par

\noindent\textbf{5：不適切} 1/1200。式・定義・対応の独立判定と一致しない。\par

\paragraph{再利用する考え方} 理想気体モル体積、分子量、密度と体積、単位。 指定年度・対象範囲・単位・否定を先に固定し、他年度の番号や現行値を転用しない。

\paragraph{共通理論} \hyperref[theory:env:lng]{第\ref*{theory:env:lng}節}

\paragraph{関連問題} \hyperref[q:r05-reg-basic-5-4]{R05-REG-BASIC-5-4}

\par 難易度 2（校準初稿）、優先度未評価。

\subsection{平成29年度 問5-3}
\label{q:h29-reg-basic-5-3}

\noindent\textbf{原典：} sources/H29/basic.pdf、PDF物理19頁。

\textbf{公式正答：} 4（single）。

\paragraph{問題の要旨} 標準気体と液体の密度比からLNG体積比。

\Needspace{4\baselineskip}\paragraph{確認した条件・記述（独自要約）}\mbox{}\par

\noindent\textbf{標準状態：} 0℃1気圧、純CH4理想気体。\par

\noindent\textbf{液体：} 425kg/m\textsuperscript{3}一定。\par

\paragraph{問われる核心} 標準気体と液体の密度比からLNG体積比。

\paragraph{必要知識} 理想気体、密度、同質量の体積比較。

\paragraph{考え方・正答を選ぶ理由（査読待ち初稿）} 1molのCH4は約16gで、0℃1気圧の体積22.4L。液体では.016/425=\allowbreak{}3.76×10\textasciicircum{}-5m\textsuperscript{3}、気体.0224m\textsuperscript{3}で比≈1/595。1/600の4が近い。平成23年度の候補順をそのまま使わない。

\Needspace{4\baselineskip}\paragraph{選択肢の検討（査読待ち初稿）}\mbox{}\par

\noindent\textbf{1：不適切} 1/1200。式・定義・対応の独立判定と一致しない。\par

\noindent\textbf{2：不適切} 1/1000。式・定義・対応の独立判定と一致しない。\par

\noindent\textbf{3：不適切} 1/800。式・定義・対応の独立判定と一致しない。\par

\noindent\textbf{4：適切・選択対象} 1/600。判定と一致する。\par

\noindent\textbf{5：不適切} 1/400。式・定義・対応の独立判定と一致しない。\par

\paragraph{再利用する考え方} 理想気体、密度、同質量の体積比較。 集計対象・時点・法の主体・選択方向を固定し、同型でも元の五肢へ戻す。

\paragraph{共通理論} \hyperref[theory:env:lng]{第\ref*{theory:env:lng}節}

\paragraph{関連問題} \hyperref[q:r05-reg-basic-5-4]{R05-REG-BASIC-5-4}

\par 難易度 2（校準初稿）、優先度未評価。

\subsection{令和5年度 問5-4}
\label{q:r05-reg-basic-5-4}

\noindent\textbf{原典：} sources/R05/basic.pdf、PDF物理17頁。

\textbf{公式正答：} 2（single）。

\paragraph{問題の要旨} 0℃1気圧のメタンと液体密度から体積比。

\Needspace{4\baselineskip}\paragraph{確認した条件・記述（独自要約）}\mbox{}\par

\noindent\textbf{気体：} 純CH4理想気体、0℃・1気圧。\par

\noindent\textbf{液体：} 密度425kg/m\textsuperscript{3}一定。\par

\paragraph{問われる核心} 0℃1気圧のメタンと液体密度から体積比。

\paragraph{必要知識} 同質量の体積比、理想気体の標準体積、密度。

\paragraph{考え方・正答を選ぶ理由（査読待ち初稿）} メタン1molの質量は約0.016kg。指定液体密度425kg/m\textsuperscript{3}より液体体積0.016/425=\allowbreak{}3.76×10\textasciicircum{}-5m\textsuperscript{3}、0℃1気圧の理想気体体積は約0.0224m\textsuperscript{3}。比は0.00168≈1/595で1/600の2。液体側に気体の状態方程式を適用しない。

\Needspace{4\baselineskip}\paragraph{選択肢の検討（査読待ち初稿）}\mbox{}\par

\noindent\textbf{1：不適切} 1/400。独立に求めた組合せ 1/600 と異なる。\par

\noindent\textbf{2：適切・選択対象} 1/600。独立に求めた組合せ 1/600 と一致する。\par

\noindent\textbf{3：不適切} 1/800。独立に求めた組合せ 1/600 と異なる。\par

\noindent\textbf{4：不適切} 1/1000。独立に求めた組合せ 1/600 と異なる。\par

\noindent\textbf{5：不適切} 1/1200。独立に求めた組合せ 1/600 と異なる。\par

\paragraph{再利用する考え方} 同質量の体積比、理想気体の標準体積、密度。 否定方向と資料年、分母と単位を確認して五肢へ戻る。

\paragraph{共通理論} \hyperref[theory:env:lng]{第\ref*{theory:env:lng}節}

\paragraph{関連問題} \hyperref[q:h23-reg-basic-5-3]{H23-REG-BASIC-5-3}, \hyperref[q:h29-reg-basic-5-3]{H29-REG-BASIC-5-3}

\par 難易度 2（校準初稿）、優先度未評価。

\section{2012年石油統計の対象範囲を固定}

\subsection{平成25年度 問5-1}
\label{q:h25-reg-basic-5-1}

\noindent\textbf{原典：} sources/H25/basic.pdf、PDF物理19頁。

\textbf{公式正答：} 5（single）。

\par\textbf{検討未確定：} 2012暦年の財務省四項目原表を全て取得していない。取得白書は2013年度83.6\%とサウジ首位を支持する補助で、指定年の確定証跡とは区別する。

\paragraph{問題の要旨} 2012年石油統計の対象範囲を固定。

\Needspace{4\baselineskip}\paragraph{確認した条件・記述（独自要約）}\mbox{}\par

\noindent\textbf{時点：} 平成25年度に指定された資料・制度。\par

\noindent\textbf{照合：} 原典21画像で否定・五肢・組合せ表を確認。後年値で置換しない。\par

\paragraph{問われる核心} 2012年石油統計の対象範囲を固定。

\paragraph{必要知識} 暦年と年度、数量と金額、統計品目、輸送経路。

\paragraph{考え方・正答を選ぶ理由（査読待ち初稿）} 2012年の指定統計では中東依存度約83\%、主要経路ホルムズ海峡、最大輸入先サウジアラビア、原油と石油製品の総輸入額比約21\%で5。数量の依存度と金額比は別であり、石炭・LNGを含む鉱物性燃料全体の比と原油・製品だけの比も違う。

\Needspace{4\baselineskip}\paragraph{選択肢の検討（査読待ち初稿）}\mbox{}\par

\noindent\textbf{1：不適切} 93/マラッカ/クウェート/21。式・定義・対応の独立判定と一致しない。\par

\noindent\textbf{2：不適切} 93/ホルムズ/サウジ/42。式・定義・対応の独立判定と一致しない。\par

\noindent\textbf{3：不適切} 83/マラッカ/サウジ/42。式・定義・対応の独立判定と一致しない。\par

\noindent\textbf{4：不適切} 83/ホルムズ/クウェート/42。式・定義・対応の独立判定と一致しない。\par

\noindent\textbf{5：適切・選択対象} 83/ホルムズ/サウジ/21。判定と一致する。\par

\paragraph{再利用する考え方} 暦年と年度、数量と金額、統計品目、輸送経路。 指定年度・対象範囲・単位・否定を先に固定し、他年度の番号や現行値を転用しない。

\paragraph{共通理論} \hyperref[theory:env:oil2012]{第\ref*{theory:env:oil2012}節}

\par 参考：\cite{env58_oil}。

\par 難易度 3（校準初稿）、優先度未評価。

\section{リサイクル制度の対象と責任主体}

\subsection{平成24年度 問5-1}
\label{q:h24-reg-basic-5-1}

\noindent\textbf{原典：} sources/H24/basic.pdf、PDF物理15頁。

\textbf{公式正答：} 3（single）。

\par\textbf{検討未確定：} 出題時点の建設規模・バーゼル条約対象と例外は追加の法令履歴照合。現行FAQだけで全過去条文の確定とはしない。

\paragraph{問題の要旨} リサイクル制度の対象と責任主体。

\Needspace{4\baselineskip}\paragraph{確認した条件・記述（独自要約）}\mbox{}\par

\noindent\textbf{時点：} 平成24年度に指定された資料・制度。\par

\noindent\textbf{照合：} 原典21画像で否定・五肢・組合せ表を確認。後年値で置換しない。\par

\paragraph{問われる核心} リサイクル制度の対象と責任主体。

\paragraph{必要知識} 法の目的、対象条件、排出者・自治体・小売・製造者の役割。

\paragraph{考え方・正答を選ぶ理由（査読待ち初稿）} 3が家電リサイクル法の小売業者の役割に対応。基本法はまず発生抑制。容器包装は三品目のみではなく、分別収集と再商品化の役割・費用も異なる。建設は規模等の対象条件があり全工事・発注者への実施義務という全称が違う。バーゼルは再利用目的の有害廃棄物でも規制を受け、条約成立の背景の輸送方向も逆。

\Needspace{4\baselineskip}\paragraph{選択肢の検討（査読待ち初稿）}\mbox{}\par

\noindent\textbf{1：不適切} 基本法は再生利用最優先。式・定義・対応の独立判定と一致しない。\par

\noindent\textbf{2：不適切} 容器三品目の全費用を販売者。式・定義・対応の独立判定と一致しない。\par

\noindent\textbf{3：適切・選択対象} 家電小売の引取・製造者への引渡し。判定と一致する。\par

\noindent\textbf{4：不適切} 全建設工事の発注者が分別再資源化。式・定義・対応の独立判定と一致しない。\par

\noindent\textbf{5：不適切} 再利用なら有害物輸出規制なし。式・定義・対応の独立判定と一致しない。\par

\paragraph{再利用する考え方} 法の目的、対象条件、排出者・自治体・小売・製造者の役割。 指定年度・対象範囲・単位・否定を先に固定し、他年度の番号や現行値を転用しない。

\paragraph{共通理論} \hyperref[theory:env:recyclelaw]{第\ref*{theory:env:recyclelaw}節}

\par 参考：\cite{env_circular_law,env_appliance}。

\par 難易度 3（校準初稿）、優先度未評価。

\subsection{令和元年度再試験 問5-2}
\label{q:r01-ret-basic-5-2}

\noindent\textbf{原典：} sources/R01\_retest/basic.pdf、PDF物理18頁。

\textbf{公式正答：} 1（single）。

\par\textbf{検討未確定：} 出題時バーゼル・建設の法令履歴と規模条件全文は追加照合。

\paragraph{問題の要旨} 家電小売の役割と循環制度の限定。

\Needspace{4\baselineskip}\paragraph{確認した条件・記述（独自要約）}\mbox{}\par

\noindent\textbf{時点：} 令和元年度再試験に指定された資料・制度。\par

\noindent\textbf{照合：} 原典18画像で否定・五肢・組合せ表を確認。後年値で置換しない。\par

\paragraph{問われる核心} 家電小売の役割と循環制度の限定。

\paragraph{必要知識} 役割、対象と例外、優先順位、国境移動。

\paragraph{考え方・正答を選ぶ理由（査読待ち初稿）} 1が家電小売の役割。再利用の有害廃棄物取引も規制対象となり得る。容器包装は三品目・販売者全費用ではなく、建設には対象条件と実施主体、循環の優先順には発生抑制がある。

\Needspace{4\baselineskip}\paragraph{選択肢の検討（査読待ち初稿）}\mbox{}\par

\noindent\textbf{1：適切・選択対象} 家電小売の引取と製造者への引渡し。判定と一致する。\par

\noindent\textbf{2：不適切} 再利用ならバーゼル規制なし。式・定義・対応の独立判定と一致しない。\par

\noindent\textbf{3：不適切} 容器三品目を販売者全費用。式・定義・対応の独立判定と一致しない。\par

\noindent\textbf{4：不適切} 全工事の発注者に分別実施。式・定義・対応の独立判定と一致しない。\par

\noindent\textbf{5：不適切} 基本法は再生最優先。式・定義・対応の独立判定と一致しない。\par

\paragraph{再利用する考え方} 役割、対象と例外、優先順位、国境移動。 集計対象・時点・法の主体・選択方向を固定し、同型でも元の五肢へ戻す。

\paragraph{共通理論} \hyperref[theory:env:recyclelaw]{第\ref*{theory:env:recyclelaw}節}

\par 参考：\cite{env_circular_law,env_appliance}。

\par 難易度 3（校準初稿）、優先度未評価。

\section{資源量と太陽光面積の桁を比較}

\subsection{平成26年度 問5-4}
\label{q:h26-reg-basic-5-4}

\noindent\textbf{原典：} sources/H26/basic.pdf、PDF物理17頁。

\textbf{公式正答：} 4（single）。

\par\textbf{検討未確定：} 概算の日射・需要は桁比較であり厳密な上限証明ではない。石炭石油の指定時点埋蔵量、海水U総量、全球純生産と人類エネルギーの各原表は追加照合。初回記録の昼夜連続1kW仮定を解説では日照条件を示す概算へ改めた。

\paragraph{問題の要旨} 資源量と太陽光面積の桁を比較。

\Needspace{4\baselineskip}\paragraph{確認した条件・記述（独自要約）}\mbox{}\par

\noindent\textbf{時点：} 平成26年度に指定された資料・制度。\par

\noindent\textbf{照合：} 原典21画像で否定・五肢・組合せ表を確認。後年値で置換しない。\par

\paragraph{問われる核心} 資源量と太陽光面積の桁を比較。

\paragraph{必要知識} 面積換算、電力とエネルギー、存在量と可採埋蔵量。

\paragraph{考え方・正答を選ぶ理由（査読待ち初稿）} 100km\textsuperscript{2}=10\textasciicircum{}8m\textsuperscript{2}。例えば1kW/m\textsuperscript{2}を昼間平均12h/日、効率を非現実的に100\%としても年約438TWhで、当時の日本の年間電力需要約1000TWhに届かない。現実には夜間・天候・入射角・変換損失でさらに小さいので4。資源の存在量と経済的に採れる確認埋蔵量も同じではなく、後者は探鉱・価格・技術で増え得る。

\Needspace{4\baselineskip}\paragraph{選択肢の検討（査読待ち初稿）}\mbox{}\par

\noindent\textbf{1：条件内で適切} 石炭埋蔵エネルギーが石油超。原典の条件内では適切な記述であり、不適切なものを選ぶ問の選択対象ではない。\par

\noindent\textbf{2：条件内で適切} 石油確認埋蔵量は危機時より増。原典の条件内では適切な記述であり、不適切なものを選ぶ問の選択対象ではない。\par

\noindent\textbf{3：条件内で適切} 海水Uは1億t超。原典の条件内では適切な記述であり、不適切なものを選ぶ問の選択対象ではない。\par

\noindent\textbf{4：不適切・選択対象} 100km\textsuperscript{2}太陽光が日本需要超。判定と一致する。\par

\noindent\textbf{5：条件内で適切} 全球の光合成固定が人類需要超。原典の条件内では適切な記述であり、不適切なものを選ぶ問の選択対象ではない。\par

\paragraph{再利用する考え方} 面積換算、電力とエネルギー、存在量と可採埋蔵量。 指定年度・対象範囲・単位・否定を先に固定し、他年度の番号や現行値を転用しない。

\paragraph{共通理論} \hyperref[theory:env:resource]{第\ref*{theory:env:resource}節}

\par 難易度 3（校準初稿）、優先度未評価。

\section{下水処理の段階と環境対策}

\subsection{平成25年度 問5-4}
\label{q:h25-reg-basic-5-4}

\noindent\textbf{原典：} sources/H25/basic.pdf、PDF物理21頁。

\textbf{公式正答：} 2（single）。

\par\textbf{検討未確定：} 当時土壌対策の掘削除去割合と施設の詳細法定仕様は追加照合。緑化の効果は気象・水・建物条件に依存。

\paragraph{問題の要旨} 下水処理の段階と環境対策。

\Needspace{4\baselineskip}\paragraph{確認した条件・記述（独自要約）}\mbox{}\par

\noindent\textbf{時点：} 平成25年度に指定された資料・制度。\par

\noindent\textbf{照合：} 原典21画像で否定・五肢・組合せ表を確認。後年値で置換しない。\par

\paragraph{問われる核心} 下水処理の段階と環境対策。

\paragraph{必要知識} 一次・二次・高度処理、物理と生物反応、汚染経路。

\paragraph{考え方・正答を選ぶ理由（査読待ち初稿）} 活性汚泥による生物処理は一般に二次処理。一次は沈殿等の物理処理であり、2が不適切。管理型処分場は浸出水等による汚染を防ぐ。緑化は遮蔽・蒸発散、ダイオキシン対策は燃焼と再合成防止の工程を分けて説明する。

\Needspace{4\baselineskip}\paragraph{選択肢の検討（査読待ち初稿）}\mbox{}\par

\noindent\textbf{1：条件内で適切} 管理型の遮水・浸出水処理。原典の条件内では適切な記述であり、不適切なものを選ぶ問の選択対象ではない。\par

\noindent\textbf{2：不適切・選択対象} 活性汚泥は一次処理。判定と一致する。\par

\noindent\textbf{3：条件内で適切} 緑化の表面温度と空調低減。原典の条件内では適切な記述であり、不適切なものを選ぶ問の選択対象ではない。\par

\noindent\textbf{4：条件内で適切} 土壌対策に掘削除去が多い。原典の条件内では適切な記述であり、不適切なものを選ぶ問の選択対象ではない。\par

\noindent\textbf{5：条件内で適切} 完全燃焼と排ガス急冷。原典の条件内では適切な記述であり、不適切なものを選ぶ問の選択対象ではない。\par

\paragraph{再利用する考え方} 一次・二次・高度処理、物理と生物反応、汚染経路。 指定年度・対象範囲・単位・否定を先に固定し、他年度の番号や現行値を転用しない。

\paragraph{共通理論} \hyperref[theory:env:treatment]{第\ref*{theory:env:treatment}節}

\par 難易度 3（校準初稿）、優先度未評価。

\subsection{令和3年度 問5-2}
\label{q:r03-reg-basic-5-2}

\noindent\textbf{原典：} sources/R03/basic.pdf、PDF物理22頁。

\textbf{公式正答：} 5（single）。

\par\textbf{検討未確定：} 土壌汚染対策実績の時点別割合、管理型処分場設備の当時法令条文は追加照合。

\paragraph{問題の要旨} 活性汚泥による生物処理は二次処理。

\Needspace{4\baselineskip}\paragraph{確認した条件・記述（独自要約）}\mbox{}\par

\noindent\textbf{原典：} 全五肢と否定方向・組合せ表を画像確認。\par

\noindent\textbf{資料時点：} 令和3年度の指定資料・統計年を保持し、後年値に置換しない。\par

\paragraph{問われる核心} 活性汚泥による生物処理は二次処理。

\paragraph{必要知識} 一次・二次・高度処理、排ガス急冷、遮水と浸出水。

\paragraph{考え方・正答を選ぶ理由（査読待ち初稿）} 5が不適切。沈砂・沈殿で浮遊固形物等を除く段階と、微生物が有機物を分解する段階を分ける。高度な窒素・リン除去等と生物処理一般の二次処理を混同しない。

\Needspace{4\baselineskip}\paragraph{選択肢の検討（査読待ち初稿）}\mbox{}\par

\noindent\textbf{1：条件内で適切} 焼却排ガスを200℃以下へ急冷。完全燃焼と急冷により生成・再合成を抑える対策。\par

\noindent\textbf{2：条件内で適切} 緑化により熱侵入を抑制。日射遮蔽や蒸散により建物への熱負荷を下げる。\par

\noindent\textbf{3：条件内で適切} 管理型処分場に遮水・浸出水処理。浸出水の環境放出を管理する設備が必要。\par

\noindent\textbf{4：条件内で適切} 原位置浄化もあるが掘削除去実績が多い。原位置は掘削せず処理する意味。実績の多寡は当時資料に依存。\par

\noindent\textbf{5：不適切・選択対象} 活性汚泥は一次処理。一次処理は沈殿等の物理的処理。活性汚泥の生物処理は一般に二次処理である。\par

\paragraph{再利用する考え方} 一次・二次・高度処理、排ガス急冷、遮水と浸出水。 否定方向と資料年、分母と単位を確認して五肢へ戻る。

\paragraph{共通理論} \hyperref[theory:env:treatment]{第\ref*{theory:env:treatment}節}

\par 難易度 3（校準初稿）、優先度未評価。

\section{温対法の排出抑制と吸収作用}

\subsection{平成27年度 問5-2}
\label{q:h27-reg-basic-5-2}

\noindent\textbf{原典：} sources/H27/basic.pdf、PDF物理15頁。

\textbf{公式正答：} 1（single）。

\par\textbf{検討未確定：} 2015施行時の全五条項を履歴条文で一括確認する作業は未完了。現行の2050目標を当時条文へ追加しない。

\paragraph{問題の要旨} 温対法の排出抑制と吸収作用。

\Needspace{4\baselineskip}\paragraph{確認した条件・記述（独自要約）}\mbox{}\par

\noindent\textbf{時点：} 平成27年度に指定された資料・制度。\par

\noindent\textbf{照合：} 原典21画像で否定・五肢・組合せ表を確認。後年値で置換しない。\par

\paragraph{問われる核心} 温対法の排出抑制と吸収作用。

\paragraph{必要知識} 目的規定、吸収と排出、政令対象、計画。

\paragraph{考え方・正答を選ぶ理由（査読待ち初稿）} 目的は排出抑制だけでなく吸収作用の保全・強化も含む。言及しないという1が誤り。対象ガス、算定報告の主体・規模、地方推進体制、計画見直しは別の制度事項なので、吸収と排出の二側面を軸に確認する。

\Needspace{4\baselineskip}\paragraph{選択肢の検討（査読待ち初稿）}\mbox{}\par

\noindent\textbf{1：不適切・選択対象} 森林吸収の保全に言及なし。判定と一致する。\par

\noindent\textbf{2：条件内で適切} 対象七種の温室効果ガス。原典の条件内では適切な記述であり、不適切なものを選ぶ問の選択対象ではない。\par

\noindent\textbf{3：条件内で適切} 一定規模事業所別の算定報告。原典の条件内では適切な記述であり、不適切なものを選ぶ問の選択対象ではない。\par

\noindent\textbf{4：条件内で適切} 知事が推進員委嘱とセンター指定。原典の条件内では適切な記述であり、不適切なものを選ぶ問の選択対象ではない。\par

\noindent\textbf{5：条件内で適切} 少なくとも三年ごと計画検討。原典の条件内では適切な記述であり、不適切なものを選ぶ問の選択対象ではない。\par

\paragraph{再利用する考え方} 目的規定、吸収と排出、政令対象、計画。 指定年度・対象範囲・単位・否定を先に固定し、他年度の番号や現行値を転用しない。

\paragraph{共通理論} \hyperref[theory:env:warminglaw]{第\ref*{theory:env:warminglaw}節}

\par 難易度 3（校準初稿）、優先度未評価。

\section{排出抑制から適正処分までの優先順}

\subsection{平成23年度 問5-2}
\label{q:h23-reg-basic-5-2}

\noindent\textbf{原典：} sources/H23/basic.pdf、PDF物理16頁。

\textbf{公式正答：} 2（single）。

\par\textbf{検討未確定：} 取得した2001資料は基本方針案と施策実績であり、2010白書の引用箇所全文の一致は追加照合。

\paragraph{問題の要旨} 排出抑制から適正処分までの優先順。

\Needspace{4\baselineskip}\paragraph{確認した条件・記述（独自要約）}\mbox{}\par

\noindent\textbf{時点：} 平成23年度に指定された資料・制度。\par

\noindent\textbf{照合：} 原典21画像で否定・五肢・組合せ表を確認。後年値で置換しない。\par

\paragraph{問われる核心} 排出抑制から適正処分までの優先順。

\paragraph{必要知識} 3R、循環利用と処分、基準年度と目標年度。

\paragraph{考え方・正答を選ぶ理由（査読待ち初稿）} 旧来の廃棄物行政の公衆衛生と公害対策に、循環型社会形成が加わる。優先は排出抑制→再使用→再生利用→熱回収→残るものの適正処分。2001基本方針の最終処分目標は1997基準の概ね半分なので2。処分を埋立の一方式だけへ狭めない。

\Needspace{4\baselineskip}\paragraph{選択肢の検討（査読待ち初稿）}\mbox{}\par

\noindent\textbf{1：不適切} 生活環境/再資源化/適正処理/半分。式・定義・対応の独立判定と一致しない。\par

\noindent\textbf{2：適切・選択対象} 公衆衛生/再使用/適正処分/半分。判定と一致する。\par

\noindent\textbf{3：不適切} 生活環境/再資源化/適正処分/1/3。式・定義・対応の独立判定と一致しない。\par

\noindent\textbf{4：不適切} 公衆衛生/再資源化/適正処理/1/3。式・定義・対応の独立判定と一致しない。\par

\noindent\textbf{5：不適切} 生活環境/再使用/適正埋立/1/3。式・定義・対応の独立判定と一致しない。\par

\paragraph{再利用する考え方} 3R、循環利用と処分、基準年度と目標年度。 指定年度・対象範囲・単位・否定を先に固定し、他年度の番号や現行値を転用しない。

\paragraph{共通理論} \hyperref[theory:env:wasteorder]{第\ref*{theory:env:wasteorder}節}

\par 参考：\cite{env_circular_law,env58_wastegoal}。

\par 難易度 3（校準初稿）、優先度未評価。

\section{廃棄物用語の略称と対象}

\subsection{平成27年度 問5-1}
\label{q:h27-reg-basic-5-1}

\noindent\textbf{原典：} sources/H27/basic.pdf、PDF物理15頁。

\textbf{公式正答：} 2（single）。

\par\textbf{検討未確定：} 当時のRDF添加法・産廃業種条件・硫酸ピッチのガス生成に関する技術資料を追加照合。全RDFが同一添加処方という一般化をしない。

\paragraph{問題の要旨} 廃棄物用語の略称と対象。

\Needspace{4\baselineskip}\paragraph{確認した条件・記述（独自要約）}\mbox{}\par

\noindent\textbf{時点：} 平成27年度に指定された資料・制度。\par

\noindent\textbf{照合：} 原典21画像で否定・五肢・組合せ表を確認。後年値で置換しない。\par

\paragraph{問われる核心} 廃棄物用語の略称と対象。

\paragraph{必要知識} 略称、法定分類、処理特性と有害性。

\paragraph{考え方・正答を選ぶ理由（査読待ち初稿）} E-wasteはelectrical/electronic waste、廃電気電子機器。環境に良い又はほぼ全量再資源化可能という定義ではなく2が誤り。RDFは可燃廃物の加工燃料、バイオマスは化石資源を除く。産業廃棄物は事業活動由来なら全てという分類でなく法定種類と業種条件を確認する。

\Needspace{4\baselineskip}\paragraph{選択肢の検討（査読待ち初稿）}\mbox{}\par

\noindent\textbf{1：条件内で適切} RDFはごみ固形化燃料。原典の条件内では適切な記述であり、不適切なものを選ぶ問の選択対象ではない。\par

\noindent\textbf{2：不適切・選択対象} E-wasteはほぼ全て再資源化できるエコ廃物。判定と一致する。\par

\noindent\textbf{3：条件内で適切} バイオマスは化石除く生物資源。原典の条件内では適切な記述であり、不適切なものを選ぶ問の選択対象ではない。\par

\noindent\textbf{4：条件内で適切} 産廃は事業由来の法定20種類。原典の条件内では適切な記述であり、不適切なものを選ぶ問の選択対象ではない。\par

\noindent\textbf{5：条件内で適切} 硫酸ピッチは強酸性油泥。原典の条件内では適切な記述であり、不適切なものを選ぶ問の選択対象ではない。\par

\paragraph{再利用する考え方} 略称、法定分類、処理特性と有害性。 指定年度・対象範囲・単位・否定を先に固定し、他年度の番号や現行値を転用しない。

\paragraph{共通理論} \hyperref[theory:env:wasteterm]{第\ref*{theory:env:wasteterm}節}

\par 難易度 3（校準初稿）、優先度未評価。

\subsection{令和4年度 問5-2}
\label{q:r04-reg-basic-5-2}

\noindent\textbf{原典：} sources/R04/basic.pdf、PDF物理17頁。

\textbf{公式正答：} 1（single）。

\par\textbf{検討未確定：} 当時の全国両区分の総量とメダル事業の回収・製造記録全文を追加確認。

\paragraph{問題の要旨} 一般廃棄物と産業廃棄物の総量を逆転しない。

\Needspace{4\baselineskip}\paragraph{確認した条件・記述（独自要約）}\mbox{}\par

\noindent\textbf{原典：} 全五肢と否定方向・組合せ表を画像確認。\par

\noindent\textbf{資料時点：} 令和4年度の指定資料・統計年を保持し、後年値に置換しない。\par

\paragraph{問われる核心} 一般廃棄物と産業廃棄物の総量を逆転しない。

\paragraph{必要知識} 廃棄物区分・量・性状、再生燃料、資源回収。

\paragraph{考え方・正答を選ぶ理由（査読待ち初稿）} 1が不適切。量は区分別の全国年間排出で比較する。物質の有害性を区分する特別管理、資源由来のバイオマス、再生燃料RPF、回収材利用はそれぞれ別の分類軸である。

\Needspace{4\baselineskip}\paragraph{選択肢の検討（査読待ち初稿）}\mbox{}\par

\noindent\textbf{1：不適切・選択対象} 一般廃棄物の総量が産業廃棄物より多い。同じ近年の年間排出総量では産業廃棄物が一般廃棄物より大きく、逆の記述が不適切。\par

\noindent\textbf{2：条件内で適切} 特別管理産業廃棄物は爆発性・毒性・感染性等。危険性に応じ通常とは別の管理が必要となる。\par

\noindent\textbf{3：条件内で適切} バイオマスは化石資源を除く生物由来有機資源。生物起源でも長期地質過程で生じた化石燃料はここに含めない。\par

\noindent\textbf{4：条件内で適切} RPFは主に廃紙・廃プラスチック固形燃料。略称と対象材料を照合。RDF等の別燃料を混同しない。\par

\noindent\textbf{5：条件内で適切} 東京大会メダルは使用済小型家電の金属を利用。回収金属を資源として使う事業であり、新規採掘だけで製造したとの説明ではない。\par

\paragraph{再利用する考え方} 廃棄物区分・量・性状、再生燃料、資源回収。 否定方向と資料年、分母と単位を確認して五肢へ戻る。

\paragraph{共通理論} \hyperref[theory:env:wasteterm]{第\ref*{theory:env:wasteterm}節}

\par 難易度 3（校準初稿）、優先度未評価。

\section{工業型から都市型の汚染と二次生成}

\subsection{令和元年度 問5-1}
\label{q:r01-reg-basic-5-1}

\noindent\textbf{原典：} sources/R01/basic.pdf、PDF物理19頁。

\textbf{公式正答：} 3（single）。

\paragraph{問題の要旨} 工業型から都市型の汚染と二次生成。

\Needspace{4\baselineskip}\paragraph{確認した条件・記述（独自要約）}\mbox{}\par

\noindent\textbf{時点：} 令和元年度に指定された資料・制度。\par

\noindent\textbf{照合：} 原典18画像で否定・五肢・組合せ表を確認。後年値で置換しない。\par

\paragraph{問われる核心} 工業型から都市型の汚染と二次生成。

\paragraph{必要知識} 発生源、硫黄酸化物、粒子、光化学反応。

\paragraph{考え方・正答を選ぶ理由（査読待ち初稿）} 工業地域の燃料硫黄によるSOx、自動車のNOxとSPM、NOxと炭化水素の光反応から生じる光化学オキシダント。SOx/NOx/SPM/Oxの3。最後の二次生成物をCO2とすると前駆物質と日射の条件に対応しない。

\Needspace{4\baselineskip}\paragraph{選択肢の検討（査読待ち初稿）}\mbox{}\par

\noindent\textbf{1：不適切} SOx/Ox/SPM/CO2。式・定義・対応の独立判定と一致しない。\par

\noindent\textbf{2：不適切} NOx/Ox/CO2/SPM。式・定義・対応の独立判定と一致しない。\par

\noindent\textbf{3：適切・選択対象} SOx/NOx/SPM/Ox。判定と一致する。\par

\noindent\textbf{4：不適切} NOx/SOx/CO2/Ox。式・定義・対応の独立判定と一致しない。\par

\noindent\textbf{5：不適切} SOx/NOx/SPM/CO2。式・定義・対応の独立判定と一致しない。\par

\paragraph{再利用する考え方} 発生源、硫黄酸化物、粒子、光化学反応。 集計対象・時点・法の主体・選択方向を固定し、同型でも元の五肢へ戻す。

\paragraph{共通理論} \hyperref[theory:env:airhistory]{第\ref*{theory:env:airhistory}節}

\par 参考：\cite{env58_pm25}。

\par 難易度 2（校準初稿）、優先度未評価。

\section{緩和・適応とEco-DRRの共便益}

\subsection{令和元年度再試験 問5-1}
\label{q:r01-ret-basic-5-1}

\noindent\textbf{原典：} sources/R01\_retest/basic.pdf、PDF物理17頁。

\textbf{公式正答：} 1（single）。

\paragraph{問題の要旨} 緩和・適応とEco-DRRの共便益。

\Needspace{4\baselineskip}\paragraph{確認した条件・記述（独自要約）}\mbox{}\par

\noindent\textbf{時点：} 令和元年度再試験に指定された資料・制度。\par

\noindent\textbf{照合：} 原典18画像で否定・五肢・組合せ表を確認。後年値で置換しない。\par

\paragraph{問われる核心} 緩和・適応とEco-DRRの共便益。

\paragraph{必要知識} 緩和適応、地域特性、生態系の機能、共便益。

\paragraph{考え方・正答を選ぶ理由（査読待ち初稿）} 排出原因を減らす緩和、避けられない影響に備える適応。湿原の遊水や森林の保全機能を使うEco-DRRは生態系を活用する減災で、生物多様性保全も共便益となる。四語を合わせる1。人材が不要という意味ではなく、略称が指す手法の核心を読む。

\Needspace{4\baselineskip}\paragraph{選択肢の検討（査読待ち初稿）}\mbox{}\par

\noindent\textbf{1：適切・選択対象} 緩和/適応/生態系/生物多様性。判定と一致する。\par

\noindent\textbf{2：不適切} 削減/対応/生態系/地域資源。式・定義・対応の独立判定と一致しない。\par

\noindent\textbf{3：不適切} 緩和/適応/地域人材/地域資源。式・定義・対応の独立判定と一致しない。\par

\noindent\textbf{4：不適切} 緩和/対応/生態系/生物多様性。式・定義・対応の独立判定と一致しない。\par

\noindent\textbf{5：不適切} 削減/対応/地域人材/地域資源。式・定義・対応の独立判定と一致しない。\par

\paragraph{再利用する考え方} 緩和適応、地域特性、生態系の機能、共便益。 集計対象・時点・法の主体・選択方向を固定し、同型でも元の五肢へ戻す。

\paragraph{共通理論} \hyperref[theory:env:ecodrr]{第\ref*{theory:env:ecodrr}節}

\par 参考：\cite{env59_ecodrr}。

\par 難易度 2（校準初稿）、優先度未評価。

\section{部門別エネルギーの歴史的推移}

\subsection{令和2年度 問5-3}
\label{q:r02-reg-basic-5-3}

\noindent\textbf{原典：} sources/R02/basic.pdf、PDF物理20頁。

\textbf{公式正答：} 4（single）。

\par\textbf{検討未確定：} 原典が参照する版のピーク年2005・2002と全五グラフの原表は未取得。後年の統計遡及改訂と区別して追加照合。

\paragraph{問題の要旨} 部門別エネルギーの歴史的推移。

\Needspace{4\baselineskip}\paragraph{確認した条件・記述（独自要約）}\mbox{}\par

\noindent\textbf{時点：} 令和2年度に指定された資料・制度。\par

\noindent\textbf{照合：} 原典18画像で否定・五肢・組合せ表を確認。後年値で置換しない。\par

\paragraph{問われる核心} 部門別エネルギーの歴史的推移。

\paragraph{必要知識} 時系列、部門、GDP原単位、複数要因。

\paragraph{考え方・正答を選ぶ理由（査読待ち初稿）} 家庭は省エネ機器・節電・世帯人員・気温等の影響で、震災後に継続増加したという4が誤り。2019統計の要因分解にも年ごとの増減があり、一様な増加ではない。総量、原単位、世帯数、産業部門の個別と合計を分ける。

\Needspace{4\baselineskip}\paragraph{選択肢の検討（査読待ち初稿）}\mbox{}\par

\noindent\textbf{1：条件内で適切} 最終消費2005ピーク後減。条件内で適切な記述であり、不適切な記述を選ぶ対象ではない。\par

\noindent\textbf{2：条件内で適切} 産業業務全体は省エネで同程度。条件内で適切な記述であり、不適切な記述を選ぶ対象ではない。\par

\noindent\textbf{3：条件内で適切} GDP当たり日本は世界平均未満。条件内で適切な記述であり、不適切な記述を選ぶ対象ではない。\par

\noindent\textbf{4：不適切・選択対象} 家庭は震災後も継続増加。判定と一致する。\par

\noindent\textbf{5：条件内で適切} 旅客2002ピーク後燃費等で減。条件内で適切な記述であり、不適切な記述を選ぶ対象ではない。\par

\paragraph{再利用する考え方} 時系列、部門、GDP原単位、複数要因。 集計対象・時点・法の主体・選択方向を固定し、同型でも元の五肢へ戻す。

\paragraph{共通理論} \hyperref[theory:env:energytrends]{第\ref*{theory:env:energytrends}節}

\par 参考：\cite{env59_energy2019}。

\par 難易度 3（校準初稿）、優先度未評価。

\section{人物の業績とサイクロトロンの発明者}

\subsection{平成29年度 問5-6}
\label{q:h29-reg-basic-5-6}

\noindent\textbf{原典：} sources/H29/basic.pdf、PDF物理20頁。

\textbf{公式正答：} 5（single）。

\par\textbf{検討未確定：} Lawrenceの公式授賞説明を確認。他の四人の全原報は未取得で業績対応は一次史料追加照合を要する。

\paragraph{問題の要旨} 人物の業績とサイクロトロンの発明者。

\Needspace{4\baselineskip}\paragraph{確認した条件・記述（独自要約）}\mbox{}\par

\noindent\textbf{時点：} 平成29年度に指定された資料・制度。\par

\noindent\textbf{照合：} 原典18画像で否定・五肢・組合せ表を確認。後年値で置換しない。\par

\paragraph{問われる核心} 人物の業績とサイクロトロンの発明者。

\paragraph{必要知識} 電池、抗菌物質、放射線、遺伝、加速器。

\paragraph{考え方・正答を選ぶ理由（査読待ち初稿）} サイクロトロンはLawrenceらが開発した加速器で、Edisonという5が不適切。Voltaは定常電流を得る電堆、Flemingはリゾチームとペニシリン、RöntgenはX線、Mendelは形質の遺伝の規則性という対応。発明者名の誤りを、他の正しい加速原理の記述と混ぜて見落とさない。

\Needspace{4\baselineskip}\paragraph{選択肢の検討（査読待ち初稿）}\mbox{}\par

\noindent\textbf{1：条件内で適切} Volta電堆。条件内で適切な記述であり、不適切な記述を選ぶ対象ではない。\par

\noindent\textbf{2：条件内で適切} Flemingリゾチームとペニシリン。条件内で適切な記述であり、不適切な記述を選ぶ対象ではない。\par

\noindent\textbf{3：条件内で適切} RöntgenX線。条件内で適切な記述であり、不適切な記述を選ぶ対象ではない。\par

\noindent\textbf{4：条件内で適切} Mendel遺伝法則。条件内で適切な記述であり、不適切な記述を選ぶ対象ではない。\par

\noindent\textbf{5：不適切・選択対象} Edisonサイクロトロン。判定と一致する。\par

\paragraph{再利用する考え方} 電池、抗菌物質、放射線、遺伝、加速器。 集計対象・時点・法の主体・選択方向を固定し、同型でも元の五肢へ戻す。

\paragraph{共通理論} \hyperref[theory:env:historydiscover]{第\ref*{theory:env:historydiscover}節}

\par 参考：\cite{env59_lawrence}。

\par 難易度 3（校準初稿）、優先度未評価。

\section{家庭エネルギーの用途と消費範囲}

\subsection{平成29年度 問5-4}
\label{q:h29-reg-basic-5-4}

\noindent\textbf{原典：} sources/H29/basic.pdf、PDF物理19頁。

\textbf{公式正答：} 2（single）。

\par\textbf{検討未確定：} 原典「近年」の指定統計年が曖昧で、給湯三割・電気五割・待機五\%の全原表は未取得。比率を現行値の普遍定数にはしない。

\paragraph{問題の要旨} 家庭エネルギーの用途と消費範囲。

\Needspace{4\baselineskip}\paragraph{確認した条件・記述（独自要約）}\mbox{}\par

\noindent\textbf{時点：} 平成29年度に指定された資料・制度。\par

\noindent\textbf{照合：} 原典18画像で否定・五肢・組合せ表を確認。後年値で置換しない。\par

\paragraph{問われる核心} 家庭エネルギーの用途と消費範囲。

\paragraph{必要知識} 用途別消費、エネルギー種、機器効率、待機。

\paragraph{考え方・正答を選ぶ理由（査読待ち初稿）} 全国総和では冷房の年間消費は暖房より小さく、十倍とする2が不適切。電力だけと全燃料を含む家庭エネルギー、機器の効率と使用時間、総量と待機電力の比率を分ける。トップランナーは将来の効率目標の設定方式で、使用する一台ごとの達成義務と同じではない。

\Needspace{4\baselineskip}\paragraph{選択肢の検討（査読待ち初稿）}\mbox{}\par

\noindent\textbf{1：条件内で適切} 給湯が約三割。条件内で適切な記述であり、不適切な記述を選ぶ対象ではない。\par

\noindent\textbf{2：不適切・選択対象} 冷房が暖房の約十倍。判定と一致する。\par

\noindent\textbf{3：条件内で適切} 電気が約五割。条件内で適切な記述であり、不適切な記述を選ぶ対象ではない。\par

\noindent\textbf{4：条件内で適切} 冷蔵庫TV空調はトップランナー。条件内で適切な記述であり、不適切な記述を選ぶ対象ではない。\par

\noindent\textbf{5：条件内で適切} 電力約5\%が待機。条件内で適切な記述であり、不適切な記述を選ぶ対象ではない。\par

\paragraph{再利用する考え方} 用途別消費、エネルギー種、機器効率、待機。 集計対象・時点・法の主体・選択方向を固定し、同型でも元の五肢へ戻す。

\paragraph{共通理論} \hyperref[theory:env:home]{第\ref*{theory:env:home}節}

\par 参考：\cite{env59_energy2019}。

\par 難易度 3（校準初稿）、優先度未評価。

\section{戦前技術史の曖昧な比較と公式複数正答}

\subsection{令和2年度 問5-5}
\label{q:r02-reg-basic-5-5}

\noindent\textbf{原典：} sources/R02/basic.pdf、PDF物理21頁。

\textbf{公式正答：} 4|3（multiple\_accepted）。

\par\textbf{検討未確定：} 公式の複数受理4|3を保持。独立解4は受理対象だが全受理集合を独立に列挙したものではない。3の欧米との差という一般化、戦前統計利用の個別例を追加史料で査読。

\paragraph{問題の要旨} 戦前技術史の曖昧な比較と公式複数正答。

\Needspace{4\baselineskip}\paragraph{確認した条件・記述（独自要約）}\mbox{}\par

\noindent\textbf{時点：} 令和2年度に指定された資料・制度。\par

\noindent\textbf{照合：} 原典18画像で否定・五肢・組合せ表を確認。後年値で置換しない。\par

\paragraph{問われる核心} 戦前技術史の曖昧な比較と公式複数正答。

\paragraph{必要知識} 技術移転、学協会史、科学的管理と品質管理、公式訂正。

\paragraph{考え方・正答を選ぶ理由（査読待ち初稿）} 独立の主な判断は4。日本で統計的品質管理が組織的に普及した節目は日科技連による1950のDeming講義等の戦後であり、戦前のTaylor導入と同時に普及したというまとめは不適切。ただし戦前の個別統計手法の利用まで無かったとはしない。公式は誤解を招く記述があったとして当初4に加え3も正答とした。3の「欧米とは異なり」という一般化は学協会の時期・構成の違いを落とす。独立解4、受理正答3と4をそれぞれ残す。

\Needspace{4\baselineskip}\paragraph{選択肢の検討（査読待ち初稿）}\mbox{}\par

\noindent\textbf{1：条件内で適切} 江戸手工業が西洋技術と発展。条件内で適切な記述であり、不適切な記述を選ぶ対象ではない。\par

\noindent\textbf{2：条件内で適切} 標準化技術を輸入し低賃金競争。条件内で適切な記述であり、不適切な記述を選ぶ対象ではない。\par

\noindent\textbf{3：不適切・公式追加受理} 学協会は欧米と異なり高学歴出身者。欧米との比較表現に誤解を招く記述があるとして公式が3も受理。4との両受理を維持。\par

\noindent\textbf{4：不適切・選択対象} Taylorと同時に統計的品質管理導入。判定と一致する。\par

\noindent\textbf{5：条件内で適切} 技術官僚の地位向上運動。条件内で適切な記述であり、不適切な記述を選ぶ対象ではない。\par

\paragraph{再利用する考え方} 技術移転、学協会史、科学的管理と品質管理、公式訂正。 集計対象・時点・法の主体・選択方向を固定し、同型でも元の五肢へ戻す。

\paragraph{共通理論} \hyperref[theory:env:japanindustrial]{第\ref*{theory:env:japanindustrial}節}

\par 参考：\cite{env59_juse}。

\par 難易度 3（校準初稿）、優先度未評価。

\section{非エネルギー利用を除いた炭素量からCO2へ}

\subsection{令和元年度 問5-4}
\label{q:r01-reg-basic-5-4}

\noindent\textbf{原典：} sources/R01/basic.pdf、PDF物理21頁。

\textbf{公式正答：} 4（single）。

\paragraph{問題の要旨} 非エネルギー利用を除いた炭素量からCO2へ。

\Needspace{4\baselineskip}\paragraph{確認した条件・記述（独自要約）}\mbox{}\par

\noindent\textbf{燃料PJ：} 石炭5044、石油7831、天然ガス都市ガス4696。\par

\noindent\textbf{非燃焼石油：} 約1600PJ。\par

\noindent\textbf{係数t-C/TJ：} 石炭24、石油19、ガス14。\par

\noindent\textbf{換算：} C12、CO2 44、PJ=\allowbreak{}1000TJ。\par

\paragraph{問われる核心} 非エネルギー利用を除いた炭素量からCO2へ。

\paragraph{必要知識} システム境界、非エネルギー利用、PJ/TJ、分子量換算。

\paragraph{考え方・正答を選ぶ理由（査読待ち初稿）} 石油は燃焼に使わない1600PJを引いて6231PJ。炭素排出は[5044×24+6231×19+4696×14]×1000=\allowbreak{}305189000t-C。PJ→TJの1000倍を忘れない。CO2は44/12倍で約1119Mt-CO2なので1100の4。炭素換算値305Mtをそのまま選ばない。

\Needspace{4\baselineskip}\paragraph{選択肢の検討（査読待ち初稿）}\mbox{}\par

\noindent\textbf{1：不適切} 100Mt。式・定義・対応の独立判定と一致しない。\par

\noindent\textbf{2：不適切} 300Mt。式・定義・対応の独立判定と一致しない。\par

\noindent\textbf{3：不適切} 500Mt。式・定義・対応の独立判定と一致しない。\par

\noindent\textbf{4：適切・選択対象} 1100Mt。判定と一致する。\par

\noindent\textbf{5：不適切} 1600Mt。式・定義・対応の独立判定と一致しない。\par

\paragraph{再利用する考え方} システム境界、非エネルギー利用、PJ/TJ、分子量換算。 集計対象・時点・法の主体・選択方向を固定し、同型でも元の五肢へ戻す。

\paragraph{共通理論} \hyperref[theory:env:nationalco2]{第\ref*{theory:env:nationalco2}節}

\par 難易度 2（校準初稿）、優先度未評価。

\section{2017年の原油数量と輸入額比}

\subsection{平成30年度 問5-3}
\label{q:h30-reg-basic-5-3}

\noindent\textbf{原典：} sources/H30/basic.pdf、PDF物理19頁。

\textbf{公式正答：} 3（single）。

\par\textbf{検討未確定：} 2019白書の公式検索本文で中東約87\%を確認。2017財務省暦年の最大輸入先と金額12\%の全原表を取得していない。

\paragraph{問題の要旨} 2017年の原油数量と輸入額比。

\Needspace{4\baselineskip}\paragraph{確認した条件・記述（独自要約）}\mbox{}\par

\noindent\textbf{時点：} 平成30年度に指定された資料・制度。\par

\noindent\textbf{照合：} 原典18画像で否定・五肢・組合せ表を確認。後年値で置換しない。\par

\paragraph{問われる核心} 2017年の原油数量と輸入額比。

\paragraph{必要知識} 統計年度・品目・数量と金額・経路。

\paragraph{考え方・正答を選ぶ理由（査読待ち初稿）} 指定年の中東依存は約87\%、主要経路ホルムズ、最大相手サウジアラビア、石油・石油製品の総輸入額比約12\%で3。依存度は数量、最後は金額であり、他の化石燃料を含む鉱物性燃料の比率へ拡張しない。

\Needspace{4\baselineskip}\paragraph{選択肢の検討（査読待ち初稿）}\mbox{}\par

\noindent\textbf{1：不適切} 67/マラッカ/クウェート/12。式・定義・対応の独立判定と一致しない。\par

\noindent\textbf{2：不適切} 67/ホルムズ/サウジ/32。式・定義・対応の独立判定と一致しない。\par

\noindent\textbf{3：適切・選択対象} 87/ホルムズ/サウジ/12。判定と一致する。\par

\noindent\textbf{4：不適切} 87/マラッカ/クウェート/32。式・定義・対応の独立判定と一致しない。\par

\noindent\textbf{5：不適切} 87/ホルムズ/クウェート/12。式・定義・対応の独立判定と一致しない。\par

\paragraph{再利用する考え方} 統計年度・品目・数量と金額・経路。 集計対象・時点・法の主体・選択方向を固定し、同型でも元の五肢へ戻す。

\paragraph{共通理論} \hyperref[theory:env:oil2017]{第\ref*{theory:env:oil2017}節}

\par 参考：\cite{env59_energy2019}。

\par 難易度 3（校準初稿）、優先度未評価。

\section{パリ協定の削減目標は緩和に対応}

\subsection{平成29年度 問5-2}
\label{q:h29-reg-basic-5-2}

\noindent\textbf{原典：} sources/H29/basic.pdf、PDF物理18頁。

\textbf{公式正答：} 1（single）。

\par\textbf{検討未確定：} 取得UNFCCCはNDC五年と緩和の公式説明。協定2・5・6・13条の本文全体を通して照合する作業は未完了。

\paragraph{問題の要旨} パリ協定の削減目標は緩和に対応。

\Needspace{4\baselineskip}\paragraph{確認した条件・記述（独自要約）}\mbox{}\par

\noindent\textbf{時点：} 平成29年度に指定された資料・制度。\par

\noindent\textbf{照合：} 原典18画像で否定・五肢・組合せ表を確認。後年値で置換しない。\par

\paragraph{問われる核心} パリ協定の削減目標は緩和に対応。

\paragraph{必要知識} NDC、緩和と適応、温度目標、透明性、協力。

\paragraph{考え方・正答を選ぶ理由（査読待ち初稿）} 削減目標NDCの五年提出は緩和の取組であり、適応と呼ぶ1が不適切。適応は影響への備え。温度目標、透明性の情報・レビュー、協力的アプローチと市場、森林の吸収・排出抑制は条文の異なる要素を読む。パリ協定本文がJCMという日本の制度名を直接列挙したとの意味にはしない。

\Needspace{4\baselineskip}\paragraph{選択肢の検討（査読待ち初稿）}\mbox{}\par

\noindent\textbf{1：不適切・選択対象} 五年ごと削減目標で適応を推進。判定と一致する。\par

\noindent\textbf{2：条件内で適切} 2℃十分下方と1.5℃努力。条件内で適切な記述であり、不適切な記述を選ぶ対象ではない。\par

\noindent\textbf{3：条件内で適切} 情報の専門家レビュー。条件内で適切な記述であり、不適切な記述を選ぶ対象ではない。\par

\noindent\textbf{4：条件内で適切} 市場メカニズムの利用。条件内で適切な記述であり、不適切な記述を選ぶ対象ではない。\par

\noindent\textbf{5：条件内で適切} 森林減少・劣化抑制の支援。条件内で適切な記述であり、不適切な記述を選ぶ対象ではない。\par

\paragraph{再利用する考え方} NDC、緩和と適応、温度目標、透明性、協力。 集計対象・時点・法の主体・選択方向を固定し、同型でも元の五肢へ戻す。

\paragraph{共通理論} \hyperref[theory:env:paris]{第\ref*{theory:env:paris}節}

\par 参考：\cite{env59_ndc}。

\par 難易度 3（校準初稿）、優先度未評価。

\section{海洋プラスチックの定義・流出・循環}

\subsection{令和2年度 問5-1}
\label{q:r02-reg-basic-5-1}

\noindent\textbf{原典：} sources/R02/basic.pdf、PDF物理19頁。

\textbf{公式正答：} 2（single）。

\par\textbf{検討未確定：} 原典の2019年6月は政府発表5月31日と相違。厳密に日付も判定すると想定組合せが崩れるためAstra重点査読。2017国内排出・海外再生の全原表も追加照合。

\paragraph{問題の要旨} 海洋プラスチックの定義・流出・循環。

\Needspace{4\baselineskip}\paragraph{確認した条件・記述（独自要約）}\mbox{}\par

\noindent\textbf{時点：} 令和2年度に指定された資料・制度。\par

\noindent\textbf{照合：} 原典18画像で否定・五肢・組合せ表を確認。後年値で置換しない。\par

\paragraph{問われる核心} 海洋プラスチックの定義・流出・循環。

\paragraph{必要知識} 粒径、陸域流出、資源循環、代替資源、日付。

\paragraph{考え方・正答を選ぶ理由（査読待ち初稿）} 原典の主題判定はア正（概ね5mm以下）、イ誤（先進国だけ又は特に先進国由来へ一般化）、ウ誤（海外再生が半数以下という2017説明が違う）、エ正（3R+Renewable）、オ正（陸域の投棄や河川流出への対策）でOXXOOの2。ただし戦略の策定日は政府資料では2019年5月31日で、原典エの「6月」と一致しない。公式2を保持し、内容の判定と日付の疑義を分ける。

\Needspace{4\baselineskip}\paragraph{選択肢の検討（査読待ち初稿）}\mbox{}\par

\noindent\textbf{1：不適切} OOXOX。式・定義・対応の独立判定と一致しない。\par

\noindent\textbf{2：適切・選択対象} OXXOO。判定と一致する。\par

\noindent\textbf{3：不適切} OOOXX。式・定義・対応の独立判定と一致しない。\par

\noindent\textbf{4：不適切} XXOOO。式・定義・対応の独立判定と一致しない。\par

\noindent\textbf{5：不適切} XOXXO。式・定義・対応の独立判定と一致しない。\par

\paragraph{再利用する考え方} 粒径、陸域流出、資源循環、代替資源、日付。 集計対象・時点・法の主体・選択方向を固定し、同型でも元の五肢へ戻す。

\paragraph{共通理論} \hyperref[theory:env:plastics]{第\ref*{theory:env:plastics}節}

\par 参考：\cite{env59_plastic}。

\par 難易度 3（校準初稿）、優先度未評価。

\subsection{令和7年度 問5-2}
\label{q:r07-reg-basic-5-2}

\noindent\textbf{原典：} sources/R07/basic.pdf、PDF物理16頁。

\textbf{公式正答：} 3（single）。

\par\textbf{検討未確定：} 原典オは2019年6月と記すが政府戦略は2019年5月31日。公式選択肢3を保持し日付差を査読事項へ残す。2010推計と2017海外量の元表は追加照合。

\paragraph{問題の要旨} 海洋プラスチック五文は誤誤正誤正。

\Needspace{4\baselineskip}\paragraph{確認した条件・記述（独自要約）}\mbox{}\par

\noindent\textbf{原典：} 全五肢と否定方向・組合せ表を画像確認。\par

\noindent\textbf{資料時点：} 令和7年度の指定資料・統計年を保持し、後年値に置換しない。\par

\paragraph{問われる核心} 海洋プラスチック五文は誤誤正誤正。

\paragraph{必要知識} 微小プラスチック、流出と発生、国内循環と輸出、3R+Renewable。

\paragraph{考え方・正答を選ぶ理由（査読待ち初稿）} アは通常5mm未満を10mmとするため誤。イは2010流出推計を先進国主体とする一般化で誤。ウは陸のごみが川を通じ海へ流れるので陸域対策が有効で正。エは2017海外リサイクルを10万t未満とする桁が誤。オの3R+Renewableという内容は正として誤誤正誤正の3。ただし原典の策定月2019年6月は政府公表の5月31日と区別して留保する。

\Needspace{4\baselineskip}\paragraph{選択肢の検討（査読待ち初稿）}\mbox{}\par

\noindent\textbf{1：不適切} 正正誤誤正。独立に求めた組合せ 誤誤正誤正 と異なる。\par

\noindent\textbf{2：不適切} 正正誤正誤。独立に求めた組合せ 誤誤正誤正 と異なる。\par

\noindent\textbf{3：適切・選択対象} 誤誤正誤正。独立に求めた組合せ 誤誤正誤正 と一致する。\par

\noindent\textbf{4：不適切} 誤正正誤誤。独立に求めた組合せ 誤誤正誤正 と異なる。\par

\noindent\textbf{5：不適切} 誤誤正正正。独立に求めた組合せ 誤誤正誤正 と異なる。\par

\paragraph{再利用する考え方} 微小プラスチック、流出と発生、国内循環と輸出、3R+Renewable。 否定方向と資料年、分母と単位を確認して五肢へ戻る。

\paragraph{共通理論} \hyperref[theory:env:plastics]{第\ref*{theory:env:plastics}節}

\par 参考：\cite{env59_plastic}。

\par 難易度 3（校準初稿）、優先度未評価。

\section{一次供給の再エネ比率を境界に合わせる}

\subsection{令和元年度再試験 問5-4}
\label{q:r01-ret-basic-5-4}

\noindent\textbf{原典：} sources/R01\_retest/basic.pdf、PDF物理19頁。

\textbf{公式正答：} 3（single）。

\par\textbf{検討未確定：} 再試験の本文には内訳がないため通常試験の2017内訳を補助として使った。総合エネルギー統計指定版の原表は追加取得。

\paragraph{問題の要旨} 一次供給の再エネ比率を境界に合わせる。

\Needspace{4\baselineskip}\paragraph{確認した条件・記述（独自要約）}\mbox{}\par

\noindent\textbf{時点：} 令和元年度再試験に指定された資料・制度。\par

\noindent\textbf{照合：} 原典18画像で否定・五肢・組合せ表を確認。後年値で置換しない。\par

\paragraph{問われる核心} 一次供給の再エネ比率を境界に合わせる。

\paragraph{必要知識} 一次供給、発電比との違い、境界、合計と百分率。

\paragraph{考え方・正答を選ぶ理由（査読待ち初稿）} 2017指定統計と通常試験で与えた内訳を使うと、水力710+水力除く再エネ938+未活用596=\allowbreak{}2244PJ。全20095PJで割ると11.17\%で3。これは一次エネルギー供給比であり発電量構成ではない。問題が指定する未活用を含む境界を保持する。

\Needspace{4\baselineskip}\paragraph{選択肢の検討（査読待ち初稿）}\mbox{}\par

\noindent\textbf{1：不適切} 44\%。式・定義・対応の独立判定と一致しない。\par

\noindent\textbf{2：不適切} 22\%。式・定義・対応の独立判定と一致しない。\par

\noindent\textbf{3：適切・選択対象} 11\%。判定と一致する。\par

\noindent\textbf{4：不適切} 2\%。式・定義・対応の独立判定と一致しない。\par

\noindent\textbf{5：不適切} 0.5\%。式・定義・対応の独立判定と一致しない。\par

\paragraph{再利用する考え方} 一次供給、発電比との違い、境界、合計と百分率。 集計対象・時点・法の主体・選択方向を固定し、同型でも元の五肢へ戻す。

\paragraph{共通理論} \hyperref[theory:env:primaryshare]{第\ref*{theory:env:primaryshare}節}

\par 難易度 3（校準初稿）、優先度未評価。

\subsection{令和6年度 問5-4}
\label{q:r06-reg-basic-5-4}

\noindent\textbf{原典：} sources/R06/basic.pdf、PDF物理19頁。

\textbf{公式正答：} 4（single）。

\par\textbf{検討未確定：} 2022確報の数値表から指定三区分を実際に合算する追加再計算が必要。原表索引と再エネ説明を確認した範囲。

\paragraph{問題の要旨} 2022一次供給の再エネ・水力・未活用比。

\Needspace{4\baselineskip}\paragraph{確認した条件・記述（独自要約）}\mbox{}\par

\noindent\textbf{原典：} 全五肢と否定方向・組合せ表を画像確認。\par

\noindent\textbf{資料時点：} 令和6年度の指定資料・統計年を保持し、後年値に置換しない。\par

\paragraph{問われる核心} 2022一次供給の再エネ・水力・未活用比。

\paragraph{必要知識} 一次供給と発電量、揚水除外、区分集計。

\paragraph{考え方・正答を選ぶ理由（査読待ち初稿）} 一次エネルギー供給を分母とし、水力を除く再エネ、水力（揚水除外）、未活用エネルギーを合算すると、指定2022統計では約14\%として4を選ぶ。発電量における再エネ割合や自給率を代入しない。未活用には廃棄・放散される部分の有効利用等が含まれ、区分の重複計上を避ける。

\Needspace{4\baselineskip}\paragraph{選択肢の検討（査読待ち初稿）}\mbox{}\par

\noindent\textbf{1：不適切} 0.5\%。独立に求めた組合せ 14\% と異なる。\par

\noindent\textbf{2：不適切} 2\%。独立に求めた組合せ 14\% と異なる。\par

\noindent\textbf{3：不適切} 7\%。独立に求めた組合せ 14\% と異なる。\par

\noindent\textbf{4：適切・選択対象} 14\%。独立に求めた組合せ 14\% と一致する。\par

\noindent\textbf{5：不適切} 28\%。独立に求めた組合せ 14\% と異なる。\par

\paragraph{再利用する考え方} 一次供給と発電量、揚水除外、区分集計。 否定方向と資料年、分母と単位を確認して五肢へ戻る。

\paragraph{共通理論} \hyperref[theory:env:primaryshare]{第\ref*{theory:env:primaryshare}節}

\par 参考：\cite{env60_energy2022}。

\par 難易度 3（校準初稿）、優先度未評価。

\section{LCA・汚染者負担・拡大生産者責任}

\subsection{平成29年度 問5-1}
\label{q:h29-reg-basic-5-1}

\noindent\textbf{原典：} sources/H29/basic.pdf、PDF物理18頁。

\textbf{公式正答：} 1（single）。

\paragraph{問題の要旨} LCA・汚染者負担・拡大生産者責任。

\Needspace{4\baselineskip}\paragraph{確認した条件・記述（独自要約）}\mbox{}\par

\noindent\textbf{列順：} A=\allowbreak{}LCA、B=汚染者負担、C=拡大生産者、D=税務監査。\par

\noindent\textbf{O/X：} 正/誤。\par

\paragraph{問われる核心} LCA・汚染者負担・拡大生産者責任。

\paragraph{必要知識} ライフサイクル、汚染者負担原則、拡大生産者責任、監査。

\paragraph{考え方・正答を選ぶ理由（査読待ち初稿）} Aは全過程の影響評価でLCA。Bは汚染の防止・回復費用を汚染者が負担するPPP。Cは廃棄・再利用まで生産者の責任を広げるEPR。Dは税務の処理を会計士が見る説明で、環境監査の定義ではない。OOOXの1。

\Needspace{4\baselineskip}\paragraph{選択肢の検討（査読待ち初稿）}\mbox{}\par

\noindent\textbf{1：適切・選択対象} OOOX。判定と一致する。\par

\noindent\textbf{2：不適切} XXXO。式・定義・対応の独立判定と一致しない。\par

\noindent\textbf{3：不適切} XOOX。式・定義・対応の独立判定と一致しない。\par

\noindent\textbf{4：不適切} OOXO。式・定義・対応の独立判定と一致しない。\par

\noindent\textbf{5：不適切} OXXX。式・定義・対応の独立判定と一致しない。\par

\paragraph{再利用する考え方} ライフサイクル、汚染者負担原則、拡大生産者責任、監査。 集計対象・時点・法の主体・選択方向を固定し、同型でも元の五肢へ戻す。

\paragraph{共通理論} \hyperref[theory:env:principles]{第\ref*{theory:env:principles}節}

\par 参考：\cite{env59_epr}。

\par 難易度 2（校準初稿）、優先度未評価。

\section{評価リスクと知覚リスクの隔たり}

\subsection{令和元年度再試験 問5-6}
\label{q:r01-ret-basic-5-6}

\noindent\textbf{原典：} sources/R01\_retest/basic.pdf、PDF物理20頁。

\textbf{公式正答：} 3（single）。

\paragraph{問題の要旨} 評価リスクと知覚リスクの隔たり。

\Needspace{4\baselineskip}\paragraph{確認した条件・記述（独自要約）}\mbox{}\par

\noindent\textbf{時点：} 令和元年度再試験に指定された資料・制度。\par

\noindent\textbf{照合：} 原典18画像で否定・五肢・組合せ表を確認。後年値で置換しない。\par

\paragraph{問われる核心} 評価リスクと知覚リスクの隔たり。

\paragraph{必要知識} リスク評価・管理・対話、知覚、根拠と社会判断。

\paragraph{考え方・正答を選ぶ理由（査読待ち初稿）} 科学的評価と人の知覚は、経験・信頼・自発性・公平性等により隔たり得る。隔たりがないという3が誤り。相互の情報・意見交換によって理解し、評価の不確実性と管理の価値判断も説明する。レギュラトリーサイエンスは科学的根拠を政策等の判断へつなぐが、科学だけで価値判断を消すものではない。

\Needspace{4\baselineskip}\paragraph{選択肢の検討（査読待ち初稿）}\mbox{}\par

\noindent\textbf{1：条件内で適切} 科学評価と費用技術を含む管理。条件内で適切な記述であり、不適切な記述を選ぶ対象ではない。\par

\noindent\textbf{2：条件内で適切} 情報と意見の相互交換。条件内で適切な記述であり、不適切な記述を選ぶ対象ではない。\par

\noindent\textbf{3：不適切・選択対象} 評価と知覚に隔たりなし。判定と一致する。\par

\noindent\textbf{4：条件内で適切} 規制科学が信頼性と波及評価。条件内で適切な記述であり、不適切な記述を選ぶ対象ではない。\par

\noindent\textbf{5：条件内で適切} 規制と合意形成を支援。条件内で適切な記述であり、不適切な記述を選ぶ対象ではない。\par

\paragraph{再利用する考え方} リスク評価・管理・対話、知覚、根拠と社会判断。 集計対象・時点・法の主体・選択方向を固定し、同型でも元の五肢へ戻す。

\paragraph{共通理論} \hyperref[theory:env:risk]{第\ref*{theory:env:risk}節}

\par 参考：\cite{eth54_riskcomm}。

\par 難易度 2（校準初稿）、優先度未評価。

\subsection{令和4年度 問5-5}
\label{q:r04-reg-basic-5-5}

\noindent\textbf{原典：} sources/R04/basic.pdf、PDF物理20頁。

\textbf{公式正答：} 5（single）。

\paragraph{問題の要旨} リスクコミュニケーションでは評価過程を開示。

\Needspace{4\baselineskip}\paragraph{確認した条件・記述（独自要約）}\mbox{}\par

\noindent\textbf{原典：} 全五肢と否定方向・組合せ表を画像確認。\par

\noindent\textbf{資料時点：} 令和4年度の指定資料・統計年を保持し、後年値に置換しない。\par

\paragraph{問われる核心} リスクコミュニケーションでは評価過程を開示。

\paragraph{必要知識} 評価と管理、知覚と科学的評価、双方向と透明性。

\paragraph{考え方・正答を選ぶ理由（査読待ち初稿）} 5が不適切。評価値だけでなく仮定・不確実性・手順を示し、受け手との意見交換を通じて理解を形成する。管理上の費用や実行可能性を評価そのものと混同しない。

\Needspace{4\baselineskip}\paragraph{選択肢の検討（査読待ち初稿）}\mbox{}\par

\noindent\textbf{1：条件内で適切} 評価結果と費用・技術を考え管理。科学的評価と管理上の選択は異なる段階。\par

\noindent\textbf{2：条件内で適切} レギュラトリーサイエンスが合意形成を支援。評価を社会的判断へつなげる科学的基盤。\par

\noindent\textbf{3：条件内で適切} 個人・機関・集団の相互交換。一方向の通知だけとは異なる。\par

\noindent\textbf{4：条件内で適切} 評価リスクと認知リスクに隔たり。知識・価値観・信頼等により受け止めが変わる。\par

\noindent\textbf{5：不適切・選択対象} 混乱回避のため評価過程を開示しない。根拠・不確実性・評価過程を説明して対話する。混乱を理由に非開示を原則とはしない。\par

\paragraph{再利用する考え方} 評価と管理、知覚と科学的評価、双方向と透明性。 否定方向と資料年、分母と単位を確認して五肢へ戻る。

\paragraph{共通理論} \hyperref[theory:env:risk]{第\ref*{theory:env:risk}節}

\par 難易度 2（校準初稿）、優先度未評価。

\section{SDGsの対象は先進国も含む}

\subsection{平成30年度 問5-1}
\label{q:h30-reg-basic-5-1}

\noindent\textbf{原典：} sources/H30/basic.pdf、PDF物理18頁。

\textbf{公式正答：} 2（single）。

\paragraph{問題の要旨} SDGsの対象は先進国も含む。

\Needspace{4\baselineskip}\paragraph{確認した条件・記述（独自要約）}\mbox{}\par

\noindent\textbf{時点：} 平成30年度に指定された資料・制度。\par

\noindent\textbf{照合：} 原典18画像で否定・五肢・組合せ表を確認。後年値で置換しない。\par

\paragraph{問われる核心} SDGsの対象は先進国も含む。

\paragraph{必要知識} 普遍性、統合、ゴールとターゲット、実施体制。

\paragraph{考え方・正答を選ぶ理由（査読待ち初稿）} SDGsは普遍的で全ての国を対象とするので、途上国だけへ限定した2が誤り。17と169、環境・経済・社会の統合、各主体のパートナーシップを一組として読む。MDGsと同じ対象範囲へ戻さない。

\Needspace{4\baselineskip}\paragraph{選択肢の検討（査読待ち初稿）}\mbox{}\par

\noindent\textbf{1：条件内で適切} 2015の2030アジェンダ中核。条件内で適切な記述であり、不適切な記述を選ぶ対象ではない。\par

\noindent\textbf{2：不適切・選択対象} 途上国だけの目標。判定と一致する。\par

\noindent\textbf{3：条件内で適切} 17ゴール169ターゲット。条件内で適切な記述であり、不適切な記述を選ぶ対象ではない。\par

\noindent\textbf{4：条件内で適切} 多様主体の連携。条件内で適切な記述であり、不適切な記述を選ぶ対象ではない。\par

\noindent\textbf{5：条件内で適切} 日本2016実施指針。条件内で適切な記述であり、不適切な記述を選ぶ対象ではない。\par

\paragraph{再利用する考え方} 普遍性、統合、ゴールとターゲット、実施体制。 集計対象・時点・法の主体・選択方向を固定し、同型でも元の五肢へ戻す。

\paragraph{共通理論} \hyperref[theory:env:sdgs]{第\ref*{theory:env:sdgs}節}

\par 参考：\cite{env_sdgs_un,env_sdgs2016}。

\par 難易度 2（校準初稿）、優先度未評価。

\section{スマートメーターとHEMSの機能を分ける}

\subsection{平成30年度 問5-4}
\label{q:h30-reg-basic-5-4}

\noindent\textbf{原典：} sources/H30/basic.pdf、PDF物理19頁。

\textbf{公式正答：} 4（single）。

\paragraph{問題の要旨} スマートメーターとHEMSの機能を分ける。

\Needspace{4\baselineskip}\paragraph{確認した条件・記述（独自要約）}\mbox{}\par

\noindent\textbf{時点：} 平成30年度に指定された資料・制度。\par

\noindent\textbf{照合：} 原典18画像で否定・五肢・組合せ表を確認。後年値で置換しない。\par

\paragraph{問われる核心} スマートメーターとHEMSの機能を分ける。

\paragraph{必要知識} 計測、通信、家庭管理、地域と系統の範囲。

\paragraph{考え方・正答を選ぶ理由（査読待ち初稿）} 家庭内の機器を管理・制御する説明はHEMSに対応。スマートメーターは通信機能を持つ電力量計であり、同じ装置として家庭全体の最適制御を定義した4が不適切。計測情報は管理システムの入力に使えるが、計測と制御の主体を分ける。

\Needspace{4\baselineskip}\paragraph{選択肢の検討（査読待ち初稿）}\mbox{}\par

\noindent\textbf{1：条件内で適切} 系統需給のsmartgrid。条件内で適切な記述であり、不適切な記述を選ぶ対象ではない。\par

\noindent\textbf{2：条件内で適切} 地域の総合管理community。条件内で適切な記述であり、不適切な記述を選ぶ対象ではない。\par

\noindent\textbf{3：条件内で適切} 住宅機器を組み合わせるhouse。条件内で適切な記述であり、不適切な記述を選ぶ対象ではない。\par

\noindent\textbf{4：不適切・選択対象} メーターが家庭需給を最適制御。判定と一致する。\par

\noindent\textbf{5：条件内で適切} smartはかしこい。条件内で適切な記述であり、不適切な記述を選ぶ対象ではない。\par

\paragraph{再利用する考え方} 計測、通信、家庭管理、地域と系統の範囲。 集計対象・時点・法の主体・選択方向を固定し、同型でも元の五肢へ戻す。

\paragraph{共通理論} \hyperref[theory:env:smart]{第\ref*{theory:env:smart}節}

\par 参考：\cite{env59_smart}。

\par 難易度 2（校準初稿）、優先度未評価。

\section{応用が科学的発見より先行する例}

\subsection{令和元年度 問5-5}
\label{q:r01-reg-basic-5-5}

\noindent\textbf{原典：} sources/R01/basic.pdf、PDF物理21頁。

\textbf{公式正答：} 2（single）。

\par\textbf{検討未確定：} 他の四組の初出原報・開発年を全て取得したわけではない。土星の環は観測と環構造の解釈を分ける。

\paragraph{問題の要旨} 応用が科学的発見より先行する例。

\Needspace{4\baselineskip}\paragraph{確認した条件・記述（独自要約）}\mbox{}\par

\noindent\textbf{時点：} 令和元年度に指定された資料・制度。\par

\noindent\textbf{照合：} 原典18画像で否定・五肢・組合せ表を確認。後年値で置換しない。\par

\paragraph{問われる核心} 応用が科学的発見より先行する例。

\paragraph{必要知識} 経験的技術、機序解明、観測装置、時間順。

\paragraph{考え方・正答を選ぶ理由（査読待ち初稿）} Jennerの種痘は1796頃で、ウイルスの発見・概念形成は19世紀末。従って発見を種痘の前とした2が逆。実用的な方法が機序の科学的解明より先行する例であり、科学→技術の一方向の順序だけでは歴史を説明できない。

\Needspace{4\baselineskip}\paragraph{選択肢の検討（査読待ち初稿）}\mbox{}\par

\noindent\textbf{1：条件内で適切} 核分裂後に原子力発電。条件内で適切な記述であり、不適切な記述を選ぶ対象ではない。\par

\noindent\textbf{2：不適切・選択対象} ウイルス発見後に種痘。判定と一致する。\par

\noindent\textbf{3：条件内で適切} 望遠鏡後に土星環確認。条件内で適切な記述であり、不適切な記述を選ぶ対象ではない。\par

\noindent\textbf{4：条件内で適切} 量子力学後にトランジスタ。条件内で適切な記述であり、不適切な記述を選ぶ対象ではない。\par

\noindent\textbf{5：条件内で適切} 電磁波確認後にレーダー。条件内で適切な記述であり、不適切な記述を選ぶ対象ではない。\par

\paragraph{再利用する考え方} 経験的技術、機序解明、観測装置、時間順。 集計対象・時点・法の主体・選択方向を固定し、同型でも元の五肢へ戻す。

\paragraph{共通理論} \hyperref[theory:env:technologysequence]{第\ref*{theory:env:technologysequence}節}

\par 参考：\cite{env59_jenner}。

\par 難易度 3（校準初稿）、優先度未評価。

\subsection{令和5年度 問5-6}
\label{q:r05-reg-basic-5-6}

\noindent\textbf{原典：} sources/R05/basic.pdf、PDF物理18頁。

\textbf{公式正答：} 5（single）。

\par\textbf{検討未確定：} 土星環の解釈・レーダー・発電開始の基準年は各一次史料へ追加照合。

\paragraph{問題の要旨} ワクチンの実践はウイルス発見に先行。

\Needspace{4\baselineskip}\paragraph{確認した条件・記述（独自要約）}\mbox{}\par

\noindent\textbf{原典：} 全五肢と否定方向・組合せ表を画像確認。\par

\noindent\textbf{資料時点：} 令和5年度の指定資料・統計年を保持し、後年値に置換しない。\par

\paragraph{問われる核心} ワクチンの実践はウイルス発見に先行。

\paragraph{必要知識} 観測・理論・応用の双方向性、発見と実践の順序。

\paragraph{考え方・正答を選ぶ理由（査読待ち初稿）} 5が不適切。科学が技術へ一方向に進むとは限らず、種痘の経験的実践はウイルス概念の成立より早い。望遠鏡の発明と環の正しい解釈の時期も区別する。

\Needspace{4\baselineskip}\paragraph{選択肢の検討（査読待ち初稿）}\mbox{}\par

\noindent\textbf{1：条件内で適切} 望遠鏡後に土星の環確認。観測手段が天体の新しい知識を生む例。\par

\noindent\textbf{2：条件内で適切} 量子力学後にトランジスター。20世紀前半の理論と1947の素子発明の順。\par

\noindent\textbf{3：条件内で適切} 電磁波確認後にレーダー。Hertzの確認とその後の電波応用の順。\par

\noindent\textbf{4：条件内で適切} 核分裂発見後に原子力発電。1938核分裂の後に発電利用へ進む。\par

\noindent\textbf{5：不適切・選択対象} ウイルス発見後にワクチン接種開始。Jennerの種痘1796はウイルス研究の19世紀末より先。病原体の詳細理解前にも技術が発達した。\par

\paragraph{再利用する考え方} 観測・理論・応用の双方向性、発見と実践の順序。 否定方向と資料年、分母と単位を確認して五肢へ戻る。

\paragraph{共通理論} \hyperref[theory:env:technologysequence]{第\ref*{theory:env:technologysequence}節}

\par 参考：\cite{env59_jenner}。

\par 難易度 3（校準初稿）、優先度未評価。

\section{AR6観測温暖化は約1.1℃であって3℃ではない}

\subsection{令和4年度 問5-1}
\label{q:r04-reg-basic-5-1}

\noindent\textbf{原典：} sources/R04/basic.pdf、PDF物理17頁。

\textbf{公式正答：} 2（single）。

\par\textbf{検討未確定：} WGI温度箇所を確認。WGII脆弱性とWGIII排出量の各原SPM全段落照合は未完。

\paragraph{問題の要旨} AR6観測温暖化は約1.1℃であって3℃ではない。

\Needspace{4\baselineskip}\paragraph{確認した条件・記述（独自要約）}\mbox{}\par

\noindent\textbf{原典：} 全五肢と否定方向・組合せ表を画像確認。\par

\noindent\textbf{資料時点：} 令和4年度の指定資料・統計年を保持し、後年値に置換しない。\par

\paragraph{問われる核心} AR6観測温暖化は約1.1℃であって3℃ではない。

\paragraph{必要知識} 期間平均の観測差と将来予測、気温と排出量。

\paragraph{考え方・正答を選ぶ理由（査読待ち初稿）} 2が不適切。WGIの2011〜2020対1850〜1900は1.09℃（評価幅0.95〜1.20℃）であり、将来シナリオの温度上昇や特定地域の上昇と区別する。人間影響、地域別脆弱性、排出量は各作業部会の対象を分ける。

\Needspace{4\baselineskip}\paragraph{選択肢の検討（査読待ち初稿）}\mbox{}\par

\noindent\textbf{1：条件内で適切} 人間影響による温暖化は疑う余地がない。WGIの結論は人間影響について明確な評価を示す。\par

\noindent\textbf{2：不適切・選択対象} 2011〜2020は1850〜1900より約3℃高い。指定二期間の世界平均地表温度差は1.09℃、概数1.1℃。約3℃は観測済みの差ではない。\par

\noindent\textbf{3：条件内で適切} 極端現象や酸性化・海面上昇等の影響。物理・生態系影響を一種類に限定しない。\par

\noindent\textbf{4：条件内で適切} 脆弱性は地域内外で異なる。発展形態や曝露、適応能力が脆弱性を変える。\par

\noindent\textbf{5：条件内で適切} 2010〜2019の平均排出量は過去の十年より高い。WGIIIの評価対象は年間平均の正味人為GHG排出量。濃度や年ごとの増加率とは違う。\par

\paragraph{再利用する考え方} 期間平均の観測差と将来予測、気温と排出量。 否定方向と資料年、分母と単位を確認して五肢へ戻る。

\paragraph{共通理論} \hyperref[theory:env:ar6]{第\ref*{theory:env:ar6}節}

\par 参考：\cite{env60_ar6}。

\par 難易度 3（校準初稿）、優先度未評価。

\section{里地里山の危機は利用縮小・管理低下}

\subsection{令和5年度 問5-1}
\label{q:r05-reg-basic-5-1}

\noindent\textbf{原典：} sources/R05/basic.pdf、PDF物理16頁。

\textbf{公式正答：} 3（single）。

\par\textbf{検討未確定：} 指定国家戦略の固有種比率・開発圧力・分布事例の全原文は追加確認対象。

\paragraph{問題の要旨} 里地里山の危機は利用縮小・管理低下。

\Needspace{4\baselineskip}\paragraph{確認した条件・記述（独自要約）}\mbox{}\par

\noindent\textbf{原典：} 全五肢と否定方向・組合せ表を画像確認。\par

\noindent\textbf{資料時点：} 令和5年度の指定資料・統計年を保持し、後年値に置換しない。\par

\paragraph{問われる核心} 里地里山の危機は利用縮小・管理低下。

\paragraph{必要知識} 生物多様性の危機、固有種、二次的自然と人の管理。

\paragraph{考え方・正答を選ぶ理由（査読待ち初稿）} 3が不適切。里山は人間の利用で維持されてきた生態系も含むため、手入れや利用の停止が二次的自然を変える。開発による第一の危機と利用縮小による第二の危機を分け、外来種・化学物質・気候変動の要因も別に読む。

\Needspace{4\baselineskip}\paragraph{選択肢の検討（査読待ち初稿）}\mbox{}\par

\noindent\textbf{1：条件内で適切} 爬虫類約6割・両生類約8割が固有種。固有種割合を全動植物へ一般化しない。\par

\noindent\textbf{2：条件内で適切} 近年大規模開発の圧力は低下。過去の減少と近年の開発圧力の変化を別に記す。\par

\noindent\textbf{3：不適切・選択対象} 里山の利用拡大で生息地減少。里地里山は利用・管理の縮小による環境変化が主要な危機。拡大という原因説明が逆。\par

\noindent\textbf{4：条件内で適切} 他地域からの導入種が在来へ影響。国外だけでなく国内他地域からの移動でも影響し得る。\par

\noindent\textbf{5：条件内で適切} タケや南方系チョウの北上。気候変化と分布の変化の事例。全種が一様に北上するという主張ではない。\par

\paragraph{再利用する考え方} 生物多様性の危機、固有種、二次的自然と人の管理。 否定方向と資料年、分母と単位を確認して五肢へ戻る。

\paragraph{共通理論} \hyperref[theory:env:biodiversity2023]{第\ref*{theory:env:biodiversity2023}節}

\par 参考：\cite{env60_biodiversity}。

\par 難易度 3（校準初稿）、優先度未評価。

\section{2024には藻場吸収量を国連へ報告済み}

\subsection{令和7年度 問5-1}
\label{q:r07-reg-basic-5-1}

\noindent\textbf{原典：} sources/R07/basic.pdf、PDF物理15頁。

\textbf{公式正答：} 5（single）。

\par\textbf{検討未確定：} 報告済みの藻場範囲は確認。他4肢の当時法令要件とZEH・ZEB目標全文の照合は追加対象。

\paragraph{問題の要旨} 2024には藻場吸収量を国連へ報告済み。

\Needspace{4\baselineskip}\paragraph{確認した条件・記述（独自要約）}\mbox{}\par

\noindent\textbf{原典：} 全五肢と否定方向・組合せ表を画像確認。\par

\noindent\textbf{資料時点：} 令和7年度の指定資料・統計年を保持し、後年値に置換しない。\par

\paragraph{問われる核心} 2024には藻場吸収量を国連へ報告済み。

\paragraph{必要知識} ブルーカーボン、評価年と提出年、濃度と排出吸収。

\paragraph{考え方・正答を選ぶ理由（査読待ち初稿）} 5が不適切。令和6年白書の時点では、2022年度排出・吸収量の2024年提出インベントリに海草・海藻藻場の評価が含まれた。海の全炭素吸収を漏れなく算定できたという意味ではない。濃度観測・事業者排出報告・吸収インベントリを別に整理する。

\Needspace{4\baselineskip}\paragraph{選択肢の検討（査読待ち初稿）}\mbox{}\par

\noindent\textbf{1：条件内で適切} 衛星で全球GHG濃度を継続観測。衛星は広域の濃度観測を行う。排出量を各地点で直接計量するという説明ではない。\par

\noindent\textbf{2：条件内で適切} 一定量以上排出する事業者の年度報告。算定・報告・公表制度は排出量の把握と公表を扱う。\par

\noindent\textbf{3：条件内で適切} HFCにも強い温室効果があり回収処理。オゾン層を破壊しないことと温室効果が小さいことは同義でない。\par

\noindent\textbf{4：条件内で適切} 2030以降新築にZEH・ZEB水準を目指す。省エネルギー性能の水準と、全建物が毎時ゼロであることは違う。\par

\noindent\textbf{5：不適切・選択対象} ブルーカーボン吸収を国連に未報告。2024年4月に海草・海藻藻場の吸収量を合わせて国連へ報告したため未報告は不適切。\par

\paragraph{再利用する考え方} ブルーカーボン、評価年と提出年、濃度と排出吸収。 否定方向と資料年、分母と単位を確認して五肢へ戻る。

\paragraph{共通理論} \hyperref[theory:env:blue]{第\ref*{theory:env:blue}節}

\par 参考：\cite{env60_blue}。

\par 難易度 3（校準初稿）、優先度未評価。

\section{代替フロンのオゾン保護と温暖化対策を区別}

\subsection{令和3年度 問5-1}
\label{q:r03-reg-basic-5-1}

\noindent\textbf{原典：} sources/R03/basic.pdf、PDF物理21頁。

\textbf{公式正答：} 2（single）。

\par\textbf{検討未確定：} TCFDの2020賛同200以上の数、FSB構成、各制度の当時の定義全文は追加照合対象。

\paragraph{問題の要旨} 代替フロンのオゾン保護と温暖化対策を区別。

\Needspace{4\baselineskip}\paragraph{確認した条件・記述（独自要約）}\mbox{}\par

\noindent\textbf{原典：} 全五肢と否定方向・組合せ表を画像確認。\par

\noindent\textbf{資料時点：} 令和3年度の指定資料・統計年を保持し、後年値に置換しない。\par

\paragraph{問われる核心} 代替フロンのオゾン保護と温暖化対策を区別。

\paragraph{必要知識} オゾン層破壊係数と温暖化係数、電力と一次エネルギー、目標と開示。

\paragraph{考え方・正答を選ぶ理由（査読待ち初稿）} 2が不適切。CFC・HCFCのオゾン層破壊とHFCの地球温暖化係数は別の問題である。RE100、TCFD、自治体目標、ZEHはそれぞれ電力調達・情報開示・地域目標・建築性能という異なる対象に働く。

\Needspace{4\baselineskip}\paragraph{選択肢の検討（査読待ち初稿）}\mbox{}\par

\noindent\textbf{1：条件内で適切} RE100は使用電力100\%再エネを目指す。目標は企業活動に使用する電力の再エネ化。全エネルギーを必ず電力だけで賄うという定義ではない。\par

\noindent\textbf{2：不適切・選択対象} 代替フロン転換で温暖化対策も十分。HFC等にはオゾン層を破壊しなくても強い温室効果を持つものがあり、転換だけで十分とはいえない。\par

\noindent\textbf{3：条件内で適切} TCFDは投資判断に資する開示を促す。気候関連のリスク・機会を財務情報として示す取組。賛同機関数は2020時点の値として扱う。\par

\noindent\textbf{4：条件内で適切} 自治体の2050実質ゼロ表明。ゼロカーボンシティは自治体の実質ゼロ目標表明を示す。\par

\noindent\textbf{5：条件内で適切} ZEH・ZEH-Mは年間一次エネルギー収支ゼロを目指す。断熱・高効率設備で需要を下げ再エネを導入する。設備導入だけ、又は瞬間ごとのゼロとは限らない。\par

\paragraph{再利用する考え方} オゾン層破壊係数と温暖化係数、電力と一次エネルギー、目標と開示。 否定方向と資料年、分母と単位を確認して五肢へ戻る。

\paragraph{共通理論} \hyperref[theory:env:climateaction]{第\ref*{theory:env:climateaction}節}

\par 参考：\cite{env60_blue}。

\par 難易度 3（校準初稿）、優先度未評価。

\section{太陽電池の国内生産品比率は毎年拡大ではない}

\subsection{令和5年度 問5-3}
\label{q:r05-reg-basic-5-3}

\noindent\textbf{原典：} sources/R05/basic.pdf、PDF物理17頁。

\textbf{公式正答：} 1（single）。

\par\textbf{検討未確定：} 2009〜2020国内生産比率の原表、当時の原油・LNG・風力・地熱数値を追加照合。

\paragraph{問題の要旨} 太陽電池の国内生産品比率は毎年拡大ではない。

\Needspace{4\baselineskip}\paragraph{確認した条件・記述（独自要約）}\mbox{}\par

\noindent\textbf{原典：} 全五肢と否定方向・組合せ表を画像確認。\par

\noindent\textbf{資料時点：} 令和5年度の指定資料・統計年を保持し、後年値に置換しない。\par

\paragraph{問われる核心} 太陽電池の国内生産品比率は毎年拡大ではない。

\paragraph{必要知識} 設備導入と製造比率、資源量と導入、原油とLNG。

\paragraph{考え方・正答を選ぶ理由（査読待ち初稿）} 1が不適切。太陽光発電の国内導入が増えても、国内出荷に占める国内生産品比率が同時に増えるとは限らない。原油・LNG・風力・地熱は輸入構成、設置容量、資源量を別々に評価する。

\Needspace{4\baselineskip}\paragraph{選択肢の検討（査読待ち初稿）}\mbox{}\par

\noindent\textbf{1：不適切・選択対象} 導入量も国内生産品比率も2009〜2020毎年拡大。導入設備の増加と国内製品比率は独立。輸入増加に伴う国内比率低下を落とし、毎年拡大とするのが不適切。\par

\noindent\textbf{2：条件内で適切} 2020原油中東90\%以上でサウジ・UAE上位。原油の数量ベースの依存を読む。\par

\noindent\textbf{3：条件内で適切} LNGの80\%以上が中東以外で豪州最大。原油とは異なる調達構成。LNGを原油と同じ中東依存としない。\par

\noindent\textbf{4：条件内で適切} 2020末風力400万kW超と系統拡大。導入容量の単位kWと発電量kWhは違う。\par

\noindent\textbf{5：条件内で適切} 地熱資源量は日本世界3位。資源量順位と実際の導入容量順位は別。\par

\paragraph{再利用する考え方} 設備導入と製造比率、資源量と導入、原油とLNG。 否定方向と資料年、分母と単位を確認して五肢へ戻る。

\paragraph{共通理論} \hyperref[theory:env:energy2020]{第\ref*{theory:env:energy2020}節}

\par 難易度 3（校準初稿）、優先度未評価。

\section{第7次計画は技術の選択肢を限定しない}

\subsection{令和7年度 問5-3}
\label{q:r07-reg-basic-5-3}

\noindent\textbf{原典：} sources/R07/basic.pdf、PDF物理17頁。

\textbf{公式正答：} 4（single）。

\par\textbf{検討未確定：} 2のS省略は不完全な要約として注意。計画本文と目標の各原段落全体を査読し、単に安全性を不要とする文へ読み替えない。

\paragraph{問題の要旨} 第7次計画は技術の選択肢を限定しない。

\Needspace{4\baselineskip}\paragraph{確認した条件・記述（独自要約）}\mbox{}\par

\noindent\textbf{原典：} 全五肢と否定方向・組合せ表を画像確認。\par

\noindent\textbf{資料時点：} 令和7年度の指定資料・統計年を保持し、後年値に置換しない。\par

\paragraph{問われる核心} 第7次計画は技術の選択肢を限定しない。

\paragraph{必要知識} S+3E、技術選択肢、需要予測、基準年と削減目標。

\paragraph{考え方・正答を選ぶ理由（査読待ち初稿）} 主判断4。脱炭素へ向け使える技術を幅広く活用する方針を、あらかじめ選択肢を限定する方針へ反転している。2は安全性Sが大前提であることを補って読む。電力需要の予測、削減目標、政策の価値基準は別の論点である。

\Needspace{4\baselineskip}\paragraph{選択肢の検討（査読待ち初稿）}\mbox{}\par

\noindent\textbf{1：条件内で適切} データセンター等で電力需要増の可能性。効率改善を見込んでも需要を押し上げる要因がある。\par

\noindent\textbf{2：条件内で適切} 三つのEの最適バランスが基本的視点。安全性Sを大前提に3Eのバランスを追求する。原典の三つのEだけの短縮を完全定義としない。\par

\noindent\textbf{3：条件内で適切} 2035/2040に2013比60\%/73\%削減。目標の基準年は2013年度。前年度比ではない。\par

\noindent\textbf{4：不適切・選択対象} 利用技術を選別し選択肢を限定。不確実性に対応し使える技術の選択肢を幅広く追求する方向と逆で不適切。\par

\noindent\textbf{5：条件内で適切} 脱炭素投資を経済成長へつなげる。供給構造転換と産業投資・成長を一体で扱う。\par

\paragraph{再利用する考え方} S+3E、技術選択肢、需要予測、基準年と削減目標。 否定方向と資料年、分母と単位を確認して五肢へ戻る。

\paragraph{共通理論} \hyperref[theory:env:energyplan7]{第\ref*{theory:env:energyplan7}節}

\par 参考：\cite{env60_plan7}。

\par 難易度 3（校準初稿）、優先度未評価。

\section{2022家庭CO2の積み上げ棒は電気・都市ガス・LP・灯油}

\subsection{令和6年度 問5-3}
\label{q:r06-reg-basic-5-3}

\noindent\textbf{原典：} sources/R06/basic.pdf、PDF物理19頁。

\textbf{公式正答：} 2（single）。

\paragraph{問題の要旨} 2022家庭CO2の積み上げ棒は電気・都市ガス・LP・灯油。

\Needspace{4\baselineskip}\paragraph{確認した条件・記述（独自要約）}\mbox{}\par

\noindent\textbf{指定：} 令和4年度家庭CO2確報、令和6年3月。\par

\noindent\textbf{図：} 2017〜2022年度、A〜Dは下からの順。\par

\paragraph{問われる核心} 2022家庭CO2の積み上げ棒は電気・都市ガス・LP・灯油。

\paragraph{必要知識} 積み上げ棒の凡例、排出係数、世帯平均と燃料量。

\paragraph{考え方・正答を選ぶ理由（査読待ち初稿）} 原典R4棒のA=\allowbreak{}1.74、B=.38、C=.15、D=.32（t-CO2/世帯・年）。確報の区分へ対応させA電気、B都市ガス、C LPガス、D灯油で2。合計2.59。2017等の棒の値を2022へ置換しない。電気は発電段階の排出を係数で世帯へ帰属させ、各家庭が煙突から直接排出した量だけではない。

\Needspace{4\baselineskip}\paragraph{選択肢の検討（査読待ち初稿）}\mbox{}\par

\noindent\textbf{1：不適切} 電気/都市ガス/灯油/LP。独立に求めた組合せ 電気/都市ガス/LP/灯油 と異なる。\par

\noindent\textbf{2：適切・選択対象} 電気/都市ガス/LP/灯油。独立に求めた組合せ 電気/都市ガス/LP/灯油 と一致する。\par

\noindent\textbf{3：不適切} 都市ガス/電気/LP/灯油。独立に求めた組合せ 電気/都市ガス/LP/灯油 と異なる。\par

\noindent\textbf{4：不適切} 都市ガス/電気/灯油/LP。独立に求めた組合せ 電気/都市ガス/LP/灯油 と異なる。\par

\noindent\textbf{5：不適切} 灯油/電気/都市ガス/LP。独立に求めた組合せ 電気/都市ガス/LP/灯油 と異なる。\par

\paragraph{再利用する考え方} 積み上げ棒の凡例、排出係数、世帯平均と燃料量。 否定方向と資料年、分母と単位を確認して五肢へ戻る。

\paragraph{共通理論} \hyperref[theory:env:house2022]{第\ref*{theory:env:house2022}節}

\par 参考：\cite{env60_house2022}。

\par 難易度 2（校準初稿）、優先度未評価。

\section{液体水素の温度・体積・質量発熱と還元}

\subsection{令和4年度 問5-4}
\label{q:r04-reg-basic-5-4}

\noindent\textbf{原典：} sources/R04/basic.pdf、PDF物理19頁。

\textbf{公式正答：} 5（single）。

\par\textbf{検討未確定：} 鉄鉱石全般について温度・相を省いた吸熱断定を避ける。液化値と発熱量の一次物性表は追加照合。

\paragraph{問題の要旨} 液体水素の温度・体積・質量発熱と還元。

\Needspace{4\baselineskip}\paragraph{確認した条件・記述（独自要約）}\mbox{}\par

\noindent\textbf{原典：} 全五肢と否定方向・組合せ表を画像確認。\par

\noindent\textbf{資料時点：} 令和4年度の指定資料・統計年を保持し、後年値に置換しない。\par

\paragraph{問われる核心} 液体水素の温度・体積・質量発熱と還元。

\paragraph{必要知識} 沸点、気液密度、質量と体積の発熱量、反応熱の成立条件。

\paragraph{考え方・正答を選ぶ理由（査読待ち初稿）} 1気圧で水素の液化温度は約−253℃、常温気体と比べ液体体積は約1/800。ガソリンより大きいのは単位質量当たりの発熱量で、単位体積ではない。水素を用いた鉄鉱石還元の工程は熱の供給を要する吸熱側の技術として扱い5。具体的反応エンタルピーは酸化物の種類・温度・生成水の相を指定して求める。水素製造まで含む排出が必ずゼロとの意味ではない。

\Needspace{4\baselineskip}\paragraph{選択肢の検討（査読待ち初稿）}\mbox{}\par

\noindent\textbf{1：不適切} -162/1:600/重量/吸熱。独立に求めた組合せ -253/1:800/重量/吸熱 と異なる。\par

\noindent\textbf{2：不適切} -162/1:800/重量/発熱。独立に求めた組合せ -253/1:800/重量/吸熱 と異なる。\par

\noindent\textbf{3：不適切} -253/1:600/体積/発熱。独立に求めた組合せ -253/1:800/重量/吸熱 と異なる。\par

\noindent\textbf{4：不適切} -253/1:800/体積/発熱。独立に求めた組合せ -253/1:800/重量/吸熱 と異なる。\par

\noindent\textbf{5：適切・選択対象} -253/1:800/重量/吸熱。独立に求めた組合せ -253/1:800/重量/吸熱 と一致する。\par

\paragraph{再利用する考え方} 沸点、気液密度、質量と体積の発熱量、反応熱の成立条件。 否定方向と資料年、分母と単位を確認して五肢へ戻る。

\paragraph{共通理論} \hyperref[theory:env:hydrogen]{第\ref*{theory:env:hydrogen}節}

\par 参考：\cite{env60_hydrogen}。

\par 難易度 3（校準初稿）、優先度未評価。

\section{1949年の最初の日本人ノーベル賞は湯川}

\subsection{令和7年度 問5-5}
\label{q:r07-reg-basic-5-5}

\noindent\textbf{原典：} sources/R07/basic.pdf、PDF物理18頁。

\textbf{公式正答：} 2（single）。

\par\textbf{検討未確定：} 特許庁の出願日を索引確認。Nobelの各受賞一次ページ、ウィーン公式参加、Ford工程資料は追加照合。

\paragraph{問題の要旨} 1949年の最初の日本人ノーベル賞は湯川。

\Needspace{4\baselineskip}\paragraph{確認した条件・記述（独自要約）}\mbox{}\par

\noindent\textbf{原典：} 全五肢と否定方向・組合せ表を画像確認。\par

\noindent\textbf{資料時点：} 令和7年度の指定資料・統計年を保持し、後年値に置換しない。\par

\paragraph{問われる核心} 1949年の最初の日本人ノーベル賞は湯川。

\paragraph{必要知識} 受賞者・年・業績、政府公式参加、出願と登録。

\paragraph{考え方・正答を選ぶ理由（査読待ち初稿）} 2が不適切。人物・受賞年・対象業績が一組でなければならず、1949/中間子は湯川、1965/量子電気力学は朝永。4は出願日の記述であり特許の登録日を問うものではない。

\Needspace{4\baselineskip}\paragraph{選択肢の検討（査読待ち初稿）}\mbox{}\par

\noindent\textbf{1：条件内で適切} Haberはアンモニア合成で化学賞。空中窒素からのアンモニア合成に対する1918化学賞。\par

\noindent\textbf{2：不適切・選択対象} 朝永が1949に中間子予言で最初の物理賞。1949の中間子予言による賞は湯川秀樹。朝永振一郎は1965に量子電気力学で受賞。\par

\noindent\textbf{3：条件内で適切} 明治政府初の公式参加は1873ウィーン。政府公式参加と民間・幕府の先行参加を区別。\par

\noindent\textbf{4：条件内で適切} 堀田の錆止塗料は1885年7月1日出願。特許庁の制度史は出願日7月1日を示す。登録日と入れ替えない。\par

\noindent\textbf{5：条件内で適切} Fordの流れ作業で大量低コスト生産。工程を連続化し作業を分担する生産方式。\par

\paragraph{再利用する考え方} 受賞者・年・業績、政府公式参加、出願と登録。 否定方向と資料年、分母と単位を確認して五肢へ戻る。

\paragraph{共通理論} \hyperref[theory:env:nobeljapan]{第\ref*{theory:env:nobeljapan}節}

\par 参考：\cite{env60_patent,env60_yukawa}。

\par 難易度 3（校準初稿）、優先度未評価。

\section{2019年度原油依存と輸入金額の分母}

\subsection{令和4年度 問5-3}
\label{q:r04-reg-basic-5-3}

\noindent\textbf{原典：} sources/R04/basic.pdf、PDF物理18頁。

\textbf{公式正答：} 1（single）。

\par\textbf{検討未確定：} エネルギー白書2021の2019年度元表・輸入額比の原図は追加確認対象。

\paragraph{問題の要旨} 2019年度原油依存と輸入金額の分母。

\Needspace{4\baselineskip}\paragraph{確認した条件・記述（独自要約）}\mbox{}\par

\noindent\textbf{原典：} 全五肢と否定方向・組合せ表を画像確認。\par

\noindent\textbf{資料時点：} 令和4年度の指定資料・統計年を保持し、後年値に置換しない。\par

\paragraph{問われる核心} 2019年度原油依存と輸入金額の分母。

\paragraph{必要知識} 年度と白書年、数量依存度と金額比、輸送経路。

\paragraph{考え方・正答を選ぶ理由（査読待ち初稿）} 指定白書2021では2019年度の原油中東依存度を概数90\%、経由する海峡をホルムズ、最大相手国をサウジアラビア、石油・石油製品輸入額の総輸入額比を概数10\%として1。数量に占める地域割合と金額に占める品目割合は分母も単位も違う。

\Needspace{4\baselineskip}\paragraph{選択肢の検討（査読待ち初稿）}\mbox{}\par

\noindent\textbf{1：適切・選択対象} 90/ホルムズ/サウジ/10。独立に求めた組合せ 90/ホルムズ/サウジ/10 と一致する。\par

\noindent\textbf{2：不適切} 90/マラッカ/クウェート/32。独立に求めた組合せ 90/ホルムズ/サウジ/10 と異なる。\par

\noindent\textbf{3：不適切} 90/ホルムズ/クウェート/10。独立に求めた組合せ 90/ホルムズ/サウジ/10 と異なる。\par

\noindent\textbf{4：不適切} 67/マラッカ/クウェート/10。独立に求めた組合せ 90/ホルムズ/サウジ/10 と異なる。\par

\noindent\textbf{5：不適切} 67/ホルムズ/サウジ/32。独立に求めた組合せ 90/ホルムズ/サウジ/10 と異なる。\par

\paragraph{再利用する考え方} 年度と白書年、数量依存度と金額比、輸送経路。 否定方向と資料年、分母と単位を確認して五肢へ戻る。

\paragraph{共通理論} \hyperref[theory:env:oil2019]{第\ref*{theory:env:oil2019}節}

\par 難易度 3（校準初稿）、優先度未評価。

\section{2018発電比・2020太陽光費用と風速三乗}

\subsection{令和3年度 問5-3}
\label{q:r03-reg-basic-5-3}

\noindent\textbf{原典：} sources/R03/basic.pdf、PDF物理23頁。

\textbf{公式正答：} 1（single）。

\par\textbf{検討未確定：} 当時の発電統計と調達価格委員会の2020設置単純平均の原表を追加照合。全規模毎年低下という表現は個別年の反例を排除しない。

\paragraph{問題の要旨} 2018発電比・2020太陽光費用と風速三乗。

\Needspace{4\baselineskip}\paragraph{確認した条件・記述（独自要約）}\mbox{}\par

\noindent\textbf{原典：} 全五肢と否定方向・組合せ表を画像確認。\par

\noindent\textbf{資料時点：} 令和3年度の指定資料・統計年を保持し、後年値に置換しない。\par

\paragraph{問われる核心} 2018発電比・2020太陽光費用と風速三乗。

\paragraph{必要知識} 発電割合の水力境界、設備費の単純平均、風のエネルギー流束。

\paragraph{考え方・正答を選ぶ理由（査読待ち初稿）} 水力を除く2018年度再エネ発電割合は約9\%、2020設置の10kW以上太陽光システム費用の単純平均は約25万円/kWとして読む。風の運動エネルギー流束は、通過質量流量ρAvに単位質量エネルギーv\textsuperscript{2}/2を掛けるのでρAv\textsuperscript{3}/2。9/25/3の1。静止した一定質量の運動エネルギーv\textsuperscript{2}と混同しない。

\Needspace{4\baselineskip}\paragraph{選択肢の検討（査読待ち初稿）}\mbox{}\par

\noindent\textbf{1：適切・選択対象} 9/25/3。独立に求めた組合せ 9/25/3 と一致する。\par

\noindent\textbf{2：不適切} 14/25/3。独立に求めた組合せ 9/25/3 と異なる。\par

\noindent\textbf{3：不適切} 14/15/3。独立に求めた組合せ 9/25/3 と異なる。\par

\noindent\textbf{4：不適切} 9/25/2。独立に求めた組合せ 9/25/3 と異なる。\par

\noindent\textbf{5：不適切} 14/15/2。独立に求めた組合せ 9/25/3 と異なる。\par

\paragraph{再利用する考え方} 発電割合の水力境界、設備費の単純平均、風のエネルギー流束。 否定方向と資料年、分母と単位を確認して五肢へ戻る。

\paragraph{共通理論} \hyperref[theory:env:renewables]{第\ref*{theory:env:renewables}節}

\par 難易度 3（校準初稿）、優先度未評価。

\subsection{令和7年度 問5-4}
\label{q:r07-reg-basic-5-4}

\noindent\textbf{原典：} sources/R07/basic.pdf、PDF物理18頁。

\textbf{公式正答：} 2（single）。

\par\textbf{検討未確定：} 22.9\%は官庁発表で確認。2024設置10kW以上単純平均の原表は追加照合。価格算定用想定値や住宅用平均を代入しない。全規模毎年低下の一般化にも留保。

\paragraph{問題の要旨} 2023再エネ比・2024設置費と風速三乗。

\Needspace{4\baselineskip}\paragraph{確認した条件・記述（独自要約）}\mbox{}\par

\noindent\textbf{原典：} 全五肢と否定方向・組合せ表を画像確認。\par

\noindent\textbf{資料時点：} 令和7年度の指定資料・統計年を保持し、後年値に置換しない。\par

\paragraph{問われる核心} 2023再エネ比・2024設置費と風速三乗。

\paragraph{必要知識} 水力含有、対象年、単純平均と想定費用、流束の三乗。

\paragraph{考え方・正答を選ぶ理由（査読待ち初稿）} 水力を含む2023年度の再エネ発電割合22.9\%を約23\%、指定2024設置10kW以上の単純平均設備費を約23万円/kW、風のエネルギー流束をv\textsuperscript{3}とし2。R03の水力除外・2018年度約9\%と同じ値を使わない。風速二乗は一定質量の運動エネルギーであり、単位時間の通過量がさらにvに比例する。

\Needspace{4\baselineskip}\paragraph{選択肢の検討（査読待ち初稿）}\mbox{}\par

\noindent\textbf{1：不適切} 23/23/2。独立に求めた組合せ 23/23/3 と異なる。\par

\noindent\textbf{2：適切・選択対象} 23/23/3。独立に求めた組合せ 23/23/3 と一致する。\par

\noindent\textbf{3：不適切} 23/30/3。独立に求めた組合せ 23/23/3 と異なる。\par

\noindent\textbf{4：不適切} 15/30/2。独立に求めた組合せ 23/23/3 と異なる。\par

\noindent\textbf{5：不適切} 15/23/3。独立に求めた組合せ 23/23/3 と異なる。\par

\paragraph{再利用する考え方} 水力含有、対象年、単純平均と想定費用、流束の三乗。 否定方向と資料年、分母と単位を確認して五肢へ戻る。

\paragraph{共通理論} \hyperref[theory:env:renewables]{第\ref*{theory:env:renewables}節}

\par 参考：\cite{env60_re2023,env60_pv2024}。

\par 難易度 3（校準初稿）、優先度未評価。

\section{安全・標準化・製品責任の制定順}

\subsection{令和5年度 問5-5}
\label{q:r05-reg-basic-5-5}

\noindent\textbf{原典：} sources/R05/basic.pdf、PDF物理18頁。

\textbf{公式正答：} 1（single）。

\par\textbf{検討未確定：} 安全週間1928と工場法1911は官庁史料索引で確認。他3法の当初公布本文との照合は未完。

\paragraph{問題の要旨} 安全・標準化・製品責任の制定順。

\Needspace{4\baselineskip}\paragraph{確認した条件・記述（独自要約）}\mbox{}\par

\noindent\textbf{原典：} 全五肢と否定方向・組合せ表を画像確認。\par

\noindent\textbf{資料時点：} 令和5年度の指定資料・統計年を保持し、後年値に置換しない。\par

\paragraph{問われる核心} 安全・標準化・製品責任の制定順。

\paragraph{必要知識} 制定・施行・初回実施、改称と制定を区別。

\paragraph{考え方・正答を選ぶ理由（査読待ち初稿）} ウ工場法制定1911、エ第1回全国安全週間1928、オ工業標準化法制定1949、ア労働安全衛生法制定1972、イ製造物責任法制定1994の順で1。工場法施行1916、PL法施行1995と制定年を入れ替えない。改称した産業標準化法の現在名称も当時の制定時名称と区別する。

\Needspace{4\baselineskip}\paragraph{選択肢の検討（査読待ち初稿）}\mbox{}\par

\noindent\textbf{1：適切・選択対象} ウエオアイ。独立に求めた組合せ ウエオアイ と一致する。\par

\noindent\textbf{2：不適切} ウオエアイ。独立に求めた組合せ ウエオアイ と異なる。\par

\noindent\textbf{3：不適切} エウオイア。独立に求めた組合せ ウエオアイ と異なる。\par

\noindent\textbf{4：不適切} エオウアイ。独立に求めた組合せ ウエオアイ と異なる。\par

\noindent\textbf{5：不適切} オウアイエ。独立に求めた組合せ ウエオアイ と異なる。\par

\paragraph{再利用する考え方} 制定・施行・初回実施、改称と制定を区別。 否定方向と資料年、分母と単位を確認して五肢へ戻る。

\paragraph{共通理論} \hyperref[theory:env:safetyhistory]{第\ref*{theory:env:safetyhistory}節}

\par 参考：\cite{env60_safety}。

\par 難易度 3（校準初稿）、優先度未評価。

\section{第1〜5期科学技術基本計画の特徴}

\subsection{令和3年度 問5-6}
\label{q:r03-reg-basic-5-6}

\noindent\textbf{原典：} sources/R03/basic.pdf、PDF物理26頁。

\textbf{公式正答：} 4（single）。

\par\textbf{検討未確定：} 2・3期の主要句は一次資料索引で確認。1・4・5期本文の当該段落全体は追加照合。

\paragraph{問題の要旨} 第1〜5期科学技術基本計画の特徴。

\Needspace{4\baselineskip}\paragraph{確認した条件・記述（独自要約）}\mbox{}\par

\noindent\textbf{原典：} 全五肢と否定方向・組合せ表を画像確認。\par

\noindent\textbf{資料時点：} 令和3年度の指定資料・統計年を保持し、後年値に置換しない。\par

\paragraph{問われる核心} 第1〜5期科学技術基本計画の特徴。

\paragraph{必要知識} 計画の期・開始年・特徴句をセットで記憶。

\paragraph{考え方・正答を選ぶ理由（査読待ち初稿）} 第1期1996はウポストドクター等1万人支援、第2期2001はイ社会のため・社会の中の科学技術、第3期2006はア倫理的・法的・社会的課題、第4期2011はオ震災からの復興・再生、第5期2016はエSociety5.0。ウイアオエの4。ELSIの概念が以前になかったという意味ではなく、問の特徴句を指定計画へ対応させる。

\Needspace{4\baselineskip}\paragraph{選択肢の検討（査読待ち初稿）}\mbox{}\par

\noindent\textbf{1：不適切} イアウエオ。独立に求めた組合せ ウイアオエ と異なる。\par

\noindent\textbf{2：不適切} イウアオエ。独立に求めた組合せ ウイアオエ と異なる。\par

\noindent\textbf{3：不適切} ウアイエオ。独立に求めた組合せ ウイアオエ と異なる。\par

\noindent\textbf{4：適切・選択対象} ウイアオエ。独立に求めた組合せ ウイアオエ と一致する。\par

\noindent\textbf{5：不適切} ウイエアオ。独立に求めた組合せ ウイアオエ と異なる。\par

\paragraph{再利用する考え方} 計画の期・開始年・特徴句をセットで記憶。 否定方向と資料年、分母と単位を確認して五肢へ戻る。

\paragraph{共通理論} \hyperref[theory:env:scienceplans]{第\ref*{theory:env:scienceplans}節}

\par 参考：\cite{env60_plan2,env60_plan3}。

\paragraph{関連問題} \hyperref[q:r07-reg-basic-5-6]{R07-REG-BASIC-5-6}

\par 難易度 3（校準初稿）、優先度未評価。

\subsection{令和7年度 問5-6}
\label{q:r07-reg-basic-5-6}

\noindent\textbf{原典：} sources/R07/basic.pdf、PDF物理19頁。

\textbf{公式正答：} 4（single）。

\par\textbf{検討未確定：} 2・3期主要句の一次索引を確認。4〜6期原本文の該当段落と2020改正・施行日を追加査読。

\paragraph{問題の要旨} 第2〜6期の社会・ELSI・復興・Society・総合知。

\Needspace{4\baselineskip}\paragraph{確認した条件・記述（独自要約）}\mbox{}\par

\noindent\textbf{原典：} 全五肢と否定方向・組合せ表を画像確認。\par

\noindent\textbf{資料時点：} 令和7年度の指定資料・統計年を保持し、後年値に置換しない。\par

\paragraph{問われる核心} 第2〜6期の社会・ELSI・復興・Society・総合知。

\paragraph{必要知識} 科学技術政策の期と特徴、総合知、改正と計画開始。

\paragraph{考え方・正答を選ぶ理由（査読待ち初稿）} イ第2期2001は社会のため・社会の中、ア第3期2006はELSI、オ第4期2011は震災復興、エ第5期2016はSociety5.0、ウ第6期2021は人文社会科学と自然科学の融合による総合知。イアオエウの4。2020改正法の名称変更と2021開始の第6期計画を区別する。

\Needspace{4\baselineskip}\paragraph{選択肢の検討（査読待ち初稿）}\mbox{}\par

\noindent\textbf{1：不適切} イアウオエ。独立に求めた組合せ イアオエウ と異なる。\par

\noindent\textbf{2：不適切} アウエオイ。独立に求めた組合せ イアオエウ と異なる。\par

\noindent\textbf{3：不適切} アイオウエ。独立に求めた組合せ イアオエウ と異なる。\par

\noindent\textbf{4：適切・選択対象} イアオエウ。独立に求めた組合せ イアオエウ と一致する。\par

\noindent\textbf{5：不適切} アイオエウ。独立に求めた組合せ イアオエウ と異なる。\par

\paragraph{再利用する考え方} 科学技術政策の期と特徴、総合知、改正と計画開始。 否定方向と資料年、分母と単位を確認して五肢へ戻る。

\paragraph{共通理論} \hyperref[theory:env:scienceplans]{第\ref*{theory:env:scienceplans}節}

\par 参考：\cite{env60_plan2,env60_plan3}。

\paragraph{関連問題} \hyperref[q:r03-reg-basic-5-6]{R03-REG-BASIC-5-6}

\par 難易度 3（校準初稿）、優先度未評価。

\section{ラムサール条約の対象は湿地}

\subsection{令和6年度 問5-1}
\label{q:r06-reg-basic-5-1}

\noindent\textbf{原典：} sources/R06/basic.pdf、PDF物理17頁。

\textbf{公式正答：} 1（single）。

\par\textbf{検討未確定：} 各条約・議定書本文の定義と適用範囲の逐条確認は追加査読対象。

\paragraph{問題の要旨} ラムサール条約の対象は湿地。

\Needspace{4\baselineskip}\paragraph{確認した条件・記述（独自要約）}\mbox{}\par

\noindent\textbf{原典：} 全五肢と否定方向・組合せ表を画像確認。\par

\noindent\textbf{資料時点：} 令和6年度の指定資料・統計年を保持し、後年値に置換しない。\par

\paragraph{問われる核心} ラムサール条約の対象は湿地。

\paragraph{必要知識} 国際環境条約の対象と規制行為、ABS。

\paragraph{考え方・正答を選ぶ理由（査読待ち初稿）} 1が不適切。条約名だけではなく、対象（湿地・種・物質・廃棄物・遺伝資源）、行為（利用・取引・生産・移動・取得）を対応させる。森林に湿地が含まれる事例があっても条約全体を森林保全と定義しない。

\Needspace{4\baselineskip}\paragraph{選択肢の検討（査読待ち初稿）}\mbox{}\par

\noindent\textbf{1：不適切・選択対象} ラムサールは重要森林の保全。重要な湿地の保全と賢明な利用を目的とし、森林一般に置換するのは不適切。\par

\noindent\textbf{2：条件内で適切} ワシントンは希少野生種の国際取引規制。CITESは種の国際取引を管理する。\par

\noindent\textbf{3：条件内で適切} モントリオールはODSの生産消費規制。オゾン層破壊物質を対象とする枠組み。後の改正で対象は広がる。\par

\noindent\textbf{4：条件内で適切} バーゼルは有害廃棄物の越境移動・処分。国境を越える有害廃棄物等の管理。\par

\noindent\textbf{5：条件内で適切} 名古屋は遺伝資源取得と利益配分。ABSは取得機会と利用利益の公正衡平な配分。\par

\paragraph{再利用する考え方} 国際環境条約の対象と規制行為、ABS。 否定方向と資料年、分母と単位を確認して五肢へ戻る。

\paragraph{共通理論} \hyperref[theory:env:treaties]{第\ref*{theory:env:treaties}節}

\par 参考：\cite{env60_ramsar}。

\par 難易度 3（校準初稿）、優先度未評価。

\section{独立正規変数の和差と標本平均}

\subsection{平成23年度 問5-4}
\label{q:h23-reg-basic-5-4}

\noindent\textbf{原典：} sources/H23/basic.pdf、PDF物理17頁。

\textbf{公式正答：} 3（single）。

\paragraph{問題の要旨} 独立な製品質量の標本平均と総質量について、平均と標準偏差を求める。

\Needspace{4\baselineskip}\paragraph{確認した条件・記述（独自要約）}\mbox{}\par

\noindent\textbf{判断方向：} 最も適切な記述又は組合せを選ぶ\par

\noindent\textbf{原典照合：} 14原典画像で16問の条件、数値、単位、表の列と全5候補を確認。\par

\noindent\textbf{モデル：} 独立性、正規性、一定発生率などの前提を個別に保持する。\par

\paragraph{問われる核心} 独立な製品質量の標本平均と総質量について、平均と標準偏差を求める。

\paragraph{必要知識} 独立変数の線形結合、平均の線形性、分散への係数の二乗、正規分布の閉性。

\paragraph{考え方・正答を選ぶ理由（査読待ち初稿）} 各質量を独立正規変数とする。Aの標本平均は \(\bar A=(A_1+\cdots+A_4)/4\) で、期待値80g、分散 \(4(4^2)/4^2=4\,\mathrm{g^2}\) である。総質量は \(T=\sum_{i=1}^4A_i+\sum_{i=1}^4B_i\) なので期待値 \(4(80)+4(120)=800\,\mathrm g\)、分散 \(4(4^2)+4(3^2)=100\,\mathrm{g^2}\) となる。標準偏差はそれぞれ2gと10gで、3を選ぶ。平均を取る操作と、四つの値を合計する操作を混同しない。

\Needspace{4\baselineskip}\paragraph{選択肢の検討（査読待ち初稿）}\mbox{}\par

\noindent\textbf{1：不適切} 平均20・偏差4・合計400・偏差10。Aの標本平均の期待値は80gであり、標本数で平均を割らない。総質量の期待値も800gである。\par

\noindent\textbf{2：不適切} 平均80・偏差4・合計800・偏差7。平均の偏差は4gではなく2gである。合計の偏差も標準偏差をそのまま足した7gではなく10gである。\par

\noindent\textbf{3：適切} 平均80・偏差2・合計800・偏差10。Aの4個平均の偏差は4/2=\allowbreak{}2g、合計の分散は4×16+4×9=\allowbreak{}100g\textsuperscript{2}なので10gである。\par

\noindent\textbf{4：不適切} 平均20・偏差2・合計400・偏差7。平均と総質量の期待値がともに違い、合計の偏差も違う。\par

\noindent\textbf{5：不適切} 平均80・偏差2・合計800・偏差7。平均部分は合うが、合計の標準偏差は7gではない。独立成分の分散を足して最後に平方根を取る。\par

\paragraph{再利用する考え方} 変数の意味と単位を先に定め、平均・分散・確率・判定を順に分ける。品質の要求、工程の監視、ロットの受入れを区別し、成立条件のない一般化を避ける。

\paragraph{共通理論} \hyperref[theory:quality:variance]{第\ref*{theory:quality:variance}節}, \hyperref[theory:quality:normal]{第\ref*{theory:quality:normal}節}, \hyperref[theory:quality:sample_mean]{第\ref*{theory:quality:sample_mean}節}

\par 参考：\cite{stat_nist_normal,stat_psu_linear}。

\paragraph{関連問題} \hyperref[q:r06-reg-basic-1-6]{R06-REG-BASIC-1-6}

\par 難易度 3（校準初稿）、優先度未評価。

\section{工程の品質保証と標準化}

\subsection{平成24年度 問5-4}
\label{q:h24-reg-basic-5-4}

\noindent\textbf{原典：} sources/H24/basic.pdf、PDF物理17頁。

\textbf{公式正答：} 2（single）。

\paragraph{問題の要旨} 最終検査と上流の品質保証を分け、プロセスの標準化と改善の関係を読む。

\Needspace{4\baselineskip}\paragraph{確認した条件・記述（独自要約）}\mbox{}\par

\noindent\textbf{判断方向：} 最も適切な記述又は組合せを選ぶ\par

\noindent\textbf{原典照合：} 14原典画像で16問の条件、数値、単位、表の列と全5候補を確認。\par

\noindent\textbf{モデル：} 独立性、正規性、一定発生率などの前提を個別に保持する。\par

\paragraph{問われる核心} 最終検査と上流の品質保証を分け、プロセスの標準化と改善の関係を読む。

\paragraph{必要知識} 検査と予防、上流工程、標準化、PDCA、プロセスと成果。

\paragraph{考え方・正答を選ぶ理由（査読待ち初稿）} 完成品を検査して不適合品を隔離することと、不適合を生む工程を上流から改善することを分ける。設問の文脈ではアが検査、イが源流管理、ウが標準化、エがプロセスであり2となる。標準を設定して実施し、結果を確認して更新する循環が品質保証を支える。標準化は永久に方法を固定することではなく、改善結果を次の標準へ反映することを含む。

\Needspace{4\baselineskip}\paragraph{選択肢の検討（査読待ち初稿）}\mbox{}\par

\noindent\textbf{1：不適切} 検査・三現管理・基準化・プロセス。現場・現物・現実という三現の確認と、上流で不適合原因を抑える源流管理は同じ説明ではない。\par

\noindent\textbf{2：適切} 検査・源流管理・標準化・プロセス。要求適合を判定する検査、上流からの源流管理、方法を定める標準化、プロセスを管理する流れでつながる。\par

\noindent\textbf{3：不適切} 検査・三現管理・基準化・アウトカム。三現管理・基準化の対応に加え、ここで管理する対象は一連のプロセスである。\par

\noindent\textbf{4：不適切} 評価・源流管理・標準化・アウトカム。評価は広い概念だが本問の要求適合判定は検査である。末尾も成果だけを意味するアウトカムではなくプロセスである。\par

\noindent\textbf{5：不適切} 評価・三現管理・標準化・プロセス。要求適合判定と上流の改善に対応する語が検査・源流管理になっていない。\par

\paragraph{再利用する考え方} 変数の意味と単位を先に定め、平均・分散・確率・判定を順に分ける。品質の要求、工程の監視、ロットの受入れを区別し、成立条件のない一般化を避ける。

\paragraph{共通理論} \hyperref[theory:quality:assurance]{第\ref*{theory:quality:assurance}節}, \hyperref[theory:quality:pdca]{第\ref*{theory:quality:pdca}節}

\par 参考：\cite{quality_process_jsqc,quality_iso_process}。

\paragraph{関連問題} \hyperref[q:r03-reg-basic-1-3]{R03-REG-BASIC-1-3}

\par 難易度 2（校準初稿）、優先度未評価。
